\chapter{线积分}
本章我们回到实变量函数的积分理论。现在，我们将考虑的积分不仅在区间上，而且在区间的连续可微映射，即曲线上进行。我们将看到，积分的这一推广具有重要而深远的影响。

当然，我们必须先精确化曲线的概念，这将在第 1 节中完成。此外，我们将引入弧长的概念，并导出计算弧长的积分公式。

在第 2 节中，我们讨论曲线的微分几何。特别地，我们证明相伴 $ n $ 标架的存在性，它是沿曲线定义的一组 $ n $ 个向量。对平面曲线，我们研究曲率；对空间曲线，我们还将研究挠率。本节的内容主要贡献于读者的一般数学知识，本章的其余部分将不再需要它。

第 3 节处理一阶微分形式。这里我们使第 VI 章中引入的微分的特设定义严格化。我们还将导出处理这类微分形式的若干简单法则。这些法则是线积分理论的基础，这一点我们将在第 4 节中得知，届时微分形式将作为线积分的被积函数出现。我们将证明线积分基本定理，它刻画了可作为某个势函数的梯度而得到的向量场。

第 5 节和第 6 节给出了线积分在复分析（亦称为复变函数理论）中的一个特别重要的应用。在第 5 节中，我们导出全纯函数的基本性质，特别是柯西积分定理与柯西积分公式。以这些为工具，我们证明基本结果：一个函数全纯当且仅当它解析。我们将一般理论应用于菲涅耳积分，从而说明如何利用柯西积分定理计算实积分。

在第 6 节中，我们研究亚纯函数，并证明重要的留数定理及其在同调理论中的版本。为此，我们引入绕数的概念，并导出其最重要的性质。作为本卷的结尾，我们通过计算若干傅里叶积分来展示留数定理的作用范围。

\section{曲线及其长度}

本节主要证明，连续可微的紧致曲线 $ \Gamma $ 具有由下式给出的有限长度 $ L(\Gamma) $：

\[
L(\Gamma) = \int_a^b |\dot{\gamma}(t)|\,dt .
\]

这里 $ \gamma $ 是 $ \Gamma $ 的任意参数化。

在下文中，设

• $ E = (E,|\cdot|) $ 是域 $ \mathbb{K} $ 上的巴拿赫空间，且 $ I = [a,b] $ 是紧致区间。

\needspace{3\baselineskip}\subsection*{全变差}

设 $ f : I \to E $，且 $ \mathfrak{Z} := (t_0,\ldots,t_n) $ 是 $ I $ 的一个划分。则

\[
L_{\mathfrak{Z}}(f) := \sum_{j=1}^{n}|f(t_j) - f(t_{j-1})|
\]

是 $ E $ 中 $ (f(t_0),\ldots,f(t_n)) $ 这条\textbf{分段直线道路的长度}，而

\[
\operatorname{Var}(f,I) := \sup\{ L_{\mathfrak{Z}}(f) \;;\; \mathfrak{Z} = (t_0,\ldots,t_n) \text{ 是 } I \text{ 的划分} \}
\]

称为 $ f $ 在 $ I $ 上的\textbf{全变差}（或简称\textbf{变差}）。若 $ \operatorname{Var}(f,I) < \infty $，则称 $ f $ 是\textbf{有界变差}的。

\bookfigure[0.38]{figures/396c68a1483532ec.png}

\begin{lemma}

对 $ f : [a,b] \to E $ 与 $ c \in [a,b] $，有

\begin{equation}
\operatorname{Var}(f,[a,b]) = \operatorname{Var}(f,[a,c]) + \operatorname{Var}(f,[c,b]) .\tag{1.1}
\end{equation}

证明

\begin{subitems}
\item

设 $ c \in [a,b] $。不失一般性，可设 $ \operatorname{Var}(f,[a,c]) $ 与 $ \operatorname{Var}(f,[c,b]) $ 均有限，否则函数 $ f : [a,b] \to E $ 就不是有界变差的，结论显然成立。

\item

设 $ \mathfrak{Z} $ 是 $ [a,b] $ 的一个划分，$ \tilde{\mathfrak{Z}} $ 是 $ \mathfrak{Z} $ 的含有 $ c $ 点的加细。此外，令 $ \tilde{\mathfrak{Z}}_1 := \tilde{\mathfrak{Z}} \cap [a,c] $ 与 $ \tilde{\mathfrak{Z}}_2 := \tilde{\mathfrak{Z}} \cap [c,b] $。则

\[
L_{\mathfrak{Z}}(f) \le L_{\tilde{\mathfrak{Z}}}(f) = L_{\tilde{\mathfrak{Z}}_1}(f) + L_{\tilde{\mathfrak{Z}}_2}(f) \le \operatorname{Var}(f,[a,c]) + \operatorname{Var}(f,[c,b]) .
\]

对 $ \mathfrak{Z} $ 取上确界，得

\[
\operatorname{Var}(f,[a,b]) \le \operatorname{Var}(f,[a,c]) + \operatorname{Var}(f,[c,b]) .
\]
\item

对 $ \varepsilon > 0 $，存在 $ [a,c] $ 的划分 $ \mathfrak{Z}_1 $ 与 $ [c,b] $ 的划分 $ \mathfrak{Z}_2 $，使得

\[
L_{\mathfrak{Z}_1}(f) \ge \operatorname{Var}(f,[a,c]) - \varepsilon/2 \quad \text{且}\quad L_{\mathfrak{Z}_2}(f) \ge \operatorname{Var}(f,[c,b]) - \varepsilon/2 .
\]

对 $ \mathfrak{Z} := \mathfrak{Z}_1 \vee \mathfrak{Z}_2 $，有

\[
L_{\mathfrak{Z}_1}(f) + L_{\mathfrak{Z}_2}(f) = L_{\mathfrak{Z}}(f) \le \operatorname{Var}(f,[a,b]) ,
\]

从而

\[
\operatorname{Var}(f,[a,c]) + \operatorname{Var}(f,[c,b]) \le L_{\mathfrak{Z}_1}(f) + L_{\mathfrak{Z}_2}(f) + \varepsilon \le \operatorname{Var}(f,[a,b]) + \varepsilon .
\]

结论现在由 (ii) 推出。\nobreak\hspace{0.35em}■

\end{subitems}

\end{lemma}

\needspace{3\baselineskip}\subsection*{可求长道路}

把 $ \gamma \in C(I,E) $ 理解为 $ E $ 中的连续道路，我们也称 $ \operatorname{Var}(\gamma,I) $ 为 $ \gamma $ 的\textbf{长度}（或\textbf{弧长}），并记为 $ L(\gamma) $。若 $ L(\gamma) < \infty $，即 $ \gamma $ 具有有限长度，则称 $ \gamma $ 是\textbf{可求长}的。

\begin{remark}

\begin{items}
\item

存在不可求长的连续道路。\footnote{还可以证明，$ \mathbb{R}^{2} $ 中存在连续道路，其像填满整个单位圆盘 $ \mathbb{B} $（见习题 8）。这类“填满空间”的道路称为\textbf{皮亚诺曲线}。}

\proofhead
考虑 $ \gamma : [0,1] \to \mathbb{R} $，其中 $ \gamma(0) := 0 $，且对 $ t \in (0,1] $ 有 $ \gamma(t) := t\cos^{2}(\pi/2t) $。则 $ \gamma $ 连续（见命题 III.2.24）。对 $ n \in \mathbb{N}^\times $，设 $ \mathfrak{Z}_n = (t_0,\ldots,t_{2n}) $ 是 $ [0,1] $ 的满足 $ t_0 = 0 $ 且对 $ 1 \le j \le 2n $ 有 $ t_j = 2/(2n + 1 - j) $ 的划分。因为
\begin{equation*}
\gamma(t_j) = \begin{cases} 0, & j = 2k,\quad 0 \le k \le n, \\ t_j, & j = 2k + 1,\quad 0 \le k < n, \end{cases}
\end{equation*}

有

\[
L_{\mathfrak{Z}_n}(\gamma) = \sum_{j=1}^{2n}|\gamma(t_j) - \gamma(t_{j-1})| = \sum_{k=0}^{n-1} t_{2k+1} = \frac{1}{2}\sum_{k=1}^{n}\frac{1}{k} \to \infty
\]

当 $ n \to \infty $ 时。因此 $ \gamma $ 不可求长。\nobreak\hspace{0.35em}■

\item

设 $ \gamma : [a,b] \to E $ 是利普希茨连续的，其利普希茨常数为 $ \lambda $。则 $ \gamma $ 可求长，且 $ L(\gamma) \le \lambda(b - a) $。

\proofhead
对 $ [a,b] $ 的每个划分 $ \mathfrak{Z} = (t_0,\ldots,t_n) $，有
\[
L_{\mathfrak{Z}}(\gamma) = \sum_{j=1}^{n}|\gamma(t_j) - \gamma(t_{j-1})| \le \lambda\sum_{j=1}^{n}|t_j - t_{j-1}| = \lambda(b - a) .\quad\blacksquare
\]
\item

当然，道路 $ \gamma $ 的长度依赖于 $ E $ 上的范数。然而，当换成另一个不同但等价的范数时，其可求长性并不改变。\nobreak\hspace{0.35em}■

\end{items}

注 1.2(a) 中所考虑的道路确实连续，但在 0 处不可微。下面的结果表明，连续可微的道路总是可求长的。

\end{remark}

\begin{theorem}

设 $ \gamma \in C^{1}(I,E) $。则 $ \gamma $ 可求长，且

\[
L(\gamma) = \int_a^b |\dot{\gamma}(t)|\,dt .
\]

证明

\begin{subitems}
\item

只需考虑 $ a < b $ 的情形，此时 $ I $ 不只是一个单点。

\item

$ \gamma $ 的可求长性直接由微积分基本定理得到。即，若 $ \mathfrak{Z} = (t_0,\ldots,t_n) $ 是 $ [a,b] $ 的划分，则

\begin{equation*}
\begin{aligned}
L_{\mathfrak{Z}}(\gamma) &= \sum_{j=1}^{n}|\gamma(t_j) - \gamma(t_{j-1})| = \sum_{j=1}^{n}\left|\int_{t_{j-1}}^{t_j}\dot{\gamma}(t)\,dt\right| \\ &\le \sum_{j=1}^{n}\int_{t_{j-1}}^{t_j}|\dot{\gamma}(t)|\,dt = \int_a^b |\dot{\gamma}(t)|\,dt .
\end{aligned}
\end{equation*}

因此

\begin{equation}
L(\gamma) = \operatorname{Var}(\gamma,[a,b]) \le \int_a^b |\dot{\gamma}(t)|\,dt .\tag{1.2}
\end{equation}
\item

现在设 $ s_0 \in [a,b) $。由引理 1.1 可知，对每个 $ s \in (s_0,b) $ 有

\[
\operatorname{Var}(\gamma,[a,s]) - \operatorname{Var}(\gamma,[a,s_0]) = \operatorname{Var}(\gamma,[s_0,s]) .
\]

我们还得到

\[
|\gamma(s) - \gamma(s_0)| \le \operatorname{Var}(\gamma,[s_0,s]) \le \int_{s_0}^{s}|\dot{\gamma}(t)|\,dt .
\]

这里第一个不等式成立是因为 $ (s_0,s) $ 是 $ [s_0,s] $ 的一个划分，第二个来自 (1.2)。
于是，由于 $ s_0 < s $，有

\begin{equation}
\left|\frac{\gamma(s) - \gamma(s_0)}{s - s_0}\right| \le \frac{\operatorname{Var}(\gamma,[a,s]) - \operatorname{Var}(\gamma,[a,s_0])}{s - s_0} \le \frac{1}{s - s_0}\int_{s_0}^{s}|\dot{\gamma}(t)|\,dt .\tag{1.3}
\end{equation}

因为 $ \gamma $ 有连续导数，由积分形式的中值定理与定理 VI.4.12 可得

\[
|\dot{\gamma}(s_0)| = \lim_{s \to s_0}\left|\frac{\gamma(s) - \gamma(s_0)}{s - s_0}\right| \le \lim_{s \to s_0}\left[\frac{1}{s - s_0}\int_{s_0}^{s}|\dot{\gamma}(t)|\,dt\right] = |\dot{\gamma}(s_0)| .
\]

于是，(1.3) 以及关于 $ s < s_0 $ 的类似考虑表明 $ s \mapsto \operatorname{Var}(\gamma,[a,s]) $ 可微，且

\[
\frac{d}{ds}\operatorname{Var}(\gamma,[a,s]) = |\dot{\gamma}(s)| \qquad \text{对 } s \in [a,b] .
\]

因此 $ s \mapsto \operatorname{Var}(\gamma,[a,s]) $ 属于 $ C^{1}(I,\mathbb{R}) $。此外，微积分基本定理给出

\[
\operatorname{Var}(\gamma,[a,b]) = \int_a^b |\dot{\gamma}(t)|\,dt
\]

因为 $ \operatorname{Var}(\gamma,[a,a]) = 0 $。\nobreak\hspace{0.35em}■

\end{subitems}

\end{theorem}

\begin{corollary}

对 $ \gamma = (\gamma_1,\ldots,\gamma_n) \in C^{1}(I,\mathbb{R}^n) $，有

\[
L(\gamma) = \int_a^b\sqrt{\left(\dot{\gamma}_1(t)\right)^{2} + \cdots + \left(\dot{\gamma}_n(t)\right)^{2}}\,dt .
\]
\end{corollary}

\needspace{3\baselineskip}\subsection*{可微曲线}

道路的像是 $ E $ 中不依赖于用来描述它的特定函数的点集。换言之，它具有独立于任何特定参数化的几何意义。为精确地理解这一点，我们必须明确哪些参数变换保持这个像不变。

设 $ J_1 $ 与 $ J_2 $ 是区间，且 $ q \in \mathbb{N} \cup \{\infty\} $。若 $ \varphi $ 属于\footnote{此外，若 $ J_1 $ 与 $ J_2 $ 不是开的，则这意味着 $ \varphi : J_1 \to J_2 $ 是双射，且 $ \varphi $ 与 $ \varphi^{-1} $ 属于 $ C^{q} $ 类。特别地，$ \operatorname{Diff}^{0}(J_1,J_2) $ 是 $ J_1 $ 到 $ J_2 $ 的所有拓扑映射（同胚）之集。} $ \operatorname{Diff}^{q}(J_1,J_2) $ 且严格递增，则称映射 $ \varphi : J_1 \to J_2 $ 是一个\textbf{（保持定向的）} $ C^{q} $ \textbf{参数变换}。若对 $ j = 1,2 $ 有 $ \gamma_j \in C^{q}(J_j,E) $，且存在 $ C^{q} $ 参数变换 $ \varphi $ 使得 $ \gamma_1 = \gamma_2 \circ \varphi $，则称 $ \gamma_1 $ 是 $ \gamma_2 $ 的一个\textbf{（保持定向的）} $ C^{q} $ \textbf{重新参数化}。

\begin{remark}

\begin{items}
\item

当 $ \varphi \in \operatorname{Diff}^{q}(J_1,J_2) $ 严格递减时，我们称 $ \varphi $ 是一个\textbf{反转定向的} $ C^{q} $ \textbf{参数变换}。不过，在下文中，我们总是取保持定向的参数变换。

\item

映射 $ \varphi : J_1 \to J_2 $ 是 $ C^{q} $ 参数变换，当且仅当 $ \varphi $ 属于 $ C^{q}(J_1,J_2) $、是满射，且对 $ t \in J_1 $ 满足 $ \dot{\varphi}(t) > 0 $。

\proofhead
这由定理 III.5.7 与 IV.2.8 得到。\nobreak\hspace{0.35em}■
\item

设 $ I_1 $ 与 $ I_2 $ 是紧致区间，且设 $ \gamma_1 \in C(I_1,E) $ 是 $ \gamma_2 \in C(I_2,E) $ 的连续重新参数化。则

\[
\operatorname{Var}(\gamma_1,I_1) = \operatorname{Var}(\gamma_2,I_2) .
\]

\proofhead
设 $ \varphi \in \operatorname{Diff}^{0}(I_1,I_2) $ 是满足 $ \gamma_1 = \gamma_2 \circ \varphi $ 的参数变换。若 $ \mathfrak{Z} = (t_0,\ldots,t_n) $ 是 $ I_1 $ 的一个划分，则 $ \varphi(\mathfrak{Z}) := (\varphi(t_0),\ldots,\varphi(t_n)) $ 是 $ I_2 $ 的一个划分，且
\[
L_{\mathfrak{Z}}(\gamma_1,I_1) = L_{\mathfrak{Z}}(\gamma_2 \circ \varphi,I_1) = L_{\varphi(\mathfrak{Z})}(\gamma_2,I_2) \le \operatorname{Var}(\gamma_2,I_2) .
\]

于是有关系

\begin{equation}
\operatorname{Var}(\gamma_1,I_1) = \operatorname{Var}(\gamma_2 \circ \varphi,I_1) \le \operatorname{Var}(\gamma_2,I_2) .\tag{1.4}
\end{equation}

注意到 $ \gamma_2 = (\gamma_2 \circ \varphi) \circ \varphi^{-1} $，由 (1.4)（把 $ \gamma_2 $ 换成 $ \gamma_2 \circ \varphi $、把 $ \varphi $ 换成 $ \varphi^{-1} $）可得

\[
\operatorname{Var}(\gamma_2,I_2) = \operatorname{Var}((\gamma_2 \circ \varphi) \circ \varphi^{-1},I_2) \le \operatorname{Var}(\gamma_2 \circ \varphi,I_1) = \operatorname{Var}(\gamma_1,I_1) .\quad\blacksquare
\]

在 $ E $ 中所有 $ C^{q} $ 道路的集合上，我们定义关系 $ \sim $ 如下

\[
\gamma_1 \sim \gamma_2 :\Longleftrightarrow \gamma_1 \text{ 是 } \gamma_2 \text{ 的一个 } C^{q} \text{ 重新参数化} .
\]

不难看出 $ \sim $ 是等价关系（见习题 5）。相应的等价类称为 $ E $ 中的 \textbf{$ C^{q} $ 曲线}。$ C^{q} $ 曲线 $ \Gamma $ 的每个代表元都是 $ \Gamma $ 的一个 \textbf{$ C^{q} $ 参数化}。我们可以说 $ C^{0} $ 曲线是\textbf{连续的}、$ C^{1} $ 曲线是\textbf{连续可微的}、$ C^{\infty} $ 曲线是\textbf{光滑的}。若曲线 $ \Gamma $ 有定义在紧致定义域（或紧致\textbf{参数区间}）上的参数化，则称该曲线是\textbf{紧致的}，且由定理 III.3.6，$ \Gamma $ 的像也是紧致的。若对 $ t \in \operatorname{dom}(\gamma) $ 有 $ \dot{\gamma}(t) \ne 0 $，则称 $ \Gamma $ 的参数化 $ \gamma $ 是\textbf{正则的}。当 $ \Gamma $ 有正则参数化时，称 $ \Gamma $ 为\textbf{正则曲线}。有时我们写 $ \Gamma = [\gamma] $，以强调 $ \Gamma $ 是含有 $ \gamma $ 的参数化等价类。

\end{items}

\end{remark}

\begin{remark}

\begin{items}
\item

若 $ \Gamma $ 是紧致曲线，则它的每个参数化都有紧致定义域。若 $ \Gamma $ 是正则 $ C^{1} $ 曲线，则它的每个参数化都是正则的。

\proofhead
设 $ \gamma \in C^{q}(J,E) $ 是 $ \Gamma $ 的一个参数化，$ \gamma_1 \in C^{q}(J_1,E) $ 是 $ \gamma $ 的一个重新参数化。则存在 $ \varphi \in \operatorname{Diff}^{q}(J_1,J) $，使得 $ \gamma_1 = \gamma \circ \varphi $。当 $ J $ 紧致时，对 $ J_1 = \varphi^{-1}(J) $ 也成立，因为紧致集的连续像也是紧致的（定理 III.3.6）。由链式法则，对 $ q \ge 1 $ 得
\[
\dot{\gamma}_1(t) = \dot{\gamma}(\varphi(t))\dot{\varphi}(t) \qquad \text{对 } t \in J_1 .
\]

因此，由于对 $ s \in J $ 有 $ \dot{\gamma}(s) \ne 0 $、对 $ t \in J_1 $ 有 $ \dot{\varphi}(t) \ne 0 $，故对 $ t \in J_1 $ 也有 $ \dot{\gamma}_1(t) \ne 0 $。\nobreak\hspace{0.35em}■

\item

正则曲线可以有非正则的参数化，但它与其正则参数化不等价。

\proofhead
考虑由函数
\[
[-1,1] \to \mathbb{R}^{2} , \quad t \mapsto \gamma(t) := (t,t) .
\]

参数化的正则光滑曲线 $ \Gamma $。再考虑光滑道路 $ \tilde{\gamma} : [-1,1] \to \mathbb{R}^{2} $，其中 $ \tilde{\gamma}(t) := (t^{3},t^{3}) $。则 $ \operatorname{im}(\gamma) = \operatorname{im}(\tilde{\gamma}) $，但 $ \dot{\tilde{\gamma}}(0) = (0,0) $。因此 $ \tilde{\gamma} $ 不是 $ \gamma $ 的 $ C^{1} $ 重新参数化。\nobreak\hspace{0.35em}■

\item

设 $ I \to E $, $ t \mapsto \gamma(t) $ 是曲线 $ \Gamma $ 的正则 $ C^{1} $ 参数化。则

\[
\dot{\gamma}(t) = \lim_{s \to t}\frac{\gamma(s) - \gamma(t)}{s - t} \qquad \text{对 } t \in I
\]

可解释为曲线 $ \Gamma $ 在点 $ \gamma(t) $ 处的“瞬时速度”。\nobreak\hspace{0.35em}■

\end{items}

设 $ \Gamma $ 是 $ E $ 中的连续曲线，且设 $ \gamma $ 与 $ \gamma_1 $ 是 $ \Gamma $ 的（等价）参数化。则显然 $ \operatorname{im}(\gamma) = \operatorname{im}(\gamma_1) $。因此 $ \Gamma $ 的\textbf{像}由下式良定义

\[
\operatorname{im}(\Gamma) := \operatorname{im}(\gamma) .
\]

若 $ \Gamma $ 紧致，$ \operatorname{dom}(\gamma) = [a,b] $，$ \operatorname{dom}(\gamma_1) = [a_1,b_1] $，则有关系

\[
\gamma(a) = \gamma_1(a_1) \quad \text{且}\quad \gamma(b) = \gamma_1(b_1) .
\]

因此，对紧致曲线，\textbf{起点} $ A_{\Gamma} := \gamma(a) $ 与\textbf{终点} $ E_{\Gamma} := \gamma(b) $ 良定义。

若 $ A_{\Gamma} $ 与 $ E_{\Gamma} $ 重合，则 $ \Gamma $ 是\textbf{闭的}，或是一个\textbf{环}。最后，我们常简单地用 $ p \in \Gamma $ 表示 $ p \in \operatorname{im}(\Gamma) $。

\end{remark}

\needspace{3\baselineskip}\subsection*{可求长曲线}

设 $ \Gamma $ 是 $ E $ 中的连续曲线，$ \gamma \in C(J,E) $ 是 $ \Gamma $ 的一个参数化。又设 $ \alpha := \inf J $ 与 $ \beta := \sup J $。则对 $ \alpha < a < b < \beta $，$ \operatorname{Var}(\gamma,[a,b]) $ 有定义。于是引理 1.1 蕴含，对每个 $ c \in (\alpha,\beta) $，函数

\[
[c,\beta) \to \overline{\mathbb{R}} , \quad b \mapsto \operatorname{Var}(\gamma,[c,b])
\]

递增，而映射

\[
(\alpha,c] \to \overline{\mathbb{R}} , \quad a \mapsto \operatorname{Var}(\gamma,[a,c])
\]

递减。于是由 (1.1)，极限

\begin{equation}
\operatorname{Var}(\gamma,J) := \lim_{\substack{a \downarrow \alpha \\ b \uparrow \beta}}\operatorname{Var}(\gamma,[a,b])\tag{1.5}
\end{equation}

在 $ \overline{\mathbb{R}} $ 中存在，我们称之为 $ \gamma $ 在 $ J $ 上的（\textbf{全}）\textbf{变差}。若 $ \gamma_1 \in C(J_1,E) $ 是 $ \gamma $ 的重新参数化，则由注 1.5(c) 得 $ \operatorname{Var}(\gamma,J) = \operatorname{Var}(\gamma_1,J_1) $。因此 $ \Gamma $ 的\textbf{全变差}或\textbf{长度}（或\textbf{弧长}）由下式良定义

\[
L(\Gamma) := \operatorname{Var}(\Gamma) := \operatorname{Var}(\gamma,J) .
\]

若曲线 $ \Gamma $ 具有有限长度，即 $ L(\Gamma) < \infty $，则称其为\textbf{可求长}的。

\begin{theorem}

设 $ \gamma \in C^{1}(J,E) $ 是 $ C^{1} $ 曲线 $ \Gamma $ 的一个参数化。则

\begin{equation}
L(\Gamma) = \int_J |\dot{\gamma}(t)|\,dt .\tag{1.6}
\end{equation}

若 $ \Gamma $ 紧致，则 $ \Gamma $ 可求长。

\proofhead
这直接由 (1.5)、定理 1.3 与反常积分的定义得到。\nobreak\hspace{0.35em}■
\end{theorem}

\begin{remark}

\begin{items}
\item

定理 1.7 特别表明，$ L(\Gamma) $，因而 (1.6) 中的积分，不依赖于特定的参数化 $ \gamma $。

\item

设 $ \tilde{\Gamma} $ 是按 VII.8 节所述方式嵌入 $ \mathbb{R}^{n} $ 中的 $ C^{q} $ 曲线。又设 $ (\varphi,U) $ 是 $ \tilde{\Gamma} $ 的一个坐标卡，$ (g,V) $ 是其相应的参数化。则\footnote{这里以及在不易引起混淆的其他情形中，我们简单地用 $ \Gamma $ 表示 $ \operatorname{im}(\Gamma) $。} $ \Gamma := \tilde{\Gamma} \cap U $ 是正则 $ C^{q} $ 曲线，且 $ g $ 是 $ \Gamma $ 的正则 $ C^{q} $ 参数化。

\item

正则 $ C^{q} $ 曲线一般不是嵌入 $ C^{q} $ 曲线。

\proofhead
这由例 VII.9.6(c) 得到。\nobreak\hspace{0.35em}■
\item

由于我们只允许保持定向的参数变换，我们的曲线都是\textbf{定向的}，即它们按固定方向“穿过”。\nobreak\hspace{0.35em}■

\end{items}

\end{remark}

\begin{example}

\begin{items}
\item

（实值函数的图像）设 $ f \in C^{q}(J,\mathbb{R}) $，$ \Gamma := \operatorname{graph}(f) $。则 $ \Gamma $ 是 $ \mathbb{R}^{2} $ 中的正则 $ C^{q} $ 曲线，且映射 $ J \to \mathbb{R}^{2} $, $ t \mapsto (t,f(t)) $ 是 $ \Gamma $ 的正则 $ C^{q} $ 参数化。此外

\[
L(\Gamma) = \int_J\sqrt{1 + [f'(t)]^{2}}\,dt .
\]
\item

（极坐标下的平面曲线）设 $ r,\varphi \in C^{q}(J,\mathbb{R}) $，且对 $ t \in J $ 有 $ r(t) \ge 0 $。又设

\[
\gamma(t) := r(t)(\cos(\varphi(t)),\sin(\varphi(t))) \qquad \text{对 } t \in J .
\]

将 $ \mathbb{R}^{2} $ 与 $ \mathbb{C} $ 等同起来，得到 $ \gamma $ 的表示

\[
\gamma(t) = r(t)e^{i\varphi(t)} \qquad \text{对 } t \in J .
\]

若 $ [\dot{r}(t)]^{2} + [r(t)\dot{\varphi}(t)]^{2} > 0 $，则 $ \gamma $ 是曲线 $ \Gamma $ 的正则 $ C^{q} $ 参数化，且

\[
L(\Gamma) = \int_J\sqrt{[\dot{r}(t)]^{2} + [r(t)\dot{\varphi}(t)]^{2}}\,dt .
\]

\proofhead
有
\[
\dot{\gamma}(t) := (\dot{r}(t) + ir(t)\dot{\varphi}(t))e^{i\varphi(t)} ,
\]

从而 $ |\dot{\gamma}(t)|^{2} = [\dot{r}(t)]^{2} + [r(t)\dot{\varphi}(t)]^{2} $，由此得证。\nobreak\hspace{0.35em}■

\item

设 $ 0 < b \le 2\pi $。对由

\[
\gamma : [0,b] \to \mathbb{R}^{2} , \quad t \mapsto R(\cos t,\sin t)
\]

参数化的圆弧，或经由

\[
\gamma : [0,b] \to \mathbb{C} , \quad t \mapsto Re^{it} ,
\]

得到的等价圆弧，将 $ \mathbb{R}^{2} $ 与 $ \mathbb{C} $ 等同起来给出\footnote{见习题 III.6.12。} $ L(\Gamma) = bR $。

\end{items}

\bookfigure[0.38]{figures/cb2e255f15b4f9d7.png}

\proofhead
这由 (b) 得到。\nobreak\hspace{0.35em}■

(d)（对数螺线）对 $ \lambda < 0 $ 与 $ a \in \mathbb{R} $，
\[
\gamma_{a,\infty} : [a,\infty) \to \mathbb{R}^{2} , \quad t \mapsto e^{\lambda t}(\cos t,\sin t)
\]

是曲线 $ \Gamma_{a,\infty} := [\gamma_{a,\infty}] $ 的光滑正则参数化。它具有有限长度 $ e^{\lambda a}\sqrt{1 + \lambda^{2}}/|\lambda| $。

当 $ \lambda > 0 $ 时，类似地有

\[
\gamma_{-\infty,a} : (-\infty,a] \to \mathbb{R}^{2} , \quad t \mapsto e^{\lambda t}(\cos t,\sin t)
\]

是曲线 $ \Gamma_{-\infty,a} := [\gamma_{-\infty,a}] $ 的光滑正则参数化，其有限长度为 $ e^{\lambda a}\sqrt{1 + \lambda^{2}}/\lambda $。对 $ \lambda \ne 0 $，映射

\[
\mathbb{R} \to \mathbb{R}^{2} , \quad t \mapsto e^{\lambda t}(\cos t,\sin t)
\]

是\textbf{对数螺线}的光滑正则参数化。其长度无穷。当 $ \lambda > 0 $ [或 $ \lambda < 0 $] 时，它向外[或向内]盘旋。将 $ \mathbb{R}^{2} $ 与 $ \mathbb{C} $ 等同起来，给对数螺线以简单的参数化 $ t \mapsto e^{(\lambda+i)t} $。

\proofhead
设 $ \lambda < 0 $。对 $ t \in [a,\infty) $，令 $ r(t) = e^{\lambda t} $ 与 $ \varphi(t) = t $。据 (b)，有
\[
L(\Gamma_{a,\infty}) = \lim_{b \to \infty}\int_a^b\sqrt{1 + \lambda^{2}}\,e^{\lambda t}\,dt = \frac{\sqrt{1 + \lambda^{2}}}{\lambda}\lim_{b \to \infty}(e^{\lambda b} - e^{\lambda a}) = \frac{\sqrt{1 + \lambda^{2}}}{|\lambda|}e^{\lambda a} .
\]

$ \lambda > 0 $ 的情形类似地处理。\nobreak\hspace{0.35em}■

\bookfigure[0.85]{figures/e769ab47f59684fb.png}

\begin{items}
\item[（e）]

对 $ R > 0 $ 与 $ h > 0 $，由

\begin{equation}
\gamma : \mathbb{R} \to \mathbb{R}^{3} , \quad t \mapsto (R\cos t, R\sin t, ht)\tag{1.7}
\end{equation}

参数化的正则光滑曲线 $ \Gamma $ 是半径为 $ R $、螺距为 $ 2\pi h $ 的\textbf{螺旋线}。通过将\footnote{作为度量空间} $ \mathbb{R}^{2} $ 与 $ \mathbb{C} $ 等同起来，从而将 $ \mathbb{R}^{3} $ 与 $ \mathbb{C} \times \mathbb{R} $ 等同起来，我们把 (1.7) 写成形式 $ t \mapsto (Re^{it},ht) $。有

\[
L(\gamma|[a,b]) = (b - a)\sqrt{R^{2} + h^{2}}
\]

其中 $ -\infty < a < b < \infty $。曲线 $ \Gamma $ 位于半径为 $ R $、轴为 $ z $ 轴的柱面上。每旋转一周，它沿轴向上升 $ 2\pi h $。

\proofhead
因为 $ |\dot{\gamma}(t)|^{2} = R^{2} + h^{2} $，结论由定理 1.7 得到。\nobreak\hspace{0.35em}■
\end{items}

\bookfigure[0.30]{figures/a6dd95e9161b79eb.png}
\end{example}

\exerciseshead

\begin{exercises}
\item

一个圆盘沿一条直线无滑动地滚动。求圆盘上固定的任一点（或圆盘“外”的一点）所描出的道路的一个参数化。确定圆盘转动一周时这条道路的弧长。

\item

设 $ -\infty < \alpha < \beta < \infty $。计算这些道路的长度：

\begin{equation*}
\begin{aligned}
&[0, \pi] \to \mathbb{R}^{2}, && t \mapsto (\cos t + t\sin t, \sin t - t\cos t)\ ; \\ &[1, \infty) \to \mathbb{R}^{2}, && t \mapsto t^{-2}(\cos t, \sin t)\ ; \\ &[\alpha, \beta] \to \mathbb{R}^{2}, && t \mapsto (t^{3}, 3t^{2}/2)\ ; \\ &[\alpha, \beta] \to \mathbb{R}^{2}, && t \mapsto (t, t^{2}/2)\ .
\end{aligned}
\end{equation*}
\item

逼近（例如，用梯形法则）帕斯卡蜗线的长度（见例 VII.9.1(c)）。

\item

画出空间曲线 $ \Gamma = [\gamma] $ 的草图，其中

\[
\gamma : [-3\pi, 3\pi] \to \mathbb{R}^{3} , \quad t \mapsto t(\cos t, \sin t, 1) ,
\]

并计算其弧长。

\item

设 $ \mathfrak{Z} = (\alpha_{0}, \ldots, \alpha_{m}) $（其中 $ m \in \mathbb{N}^{\times} $）是紧致区间 $ I $ 的一个划分，且 $ q \in \mathbb{N}^{\times} \cup \{\infty\} $。若对 $ j = 1, \ldots, m $ 有 $ \gamma_{j} := \gamma|[\alpha_{j-1}, \alpha_{j}] \in C^{q}([\alpha_{j-1}, \alpha_{j}], E) $，则连续道路 $ \gamma \in C(I, E) $ 是 $ E $ 中的\textbf{分段 $ C^{q} $ 道路}（关于 $ \mathfrak{Z} $）。若存在 $ C^{q} $ 参数变换 $ \varphi_{j} \in \operatorname{Diff}^{q}([\alpha_{j-1}, \alpha_{j}], [\beta_{j-1}, \beta_{j}]) $ 使得对 $ j = 1, \ldots, m $ 有 $ \gamma_{j} := \eta_{j} \circ \varphi_{j} $，则称在 $ J $ 的划分 $ \mathfrak{Z}' = (\beta_{0}, \ldots, \beta_{m}) $ 上分段 $ C^{q} $ 的道路 $ \eta \in C(J, E) $ 为 $ \gamma $ 的一个 \textbf{$ C^{q} $ 重新参数化}。在所有 $ E $ 中的分段 $ C^{q} $ 道路组成的集合上，我们通过下式定义 $ \sim $：

\[
\gamma \sim \eta :\Longleftrightarrow \eta \text{ 是 } \gamma \text{ 的重新参数化} .
\]

\begin{items}
\item

\proofhead
$ \sim $ 是一个等价关系。相应的等价类称为 $ E $ 中的\textbf{分段 $ C^{q} $ 曲线}。分段 $ C^{q} $ 曲线 $ \Gamma $ 的每个代表元都是 $ \Gamma $ 的一个\textbf{分段 $ C^{q} $ 参数化}。当 $ \gamma \in C(I, E) $ 是 $ \Gamma $ 的一个分段 $ C^{q} $ 参数化时，我们符号化地写成 $ \Gamma = \sum_{j=1}^{m}\Gamma_{j} $，其中 $ \Gamma_{j} := [\gamma_{j}] $，$ j = 1, \ldots, m $。
\item

设 $ \Gamma $ 是 $ E $ 中的曲线，其参数化 $ \gamma \in C(I, E) $ 在划分 $ \mathfrak{Z} = (\alpha_{0}, \ldots, \alpha_{m}) $ 上分段 $ C^{q} $。通过下式定义 $ \Gamma $ 的\textbf{长度}（或\textbf{弧长}）：

\[
L(\Gamma) := \operatorname{Var}(\gamma, I) .
\]

\proofhead
$ L(\Gamma) $ 是良定义的，并且
\[
L(\Gamma) = \sum_{j=1}^{m}L(\Gamma_{j}) = \sum_{j=1}^{m}\int_{\alpha_{j-1}}^{\alpha_{j}}\left|\dot{\gamma}_{j}(t)\right|\,dt .
\]

（提示：注 1.5(c) 与引理 1.1。）

\end{items}

\item

若 $ \Gamma = [\gamma] $ 是平面闭分段 $ C^{q} $ 曲线，且 $ (\alpha_{0}, \ldots, \alpha_{m}) $ 是 $ \gamma $ 的一个划分，则

\[
A(\Gamma) := \frac{1}{2}\sum_{j=1}^{m}\int_{\alpha_{j-1}}^{\alpha_{j}}\det\left[\gamma_{j}(t), \dot{\gamma}_{j}(t)\right]\,dt
\]

是 $ \Gamma $ 所含的\textbf{有向面积}。

\begin{items}
\item

\proofhead
$ A(\Gamma) $ 是良定义的。
\item

设 $ -\infty < \alpha < \beta < \infty $，并假设 $ f \in C^{1}([\alpha, \beta], \mathbb{R}) $ 满足 $ f(\alpha) = f(\beta) = 0 $。令 $ a := \alpha $ 与 $ b := 2\beta - \alpha $，并按如下方式定义 $ \gamma : [a, b] \to \mathbb{R}^{2} $：

\begin{equation*}
\gamma(t) := \begin{cases} (\alpha + \beta - t, f(\alpha + \beta - t)) , & t \in [a, \beta] , \\ (\alpha - \beta + t, 0) , & t \in [\beta, b] . \end{cases}
\end{equation*}

证明
(α) $ \Gamma := [\gamma] $ 是闭分段 $ C^{q} $ 曲线（画出它的草图）。
(β) $ A(\Gamma) = \int_{\alpha}^{\beta}f(t)\,dt $，即有向面积 $ A(\Gamma) $ 等于 $ f $ 的图像与 $ [\alpha, \beta] $ 之间的有向面积（见注 VI.3.3(a)）。

\item

求半轴为 $ a, b > 0 $ 的椭圆的有向面积。使用参数化 $ [0, 2\pi] \to \mathbb{R}^{2} $，$ t \mapsto (a\cos t, b\sin t) $。

\item

道路 $ \gamma : [-\pi/4, 3\pi/4] \to \mathbb{R}^{2} $，其中

\begin{equation*}
t \mapsto \begin{cases} \sqrt{\cos(2t)}\,(\cos t, \sin t) , & t \in [-\pi/4, \pi/4] , \\ \sqrt{\left|\cos(2t)\right|}\,(-\sin t, -\cos t) , & t \in [\pi/4, 3\pi/4] \end{cases}
\end{equation*}

参数化了\textbf{双纽线}。验证 $ \Gamma = [\gamma] $ 是平面分段 $ C^{\infty} $ 曲线，并求 $ A(\Gamma) $。

\bookfigure[0.25]{figures/b16fafa378ad83a3.png}

\end{items}

\item

若 $ \Gamma $ 是一个正方形的边界，求 $ A(\Gamma) $。

\item

设 $ f : \mathbb{R} \to [0, 1] $ 是连续且 2 周期的，其中

\begin{equation*}
f(t) = \begin{cases} 0 , & t \in [0, 1/3] , \\ 1 , & t \in [2/3, 1] . \end{cases}
\end{equation*}

又设

\[
r(t) := \sum_{k=1}^{\infty}2^{-k}f(3^{2k}t) \quad \text{且} \quad \alpha(t) := 2\pi\sum_{k=1}^{\infty}2^{-k}f(3^{2k-1}t) \quad \text{对 } t \in \mathbb{R} .
\]

证明

\begin{items}
\item

$ \gamma : \mathbb{R} \to \mathbb{R}^{2} $，$ t \mapsto r(t)(\cos\alpha(t), \sin\alpha(t)) $ 是连续的；

\item

$ \gamma([0, 1]) = \mathbb{B} $。
（提示：设 $ (r_{0}, \alpha_{0}) \in [0, 1]\times[0, 2\pi] $ 是 $ (x_{0}, y_{0}) \in \mathbb{B}\setminus\{0, 0\} $ 的极坐标。又设 $ \sum_{k=1}^{\infty}g_{k}2^{-k} $ 与 $ \sum_{k=1}^{\infty}h_{k}2^{-k} $ 分别是 $ r_{0} $ 与 $ \alpha_{0}/2\pi $ 作为二进制小数的展开。对于

\begin{equation*}
a_{n} := \begin{cases} g_{k} , & n = 2k , \\ h_{k} , & n = 2k - 1 , \end{cases}
\end{equation*}

\proofhead
$ t_{0} := 2\sum_{n=1}^{\infty}a_{n}3^{-n-1} \in [0, 1] $，对 $ n \in \mathbb{N} $ 有 $ f(3^{n}t_{0}) = a_{n} $，且 $ \gamma(t_{0}) = (x_{0}, y_{0}) $。）
\end{items}

\end{exercises}

\section{$ \mathbb{R}^{n} $ 中的曲线}

在第 VII.6 节中，我们曾允许曲线穿过无穷维向量空间。然而，有限维的情形从经典（微分几何）来看无疑更为熟悉，并且从 $ \mathbb{R}^n $ 的欧几里得结构中获得了额外的结构。现在我们来考察 $ \mathbb{R}^n $ 中曲线的局部性质。

在下文中，设

• $ n \ge 2 $ 且 $ \gamma \in C^1(J, \mathbb{R}^n) $ 是曲线 $ \Gamma $ 的一个正则参数化。

\needspace{3\baselineskip}\subsection*{单位切向量}

对 $ t \in J $，

\[
\mathbf{t}(t) := \mathbf{t}_{\gamma}(t) := (\gamma(t), \dot\gamma(t)/|\dot\gamma(t)|) \in T_{\gamma(t)}\mathbb{R}^n
\]

是 $ \mathbb{R}^n $ 中长度为 $ 1 $ 的一个切向量，称为 $ \Gamma $ 在点 $ \gamma(t) $ 处的\textbf{单位切向量}。

\[
\mathbb{R}\mathbf{t}(t) := (\gamma(t), \mathbb{R}\dot\gamma(t)) \subset T_{\gamma(t)}\mathbb{R}^n
\]

的任一元素都是 $ \Gamma $ 在 $ \gamma(t) $ 处的一条\textbf{切线}。

\begin{remark}

\begin{items}
\item

单位切向量 $ \mathbf{t} $ 在参数变换下不变。确切地说，若 $ \zeta = \gamma \circ \varphi $ 是 $ \gamma $ 的一个重新参数化，则对 $ t = \varphi(s) $ 有 $ \mathbf{t}_{\gamma}(t) = \mathbf{t}_{\zeta}(s) $。因此，谈论 $ \Gamma $ 的单位切向量是有意义的。

\proofhead
这由链式法则和 $ \dot\varphi(s) $ 的正性立即推出。\nobreak\hspace{0.35em}■
\item

根据推论 VII.9.8，每个 $ t_0 \in \mathring{J} $ 都有一个环绕它的 $ J $ 的开子区间 $ J_0 $，使得 $ \Gamma_0 := \operatorname{im}(\gamma_0) $（其中 $ \gamma_0 := \gamma|J_0 $）是 $ \mathbb{R}^n $ 的一维 $ C^1 $ 子流形。通过缩小 $ J_0 $，我们可以假设 $ \Gamma_0 $ 由单个坐标卡 $ (U_0, \varphi_0) $ 描述，其中 $ \gamma_0 = i_{\Gamma_0} \circ \varphi_0^{-1} $。于是曲线 $ \Gamma_0 $ 在 $ p := \gamma(t_0) $ 处的切线显然与\emph{流形} $ \Gamma_0 $ 在点 $ p $ 处的切空间一致，即 $ T_p\Gamma_0 = \mathbb{R}\mathbf{t}(t_0) $。
现在可能出现存在 $ \mathring{J} $ 中的一个 $ t_1 \ne t_0 $ 使得 $ \gamma(t_1) = p = \gamma(t_0) $ 的情况。于是存在一个环绕 $ t_1 $ 的开区间 $ J_1 $，使得 $ \Gamma_1 := \operatorname{im}(\gamma_1) $（其中 $ \gamma_1 := \gamma|J_1 $）是 $ \mathbb{R}^n $ 的一个 $ C^1 $ 子流形，并且满足 $ T_p\Gamma_1 = \mathbb{R}\mathbf{t}(t_1) $。因此，一般地 $ T_p\Gamma_0 = T_p\Gamma_1 $ 不成立，也就是说，一般地 $ \mathbf{t}(t_0) \ne \pm\mathbf{t}(t_1) $。这表明 $ \Gamma $ 可以有多条这样的“二重点”切线。在这种情况下，$ \Gamma $ 不是 $ \mathbb{R}^n $ 的子流形，尽管它的足够小的“片段”是子流形；这里“足够小”是在参数区间中衡量的。
例 VII.9.6(c) 中的帕斯卡蜗线给出一个具体的例子。设 $ \Gamma $ 是 $ \mathbb{R}^2 $ 中由下式参数化的紧致正则曲线

\[
\gamma : [-\pi, \pi] \to (1 + 2\cos t)(\cos t, \sin t) .
\]

令 $ t_0 := \arccos(-1/2) $ 和 $ t_1 := -t_0 $，我们有 $ \gamma(t_0) = \gamma(t_1) = (0, 0) \in \Gamma $。对于相应的单位切向量，我们得到

\[
\mathbf{t}(t_j) = ((0, 0), ((-1)^j, -\sqrt{3})/2) \quad \text{对 } j = 0, 1 .
\]

因此 $ \mathbf{t}(t_0) \ne \mathbf{t}(t_1) $。\nobreak\hspace{0.35em}■

\end{items}

\bookfigure[0.21]{figures/7bc571b83427df34.png}
\end{remark}

\needspace{3\baselineskip}\subsection*{弧长参数化}

如果 $ \gamma : J \to \mathbb{R}^n $ 是 $ \Gamma $ 的一个参数化，使得对每个 $ t \in J $ 都有 $ |\dot\gamma(t)| = 1 $，那么显然 $ L(\Gamma) = \int_J dt $，且区间 $ J $ 的长度等于 $ \Gamma $ 的长度。我们说 $ \gamma $ \textbf{按弧长参数化 $ \Gamma $}。

\begin{proposition}

$ \mathbb{R}^n $ 中每条正则 $ C^1 $ 曲线都可以按弧长参数化。这种参数化在形如 $ s \mapsto s + \mathrm{const} $ 的重新参数化下是唯一的。若 $ \eta : I \to \mathbb{R}^n $ 是一个按弧长的参数化，则对 $ s \in I $ 有 $ \mathbf{t}_{\eta}(s) = (\eta(s), \dot\eta(s)) $。

\proofhead
设 $ \gamma \in C^1(J, \mathbb{R}^n) $ 是 $ \Gamma $ 的一个正则参数化。固定一个 $ a \in J $ 并令
\[
\varphi(t) := \int_a^t |\dot\gamma(\tau)| \, d\tau \quad \text{对 } t \in J \text{ 且 } I := \varphi(J) .
\]

$ \gamma $ 的正则性以及 $ |\cdot| \in C(\mathbb{R}^n \setminus \{0\}, \mathbb{R}) $ 表明 $ \varphi $ 属于 $ C^1(J, \mathbb{R}^n) $。此外，$ \dot\varphi(t) = |\dot\gamma(t)| > 0 $。由注 1.5(b)，$ \varphi $ 是从 $ J $ 到 $ I := \varphi(J) $ 的一个 $ C^1 $ 参数变换。令 $ \eta := \gamma \circ \varphi^{-1} $，由此可得

\[
|\dot\eta(s)| = |\dot\gamma(\varphi^{-1}(s))(\varphi^{-1})^{\cdot}(s)| = \frac{|\dot\gamma(\varphi^{-1}(s))|}{|\dot\varphi(\varphi^{-1}(s))|} = \frac{|\dot\gamma(\varphi^{-1}(s))|}{|\dot\gamma(\varphi^{-1}(s))|} = 1
\]

对 $ s \in I $。因此 $ \eta $ 按弧长参数化 $ \Gamma $。

再设 $ \tilde\eta \in C^1(\tilde{I}, \mathbb{R}^n) $ 也按弧长参数化 $ \Gamma $，且 $ \psi \in \operatorname{Diff}^1(I, \tilde{I}) $ 满足 $ \eta = \tilde\eta \circ \psi $。则

\[
1 = |\dot\eta(s)| = |\dot{\tilde\eta}(\psi(s))\dot\psi(s)| = \dot\psi(s) \quad \text{对 } s \in I ,
\]

并且我们得到 $ \psi(s) = s + \mathrm{const} $。\nobreak\hspace{0.35em}■

\end{proposition}

\begin{remark}

\begin{items}
\item

设 $ \Gamma $ 是一条正则 $ C^2 $ 曲线，且 $ \gamma : I \to \mathbb{R}^n $ 按弧长参数化它。则对 $ s \in I $ 有 $ (\dot\gamma(s)|\ddot\gamma(s)) = 0 $。因此，向量 $ (\gamma(s), \ddot\gamma(s)) \subset T_{\gamma(s)}\mathbb{R}^n $ 正交于切线 $ \mathbb{R}\mathbf{t}(s) \subset T_{\gamma(s)}\mathbb{R}^n $。

\proofhead
由 $ |\dot\gamma(s)|^2 = 1 $，我们得到对 $ s \in I $ 有 $ (|\dot\gamma|^2)^{\cdot}(s) = 2(\dot\gamma(s)|\ddot\gamma(s)) = 0 $。\nobreak\hspace{0.35em}■
\item

按弧长参数化曲线是许多理论工作中的一个便捷手段；例如参见命题 2.12 的证明。然而，对于任何具体的例子，进行这样的参数化都要付出计算弧长积分的代价，而且所得的新参数化可能不是初等函数。

\proofhead
我们考虑正则参数化
\[
\gamma : [0, 2\pi] \to \mathbb{R}^2 , \quad t \mapsto (a\cos t, b\sin t) .
\]

它的像是椭圆 $ (x/a)^2 + (y/b)^2 = 1 $，其中 $ a, b > 0 $。命题 2.2 的证明表明，把 $ \gamma $ 变为弧长参数化的每一个重新参数化都具有如下形式

\[
\varphi(t) := \int_0^t \sqrt{a^2\sin^2 s + b^2\cos^2 s} \, ds + \mathrm{const} \quad \text{对 } 0 \le t \le 2\pi .
\]

可以证明 $ \varphi $ 不是初等函数。\footnote{这个 $ \varphi $ 事实上由\textbf{椭圆积分}给出。相关理论参见例如 [FB95]。} \nobreak\hspace{0.35em}■

\end{items}

\end{remark}

\needspace{3\baselineskip}\subsection*{定向基}

设 $ \mathcal{B} = (b_1, \ldots, b_n) $ 和 $ \mathcal{C} = (c_1, \ldots, c_n) $ 是维数为 $ n $ 的实向量空间 $ E $ 的两个有序基。也用 $ T_{\mathcal{B},\mathcal{C}} = [t_{jk}] $ 表示从 $ \mathcal{B} $ 到 $ \mathcal{C} $ 的变换矩阵，即对 $ 1 \le j \le n $ 有 $ c_j = \sum_{k=1}^{n} t_{jk}b_k $。我们说 $ \mathcal{B} $ 和 $ \mathcal{C} $ 是\textbf{同向}[或反向]\textbf{定向的}，如果 $ \det T_{\mathcal{B},\mathcal{C}} $ 为正[或为负]。

\begin{remark}

\begin{items}
\item

在 $ E $ 的有序基的集合上，指派

\[
\mathcal{B} \sim \mathcal{C} :\Longleftrightarrow \mathcal{B} \text{ 与 } \mathcal{C} \text{ 同向定向}
\]

定义一个等价关系。恰有两个等价类，即 $ E $ 的两个\textbf{定向}。通过选取第一个类并称之为 $ \mathcal{O}r $，我们说 $ (E, \mathcal{O}r) $ 是\textbf{定向的}。这个定向中的每个基都是\textbf{正向定向的}，而相反定向（我们称之为 $ -\mathcal{O}r $）中的每个基都是\textbf{负向定向的}。

\item

一个 $ n $ 维内积空间的两个有序标准正交基 $ \mathcal{B} $ 和 $ \mathcal{C} $ 同向定向，当且仅当 $ T_{\mathcal{B},\mathcal{C}} $ 属于 $ SO(n) $。\footnote{参见习题 VII.9.2 和 VII.9.10。}

\item

$ SO(2) $ 中的元素也称为\textbf{旋转矩阵}，因为对每个 $ T \in SO(2) $，存在唯一的\textbf{旋转角} $ \alpha \in [0, 2\pi) $，使得 $ T $ 有表示

\begin{equation*}
T = \begin{bmatrix} \cos\alpha & -\sin\alpha \\ \sin\alpha & \cos\alpha \end{bmatrix}
\end{equation*}

（见习题 10）。因此，一个二维内积空间的任意两个同向定向的基都通过一个旋转相联系。

\item

若 $ E $ 是一个内积空间且 $ \mathcal{B} $ 是一个标准正交基，则 $ T_{\mathcal{B},\mathcal{C}} = [(c_j|b_k)] $。

\proofhead
这由 $ T_{\mathcal{B},\mathcal{C}} $ 的定义立即推出。\nobreak\hspace{0.35em}■
\item

$ \mathbb{R}^n $ 的一个定向基是\textbf{正向定向的}，如果它的定向与典范基的定向相同。\nobreak\hspace{0.35em}■

\end{items}

\end{remark}

\needspace{3\baselineskip}\subsection*{Frenet $ n $-标架}

在下文中，我们将对 $ \Gamma $ 施加稍强的可微性假设，并始终假设 $ \Gamma $ 是 $ \mathbb{R}^n $ 中 $ n \ge 2 $ 的一条正则 $ C^n $ 曲线。照常，设 $ \gamma \in C^n(J, \mathbb{R}^n) $ 是 $ \Gamma $ 的一个参数化。我们说 $ \Gamma $ 是\textbf{完备的}，如果对每个 $ t \in J $，向量 $ \dot\gamma(t), \ldots, \gamma^{(n-1)}(t) $ 都线性无关。

设 $ \Gamma $ 是 $ \mathbb{R}^n $ 中的一条完备曲线。函数 $ J \to \mathbb{R}^n $ 的一个 $ n $ 元组 $ \mathbf{e} = (\mathbf{e}_1, \ldots, \mathbf{e}_n) $ 称为 $ \Gamma $ 的\textbf{活动 $ n $-标架}或 \textbf{Frenet $ n $-标架}，如果

(FB$ {}_{1} $) 对 $ 1 \le j \le n $ 有 $ \mathbf{e}_j \in C^1(J, \mathbb{R}^n) $；

(FB$ {}_{2} $) 对 $ t \in J $，$ \mathbf{e}(t) $ 是 $ \mathbb{R}^n $ 的一个正向标准正交基；

(FB$ {}_{3} $) 对 $ 1 \le k \le n - 1 $ 和 $ t \in J $ 有 $ \gamma^{(k)}(t) \in \operatorname{span}\{\mathbf{e}_1(t), \ldots, \mathbf{e}_k(t)\} $；

(FB$ {}_{4} $) 对 $ t \in J $ 和 $ k \in \{1, \ldots, n - 1\} $，$ (\dot\gamma(t), \ldots, \gamma^{(k)}(t)) $ 与 $ (\mathbf{e}_1(t), \ldots, \mathbf{e}_k(t)) $ 具有相同的定向（作为 $ \operatorname{span}\{\mathbf{e}_1(t), \ldots, \mathbf{e}_k(t)\} $ 的基）。

\begin{remark}

\begin{items}
\item

由链式法则容易推出这个概念是良定义的，也就是说，在下述意义上与 $ \Gamma $ 的特殊参数化无关：若 $ \tilde\gamma := \gamma \circ \varphi $ 是 $ \gamma $ 的一个重新参数化，且 $ \tilde{\mathbf{e}} $ 是 $ \Gamma $ 的一个活动 $ n $-标架，它在将 $ \gamma $ 替换为 $ \tilde\gamma $ 后满足 (FB$ {}_{3} $) 和 (FB$ {}_{4} $)，则对 $ t = \varphi(s) $ 有 $ \tilde{\mathbf{e}}(s) = \mathbf{e}(t) $。

\item

$ \mathbb{R}^2 $ 中的一条 $ C^2 $ 曲线完备当且仅当它是正则的。

\item

设 $ \Gamma = [\gamma] $ 是 $ \mathbb{R}^2 $ 中的一条正则 $ C^2 $ 曲线。又设 $ \mathbf{e}_1 := \dot\gamma/|\dot\gamma| =: (\mathbf{e}_1^1, \mathbf{e}_1^2) $ 和 $ \mathbf{e}_2 := (-\mathbf{e}_1^2, \mathbf{e}_1^1) $。则 $ (\mathbf{e}_1, \mathbf{e}_2) $ 是 $ \Gamma $ 的一个 Frenet 二标架。

\end{items}

\bookfigure[0.78]{figures/39e2a3f7c8873782.png}

（在这些以及类似的图中，我们把 $ \mathbf{e}_j(t) \in \mathbb{R}^2 $ 等同于 $ (\gamma(t), \mathbf{e}_j(t)) \in T_{\gamma(t)}\Gamma $ 的切部分。）

\proofhead
显然 $ (\mathbf{e}_1(t), \mathbf{e}_2(t)) $ 是 $ \mathbb{R}^2 $ 的一个标准正交基，并且是正向的，因为
\[
\det[\mathbf{e}_1, \mathbf{e}_2] = (\mathbf{e}_1^1)^2 + (\mathbf{e}_1^2)^2 = |\mathbf{e}_1|^2 = 1 .\quad\blacksquare
\]
\end{remark}

\begin{theorem}

$ \mathbb{R}^n $ 中每条完备 $ C^n $ 曲线都有唯一的活动 $ n $-标架。

证明

\begin{subitems}
\item

设 $ \Gamma = [\gamma] $ 是 $ \mathbb{R}^n $ 中的一条完备 $ C^n $ 曲线。可以用格拉姆–施密特正交化过程建立 Frenet $ n $-标架的存在性（例如参见 [Art93, VII.1.22]）。事实上，作为完备曲线，$ \Gamma $ 是正则的。因此对 $ t \in I $，$ \mathbf{e}_1(t) := \dot\gamma(t)/|\dot\gamma(t)| $ 有定义。对 $ k \in \{2, \ldots, n - 1\} $，我们递归地定义 $ \mathbf{e}_k $。利用已经构造好的 $ \mathbf{e}_1, \ldots, \mathbf{e}_{k-1} $，我们令

\begin{equation}
\widetilde{\mathbf{e}}_k := \gamma^{(k)} - \sum_{j=1}^{k-1}(\gamma^{(k)}|\mathbf{e}_j)\mathbf{e}_j \quad \text{且} \quad \mathbf{e}_k := \widetilde{\mathbf{e}}_k/|\widetilde{\mathbf{e}}_k| .\tag{2.1}
\end{equation}

因为

\begin{equation}
\operatorname{span}(\dot\gamma(t), \ldots, \gamma^{(k-1)}(t)) = \operatorname{span}(\mathbf{e}_1(t), \ldots, \mathbf{e}_{k-1}(t))\tag{2.2}
\end{equation}

且 $ \gamma^{(k)}(t) \notin \operatorname{span}(\dot\gamma(t), \ldots, \gamma^{(k-1)}(t)) $，所以对 $ t \in J $ 有 $ \widetilde{\mathbf{e}}_k(t) \ne 0 $，从而 $ \mathbf{e}_k $ 是良定义的。此外，我们有

\begin{equation}
(\mathbf{e}_j(t)|\mathbf{e}_k(t)) = \delta_{jk} \quad \text{且} \quad \gamma^{(k)}(t) \in \operatorname{span}(\mathbf{e}_1(t), \ldots, \mathbf{e}_k(t))\tag{2.3}
\end{equation}

对 $ 1 \le j, k \le n - 1 $ 和 $ t \in J $。

\item

现在设 $ T_k(t) $ 是从 $ (\dot\gamma(t), \ldots, \gamma^{(k)}(t)) $ 到 $ (\mathbf{e}_1(t), \ldots, \mathbf{e}_k(t)) $ 的变换矩阵。由 (2.1)，对 $ k \in \{2, \ldots, n - 1\} $ 我们有递推公式

\begin{equation*}
T_1(t) = |\dot\gamma(t)|^{-1} \quad \text{且} \quad T_k(t) = \begin{bmatrix} T_{k-1}(t) & 0 \\ * & |\widetilde{\mathbf{e}}_k(t)|^{-1} \end{bmatrix} .
\end{equation*}

因此 $ \det T_k(t) > 0 $，这表明 $ (\dot\gamma(t), \ldots, \gamma^{(k)}(t)) $ 与 $ (\mathbf{e}_1(t), \ldots, \mathbf{e}_k(t)) $ 具有相同的定向。

\item

最后，对 $ t \in J $，我们通过解线性方程组

\begin{equation}
(\mathbf{e}_j(t)|x) = 0 \quad \text{对 } j = 1, \ldots, n - 1 \quad \text{且} \quad \det[\mathbf{e}_1(t), \ldots \mathbf{e}_{n-1}(t), x] = 1\tag{2.4}
\end{equation}

（关于 $ x \in \mathbb{R}^n $）来定义 $ \mathbf{e}_n(t) $。因为 $ \{\mathbf{e}_1(t), \ldots \mathbf{e}_{n-1}(t)\}^{\perp} $ 是一维的，所以 (2.4) 有唯一解。由 (2.3) 和 (2.4)，现在容易推出 $ \mathbf{e} = (\mathbf{e}_1, \ldots, \mathbf{e}_n) $ 满足条件 (FB$ {}_{2} $)–(FB$ {}_{4} $)。

\item

剩下要验证 (FB$ {}_{1} $)。由 (2.1) 容易推出对 $ k = 1, \ldots, n - 1 $ 有 $ \mathbf{e}_k \in C^1(J, \mathbb{R}^n) $。于是 (2.4) 和克莱姆法则（参见 [Art93, § I.5]）表明 $ \mathbf{e}_n $ 也属于 $ C^1(J, \mathbb{R}^n) $。总之，我们已经证明了 $ \mathbf{e} $ 是 $ \Gamma $ 的一个 Frenet $ n $-标架。

\item

为了检验 $ \mathbf{e} $ 的唯一性，设 $ (\delta_1, \ldots, \delta_n) $ 是 $ \Gamma $ 的另一个 Frenet $ n $-标架。那么显然有 $ \mathbf{e}_1 = \delta_1 = \dot\gamma/|\dot\gamma| $。此外，由 (FB$ {}_{2} $) 和 (FB$ {}_{3} $) 推出，对每个 $ k \in \{1, \ldots, n - 1\} $，$ \gamma^{(k)} $ 有展开式

\[
\gamma^{(k)} = \sum_{j=1}^{k}(\gamma^{(k)}|\delta_j)\delta_j
\]

于是对 $ 1 \le j \le k - 1 $ 和 $ k \in \{2, \ldots, n - 1\} $ 有 $ \mathbf{e}_j = \delta_j $，所以 (2.1) 蕴含 $ (\gamma^{(k)}|\delta_k)\delta_k = \mathbf{e}_k|\widetilde{\mathbf{e}}_k| $。因此，对每个 $ t \in J $，存在一个实数 $ \alpha(t) $ 使得 $ |\alpha(t)| = 1 $ 且 $ \delta_k(t) = \alpha(t)\mathbf{e}_k(t) $。由于 $ (\delta_1(t), \ldots, \delta_k(t)) $ 与 $ (\mathbf{e}_1(t), \ldots, \mathbf{e}_k(t)) $ 同向定向，可得 $ \alpha(t) = 1 $。于是 $ (\mathbf{e}_1, \ldots, \mathbf{e}_{n-1}) $ 从而 $ \mathbf{e}_n $ 也被唯一确定。\nobreak\hspace{0.35em}■

\end{subitems}

\end{theorem}

\begin{corollary}

设 $ \Gamma = [\gamma] $ 是 $ \mathbb{R}^n $ 中的一条完备 $ C^n $ 曲线，且 $ \mathbf{e} = (\mathbf{e}_1, \ldots, \mathbf{e}_n) $ 是 $ \Gamma $ 的一个 Frenet $ n $-标架。则我们有 \textbf{Frenet 导数公式}

\begin{equation}
\dot{\mathbf{e}}_j = \sum_{k=1}^{n}(\dot{\mathbf{e}}_j|\mathbf{e}_k)\mathbf{e}_k \quad \text{对 } j = 1, \ldots, n .\tag{2.5}
\end{equation}

我们还有关系式

\[
(\dot{\mathbf{e}}_j|\mathbf{e}_k) = -(\mathbf{e}_j|\dot{\mathbf{e}}_k) \quad \text{对 } 1 \le j, k \le n \quad \text{且} \quad (\dot{\mathbf{e}}_j|\mathbf{e}_k) = 0 \quad \text{对 } |k - j| > 1 .
\]

\proofhead
因为对每个 $ t \in J $，$ (\mathbf{e}_1(t), \ldots, \mathbf{e}_n(t)) $ 都是 $ \mathbb{R}^n $ 的一个标准正交基，所以我们得到 (2.5)。对 $ (\mathbf{e}_j|\mathbf{e}_k) = \delta_{jk} $ 求导蕴含 $ (\dot{\mathbf{e}}_j|\mathbf{e}_k) = -(\mathbf{e}_j|\dot{\mathbf{e}}_k) $。最后，(2.1) 和 (2.2) 表明对 $ 1 \le j < n $，$ \mathbf{e}_j(t) $ 属于 $ \operatorname{span}(\dot\gamma(t), \ldots, \gamma^{(j)}(t)) $。因此，
\[
\dot{\mathbf{e}}_j(t) \in \operatorname{span}(\dot\gamma(t), \ldots, \gamma^{(j+1)}(t)) = \operatorname{span}(\mathbf{e}_1(t), \ldots, \mathbf{e}_{j+1}(t)) \quad \text{对 } t \in J
\]

且 $ 1 \le j < n - 1 $，于是我们得到对 $ k > j + 1 $ 从而对 $ |k - j| > 1 $ 有 $ (\dot{\mathbf{e}}_j|\mathbf{e}_k) = 0 $。\nobreak\hspace{0.35em}■

\end{corollary}

\begin{remark}

设 $ \Gamma = [\gamma] $ 是 $ \mathbb{R}^n $ 中的一条完备 $ C^n $ 曲线。

\begin{items}
\item

曲率

\[
\kappa_j := (\dot{\mathbf{e}}_j|\mathbf{e}_{j+1})/|\dot\gamma| \in C(J) \quad \text{满足 } 1 \le j \le n - 1 ,
\]

是良定义的，也就是说，它与参数化 $ \gamma $ 无关。

\proofhead
这是注 2.5(a) 和链式法则的一个简单推论。\nobreak\hspace{0.35em}■
\item

令 $ \omega_{jk} := (\dot{\mathbf{e}}_j|\mathbf{e}_k) $，则 Frenet 导数公式读作

\[
\dot{\mathbf{e}}_j = \sum_{k=1}^{n}\omega_{jk}\mathbf{e}_k \quad \text{对 } 1 \le j \le n ,
\]

其中矩阵 $ [\omega_{jk}] $ 为

\begin{equation*}
\left|\dot\gamma\right| \begin{bmatrix} 0 & \kappa_1 & & & & \\ -\kappa_1 & 0 & \kappa_2 & & & \\ & -\kappa_2 & 0 & \kappa_3 & & 0 \\ & & \ddots & \ddots & \ddots & \\ 0 & & & -\kappa_{n-2} & 0 & \kappa_{n-1} \\ & & & & -\kappa_{n-1} & 0 \end{bmatrix} .
\end{equation*}

\proofhead
这由推论 2.7 和 (a) 推出。\nobreak\hspace{0.35em}■
\end{items}

在本节的剩余部分，我们将更详细地讨论定理 2.6 和推论 2.7 对平面曲线和空间曲线的意义。

\end{remark}

\needspace{3\baselineskip}\subsection*{平面曲线的曲率}

在下文中，$ \Gamma = [\gamma] $ 是一条平面正则 $ C^{2} $ 曲线，$ (\mathrm{e}_{1}, \mathrm{e}_{2}) $ 是其相伴的活动二标架。于是 $ \mathrm{e}_{1} $ 与单位切向量 $ \mathrm{t} = (\gamma, \dot{\gamma}/|\dot{\gamma}|) $ 的切部分重合。对每个 $ t \in J $，

\[
\mathrm{n}(t) := (\gamma(t), \mathrm{e}_{2}(t)) \in T_{\gamma(t)}\mathbb{R}^{2}
\]

是 $ \Gamma $ 在点 $ \gamma(t) $ 处的\textbf{单位法向量}。在平面情形，只有 $ \kappa_{1} $ 有定义，并称之为曲线 $ \Gamma $ 的\textbf{曲率} $ \kappa $。因此

\begin{equation}
\kappa = (\dot{\mathrm{e}}_{1}|\mathrm{e}_{2})/|\dot{\gamma}| .\tag{2.6}
\end{equation}

\begin{remark}

\begin{items}
\item

Frenet 导数公式为\footnote{自然地，$ \dot{\mathrm{t}}(t) $ 应理解为向量 $ (\gamma(t), \dot{\mathrm{e}}_{1}(t)) \in T_{\gamma(t)}\mathbb{R}^{2} $，等等。}

\begin{equation}
\dot{\mathrm{t}} = |\dot{\gamma}|\kappa\mathrm{n} \quad \text{且} \quad \dot{\mathrm{n}} = -|\dot{\gamma}|\kappa\mathrm{t} .\tag{2.7}
\end{equation}

若 $ \eta \in C^{2}(I, \mathbb{R}^{2}) $ 按弧长参数化 $ \Gamma $，则方程 (2.7) 取如下简单形式

\begin{equation}
\dot{\mathrm{t}} = \kappa\mathrm{n} \quad \text{且} \quad \dot{\mathrm{n}} = -\kappa\mathrm{t} .\tag{2.8}
\end{equation}

此外，此时还有

\[
\kappa(s) = (\dot{\mathrm{e}}_{1}(s)|\mathrm{e}_{2}(s)) = (\dot{\mathrm{t}}(s)|\mathrm{n}(s))_{\eta(s)} \qquad \text{对 } s \in I .
\]

由 (2.8) 得 $ \ddot{\eta}(s) = \kappa(s)\mathrm{e}_{2}(s) $，从而对 $ s \in I $ 有 $ |\ddot{\eta}(s)| = |\kappa(s)| $。据此，$ \kappa(s)\mathrm{n}(s) \in T_{\eta(s)}\mathbb{R}^{2} $ 称为 $ \Gamma $ 在点 $ \eta(s) $ 处的\textbf{曲率向量}。

\end{items}

\bookfigure[0.48]{figures/c0507ec154b3fcac.png}

在 $ \kappa(s) > 0 $ [或 $ \kappa(s) < 0 $] 的情形，单位切向量 $ \mathrm{t}(s) $ 沿正 [或负] 方向旋转（因为 $ (\dot{\mathrm{t}}(s)|\mathrm{n}(s))_{\eta(s)} $ 是 $ \mathrm{t}(s) $ 的瞬时变化在法向量上的投影）。我们说单位法向量 $ \mathrm{n}(s) $ 指向 $ \Gamma $ 的\textbf{凸} [或\textbf{凹}] \textbf{侧}。

\begin{items}
\item[（b）]

当 $ \eta : I \to \mathbb{R}^{2} $，$ s \mapsto (x(s), y(s)) $ 按弧长参数化 $ \Gamma $ 时，$ \mathrm{e}_{2} = (-\dot{y}, \dot{x}) $ 且

\[
\kappa = \dot{x}\ddot{y} - \dot{y}\ddot{x} = \det[\dot{\eta}, \ddot{\eta}] .
\]

\proofhead
这由注 2.5(c) 和 (2.6) 得出。\nobreak\hspace{0.35em}■
\end{items}

对具体计算而言，知道在给定任意正则参数化时如何求曲率是有用的。

\end{remark}

\begin{proposition}

当 $ \gamma : J \to \mathbb{R}^{2} $，$ t \mapsto (x(t), y(t)) $ 是一条平面 $ C^{2} $ 曲线的正则参数化时，

\[
\kappa = \frac{\dot{x}\ddot{y} - \dot{y}\ddot{x}}{\left((\dot{x})^{2} + (\dot{y})^{2}\right)^{3/2}} = \frac{\det[\dot{\gamma}, \ddot{\gamma}]}{|\dot{\gamma}|^{3}} .
\]

\proofhead
由 $ \mathrm{e}_{1} = \dot{\gamma}/|\dot{\gamma}| $ 得
\[
\dot{\mathrm{e}}_{1} = \frac{\ddot{\gamma}}{|\dot{\gamma}|} - \frac{\dot{x}\ddot{x} + \dot{y}\ddot{y}}{|\dot{\gamma}|^{2}}\,\frac{\dot{\gamma}}{|\dot{\gamma}|} = \frac{\ddot{\gamma}}{|\dot{\gamma}|} - \frac{\dot{x}\ddot{x} + \dot{y}\ddot{y}}{|\dot{\gamma}|^{2}}\mathrm{e}_{1} .
\]

根据注 2.5(c)，有 $ \mathrm{e}_{2} = (-\dot{y}, \dot{x})/|\dot{\gamma}| $，于是定理由 (2.6) 得出。\nobreak\hspace{0.35em}■

\end{proposition}

\begin{example}

（图像的曲率）设 $ f \in C^{2}(J, \mathbb{R}) $。则对按 $ J \to \mathbb{R}^{2} $，$ x \mapsto (x, f(x)) $ 参数化的曲线，曲率为

\[
\kappa = \frac{f''}{\left(1 + (f')^{2}\right)^{3/2}} .
\]

\proofhead
这立即由命题 2.10 得出。\nobreak\hspace{0.35em}■

由推论 IV.2.13 对凸函数的刻画及上述公式可知，$ \operatorname{graph}(f) $ 在点 $ (x, f(x)) $ 处曲率为正当且仅当 $ f $ 在 $ x $ 附近是凸的。因为根据注 2.9(b)，$ \operatorname{graph}(f) $ 上的单位法向量具有正的第二分量，这就解释了注 2.9(a) 中所采用的约定，该约定定义了凸与凹的概念。
\end{example}

\needspace{3\baselineskip}\subsection*{直线与圆的识别}

下面，我们应用 Frenet 公式来证明线段与圆弧可由其曲率刻画。

\begin{proposition}

设 $ \Gamma = [\gamma] $ 是一条平面正则 $ C^{2} $ 曲线。则

\begin{subitems}
\item

$ \Gamma $ 是线段当且仅当对 $ t \in J $ 有 $ \kappa(t) = 0 $；

\item

$ \Gamma $ 是半径为 $ r $ 的圆弧当且仅当对 $ t \in J $ 有 $ |\kappa(t)| = 1/r $。

\proofhead
(i) 显然曲率在每条线段上都为零（见命题 2.10）。
反之，设 $ \kappa = 0 $。我们可以不失一般性地假设 $ \Gamma $ 有弧长参数化 $ \eta $。由注 2.9(a) 得对 $ s \in J $ 有 $ \ddot{\eta}(s) = 0 $。于是定理 IV.3.2 的泰勒公式蕴含存在 $ a, b \in \mathbb{R}^{2} $ 使得 $ \eta(s) = a + bs $。
\end{subitems}

(ii) 设 $ \Gamma $ 是半径为 $ r $ 的圆弧。则存在 $ m \in \mathbb{R}^{2} = \mathbb{C} $ 及一个区间 $ J $ 使得 $ \Gamma $ 由 $ J \to \mathbb{C} $，$ t \mapsto re^{it} + m $ 或 $ t \mapsto re^{-it} + m $ 参数化。断言现在由命题 2.10 得出。

反之，设 $ \kappa = \delta/r $，其中 $ \delta = 1 $ 或 $ \delta = -1 $。我们可以再次假设 $ \eta $ 按弧长参数化 $ \Gamma $。以 $ (\mathrm{e}_{1}, \mathrm{e}_{2}) $ 记 $ \Gamma $ 的活动二标架，由 Frenet 公式 (2.8) 得

\[
\dot{\eta} = \mathrm{e}_{1} , \quad \dot{\mathrm{e}}_{1} = (\delta/r)\mathrm{e}_{2} , \quad \dot{\mathrm{e}}_{2} = -(\delta/r)\mathrm{e}_{1} .
\]

由此得

\[
\left(\eta(s) + (r/\delta)\mathrm{e}_{2}\right)^{\cdot} = \dot{\eta}(s) + (r/\delta)\dot{\mathrm{e}}_{2}(s) = 0 \qquad \text{对 } s \in I .
\]

因此存在 $ m \in \mathbb{R}^{2} $ 使得对 $ s \in I $ 有 $ \eta(s) + (r/\delta)\mathrm{e}_{2}(s) = m $。于是我们得到

\[
|\eta(s) - m| = |r\mathrm{e}_{2}(s)| = r \qquad \text{对 } s \in I ,
\]

这证明 $ \eta $ 参数化了一条半径为 $ r $ 的圆弧。\nobreak\hspace{0.35em}■

\end{proposition}

\needspace{3\baselineskip}\subsection*{沿曲线的瞬时圆}

设 $ \Gamma = [\gamma] $ 是一条平面正则 $ C^{2} $ 曲线。若在 $ t_{0} \in J $ 处 $ \kappa(t_{0}) \ne 0 $，我们称

\[
r(t_{0}) := 1/|\kappa(t_{0})|
\]

为 $ \Gamma $ 在点 $ \gamma(t_{0}) $ 处的\textbf{瞬时半径}，并称

\[
m(t_{0}) := \gamma(t_{0}) + r(t_{0})\mathrm{e}_{2}(t_{0}) \in \mathbb{R}^{2}
\]

为 $ \Gamma $ 在点 $ \gamma(t_{0}) $ 处的\textbf{瞬时中心}。$ \mathbb{R}^{2} $ 中以 $ m(t_{0}) $ 为心、以 $ r(t_{0}) $ 为半径的圆是 $ \Gamma $ 在 $ \gamma(t_{0}) $ 处的\textbf{瞬时密切圆}。

\bookfigure[0.32]{figures/384cee7c34eb705f.png}

\begin{remark}

\begin{items}
\item

密切圆在 $ \gamma(t_{0}) $ 处有参数化

\[
[0, 2\pi] \to \mathbb{R}^{2} , \quad t \mapsto m(t_{0}) + |\kappa(t_{0})|^{-1}(\cos t, \sin t) .
\]
\item

密切圆是唯一在 $ \gamma(t_{0}) $ 处至少二阶与 $ \Gamma $ 相切的圆。

\proofhead
我们可以假设 $ \eta $ 是 $ \Gamma $ 的弧长参数化且 $ \eta(s_{0}) = \gamma(t_{0}) $。再假设 $ (a, b) $ 是 $ \mathbb{R}^{2} $ 的正定向的规范正交基。最后，令
\[
\tilde{\gamma} : [s_{0}, s_{0} + 2\pi r] \to \mathbb{R}^{2} , \quad s \mapsto \tilde{\gamma}(s)
\]

其中

\[
\tilde{\gamma}(s) := m + r\cos\left((s - s_{0})/r\right)a + r\sin\left((s - s_{0})/r\right)b
\]

按弧长参数化以 $ m $ 为心、以 $ r $ 为半径的圆 $ K $。则当对 $ 0 \le j \le 2 $ 有 $ \eta^{(j)}(s_{0}) = \tilde{\gamma}^{(j)}(s_{0}) $ 时，或等价地，当诸方程

\[
\eta(s_{0}) = m + ra , \quad \dot{\eta}(s_{0}) = b , \quad \ddot{\eta}(s_{0}) = -a/r
\]

成立时，$ \Gamma $ 与 $ K $ 在 $ \eta(s_{0}) $ 处至少二阶相切。因为 $ \mathrm{e}_{1} = \dot{\eta} $，第一个 Frenet 公式蕴含 $ \kappa(s_{0})\mathrm{e}_{2}(s_{0}) = -a/r $。因为规范正交基 $ \left(-\mathrm{e}_{2}(s_{0}), \mathrm{e}_{1}(s_{0})\right) $ 是正定向的，我们得到 $ |\kappa(s_{0})| = 1/r $ 且 $ (a, b) = \left(-\mathrm{e}_{2}(s_{0}), \mathrm{e}_{1}(s_{0})\right) $ 且 $ m = \eta(s_{0}) + r\mathrm{e}_{2}(s_{0}) = \gamma(t_{0}) + [\kappa(t_{0})]^{-1}\mathrm{e}_{2}(t_{0}) $。\nobreak\hspace{0.35em}■

\item

当对 $ t \in J $ 有 $ \kappa(t) \ne 0 $ 时，

\[
t \mapsto m(t) := \gamma(t) + \mathrm{e}_{2}(t)/\kappa(t)
\]

是一条平面曲线——$ \Gamma $ 的\textbf{渐屈线}——的连续参数化。\nobreak\hspace{0.35em}■

\end{items}

\bookfigure[0.29]{figures/47c4dc04576f5762.png}
\end{remark}

\needspace{3\baselineskip}\subsection*{向量积}

在线性代数中，$ \mathbb{R}^{3} $ 中除内积外还有“另一种积”。它把另一个向量 $ a \times b $ 指派给两个向量 $ a $ 与 $ b $，即 $ a $ 与 $ b $ 的向量积或叉积。为完整起见，我们回顾其定义方式并汇集其最重要的性质。

对 $ a, b \in \mathbb{R}^{3} $，

\[
\mathbb{R}^{3} \to \mathbb{R} , \quad c \mapsto \det[a, b, c]
\]

是 $ \mathbb{R}^{3} $ 上的一个（连续）线性形式。因此，由里斯表示定理，恰有一个向量，记作 $ a \times b $，使得

\begin{equation}
(a \times b|c) = \det[a, b, c] \qquad \text{对 } c \in \mathbb{R}^{3} .\tag{2.9}
\end{equation}

由此我们定义了\textbf{向量积}或\textbf{叉积}。

\[
\times : \mathbb{R}^{3} \times \mathbb{R}^{3} \to \mathbb{R}^{3} , \quad (a, b) \mapsto a \times b .
\]

\begin{remark}

\begin{items}
\item

叉积是一个\textbf{交错的}双线性映射，即 $ \times $ 是双线性的且斜对称的：

\[
a \times b = -b \times a \qquad \text{对 } a, b \in \mathbb{R}^{3} .
\]

\proofhead
这易由 (2.9) 得出，因为行列式是一个交错的\footnote{若交换一个 $ m $ 线性映射的两个自变量便使其乘以 $ -1 $，则该映射是\textbf{交错的}。}三线性形式。\nobreak\hspace{0.35em}■
\item

对 $ a, b \in \mathbb{R}^{3} $，$ a \times b \ne 0 $ 当且仅当 $ a $ 与 $ b $ 线性无关。若它们线性无关，则 $ (a, b, a \times b) $ 是 $ \mathbb{R}^{3} $ 的一个正定向基。

\proofhead
由 (2.9)，$ a \times b = 0 $ 当且仅当对每个 $ c \in \mathbb{R}^{3} $ 的选择 $ \det[a, b, c] $ 都为零。若 $ a $ 与 $ b $ 线性无关，则通过选取 $ c $ 使得 $ a, b $ 与 $ c $ 线性无关，我们会有 $ \det[a, b, c] \ne 0 $，而由 (2.9) 这与假设 $ a \times b = 0 $ 矛盾。这蕴含第一个断言。第二个断言由 (2.9) 给出
\[
\det[a, b, a \times b] = |a \times b|^{2} \ge 0 .\quad\blacksquare
\]
\item

对 $ a, b \in \mathbb{R}^{3} $，向量 $ a \times b $ 垂直于 $ a $ 与 $ b $。

\proofhead
这由 (2.9) 得出，因为对 $ c \in \{a, b\} $，行列式有两列相同。\nobreak\hspace{0.35em}■
\item

对 $ a, b \in \mathbb{R}^{3} $，叉积具有显式形式

\[
a \times b = (a_{2}b_{3} - a_{3}b_{2}, a_{3}b_{1} - a_{1}b_{3}, a_{1}b_{2} - a_{2}b_{1}) .
\]

\proofhead
这同样由 (2.9) 按最后一列展开行列式得出。\nobreak\hspace{0.35em}■
\item

对 $ a, b \in \mathbb{R}^{3}\setminus\{0\} $，$ a $ 与 $ b $ 之间的（无向）\textbf{夹角} $ \varphi := \angle(a, b) \in [0, \pi] $ 由 $ \cos\varphi := (a|b)/|a|\,|b| $\footnote{这个角可以在任何内积空间而不只是在 $ \mathbb{R}^{3} $ 中定义。还有 $ \angle(a, b) = \angle(b, a) $。}定义。于是我们有

\[
|a \times b| = \sqrt{|a|^{2}|b|^{2} - (a|b)^{2}} = |a|\,|b|\sin\varphi .
\]

这意味着 $ |a \times b| $ 恰是 $ a $ 与 $ b $ 张成的平行四边形的（无向）面积。

\end{items}

\bookfigure[0.30]{figures/03aee5ed17700aab.png}

\proofhead
由于柯西-施瓦茨不等式，$ \angle(a, b) \in [0, \pi] $ 是良定义的。第一个等式可直接由 (d) 验证，第二个等式由 $ \cos^{2}\varphi + \sin^{2}\varphi = 1 $ 与 $ \sin\varphi \ge 0 $ 得出。\nobreak\hspace{0.35em}■

(f) 当 $ a $ 与 $ b $ 是 $ \mathbb{R}^{3} $ 中的规范正交向量时，（有序）组 $ (a, b, a \times b) $ 是 $ \mathbb{R}^{3} $ 上的一个正定向的规范正交基。

\proofhead
这由 (b) 与 (e) 得出。\nobreak\hspace{0.35em}■
\end{remark}

\needspace{3\baselineskip}\subsection*{空间曲线的曲率与挠率}

现在在三维中工作，设 $ \Gamma = [\gamma] $ 是 $ \mathbb{R}^{3} $ 中一条完备的 $ C^{3} $ 曲线。则 $ \Gamma $ 有两个曲率 $ \kappa_{1} $ 与 $ \kappa_{2} $。我们称 $ \kappa := \kappa_{1} $ 为\textbf{曲率}，称 $ \tau := \kappa_{2} $ 为\textbf{挠率}。因此

\begin{equation}
\kappa = (\dot{\mathrm{e}}_{1}|\mathrm{e}_{2})/|\dot{\gamma}| \quad \text{且} \quad \tau = (\dot{\mathrm{e}}_{2}|\mathrm{e}_{3})/|\dot{\gamma}| .\tag{2.10}
\end{equation}

这里我们也称

\[
\mathrm{n}(t) := (\gamma(t), \mathrm{e}_{2}(t)) \in T_{\gamma(t)}\mathbb{R}^{3}
\]

为\textbf{单位法向量}，而

\[
\mathrm{b}(t) := (\gamma(t), \mathrm{e}_{3}(t)) \in T_{\gamma(t)}\mathbb{R}^{3}
\]

是 $ \Gamma $ 在点 $ \gamma(t) $ 处的\textbf{单位副法向量}。向量 $ \mathrm{t}(t) $ 与 $ \mathrm{n}(t) $ 张成 $ \Gamma $ 在点 $ \gamma(t) $ 处的\textbf{密切平面}。由 $ \mathrm{n}(t) $ 与 $ \mathrm{b}(t) $ 张成的平面是 $ \Gamma $ 在 $ \gamma(t) $ 处的\textbf{法平面}。

\begin{remark}

\begin{items}
\item

我们有 $ \mathrm{e}_{3} = \mathrm{e}_{1} \times \mathrm{e}_{2} $，即 $ \mathrm{b} = \mathrm{t} \times \mathrm{n} $。

\item

当 $ \Gamma $ 按弧长参数化时，Frenet 公式给出

\[
\dot{\mathrm{e}}_{1} = \kappa\mathrm{e}_{2} , \quad \dot{\mathrm{e}}_{2} = -\kappa\mathrm{e}_{1} + \tau\mathrm{e}_{3} , \quad \dot{\mathrm{e}}_{3} = -\tau\mathrm{e}_{2} ,
\]

且因此

\[
\dot{\mathrm{t}} = \kappa\mathrm{n} , \quad \dot{\mathrm{n}} = -\kappa\mathrm{t} + \tau\mathrm{b} , \quad \dot{\mathrm{b}} = -\tau\mathrm{n} .
\]
\item

完备曲线的曲率恒为正。当 $ \eta $ 是弧长参数化时，

\[
\kappa = |\ddot{\eta}| = |\dot{\eta} \times \ddot{\eta}| .
\]

\proofhead
因为 $ \mathrm{e}_{1} = \dot{\eta} $ 且因为 (2.1) 与注 2.3(a) 蕴含 $ \mathrm{e}_{2} = \ddot{\eta}/|\ddot{\eta}| $，我们得到
\[
\kappa = (\dot{\mathrm{e}}_{1}|\mathrm{e}_{2}) = (\ddot{\eta}|\ddot{\eta}/|\ddot{\eta}|) = |\ddot{\eta}| > 0 .
\]

所断言的第二个等式现在由注 2.14(e) 得出。\nobreak\hspace{0.35em}■

\item

当 $ \eta $ 是 $ \Gamma $ 的弧长参数化时，我们有 $ \tau = \det[\dot{\eta}, \ddot{\eta}, \dddot{\eta}]/\kappa^{2} $。

\proofhead
因为 $ \ddot{\eta} = \mathrm{e}_{2}|\ddot{\eta}| = \kappa\mathrm{e}_{2} $，且由 $ \mathrm{e}_{2} $ 的 Frenet 导数公式，我们得到 $ \dddot{\eta} = \dot{\kappa}\mathrm{e}_{2} - \kappa^{2}\mathrm{e}_{1} + \kappa\tau\mathrm{e}_{3} $。因此 $ \det[\dot{\eta}, \ddot{\eta}, \dddot{\eta}] = \det[\mathrm{e}_{1}, \kappa\mathrm{e}_{2}, \kappa\tau\mathrm{e}_{3}] = \kappa^{2}\tau $。注意到 (c) 连同 $ \Gamma $ 的完备性蕴含 $ \kappa $ 处处非零，证明便告完成。\nobreak\hspace{0.35em}■
\end{items}

挠率 $ \tau(t) $ 度量了曲线 $ \Gamma = [\gamma] $ 在 $ \gamma(t) $ 附近“绕出”密切平面的程度，密切平面在右侧以较浅的灰色示出。事实上，平面曲线由挠率为零来刻画，如以下定理所示。

\bookfigure[0.26]{figures/04501254b2847d30.png}
\end{remark}

\begin{proposition}

一条完备曲线位于一个平面中当且仅当其挠率处处为零。

\proofhead
我们可以假设 $ \eta $ 按弧长参数化 $ \Gamma $。若对每个 $ s \in I $，$ \eta(s) $ 都位于一个平面中，则 $ \dot{\eta}(s) $、$ \ddot{\eta}(s) $ 与 $ \dddot{\eta}(s) $ 也如此。由注 2.15(d)，这蕴含 $ \tau = \det[\dot{\eta}, \ddot{\eta}, \dddot{\eta}]/\kappa^{2} = 0 $。反之，若 $ \tau = 0 $，则第三个 Frenet 导数公式给出对每个 $ s \in I $ 及某个 $ \alpha \in I $ 有 $ \mathrm{e}_{3}(s) = \mathrm{e}_{3}(\alpha) $。因为 $ \mathrm{e}_{1}(s) \perp \mathrm{e}_{3}(s) $，
\[
\left(\eta(s)|\mathrm{e}_{3}(\alpha)\right)^{\cdot} = (\dot{\eta}(s)|\mathrm{e}_{3}(\alpha)) = (\mathrm{e}_{1}(s)|\mathrm{e}_{3}(s)) = 0 \qquad \text{对 } s \in I .
\]

因此 $ (\eta(\cdot)|\mathrm{e}_{3}(\alpha)) $ 在 $ I $ 上为常数，这意味着 $ \eta(s) $ 到直线 $ \mathbb{R}\mathrm{e}_{3}(\alpha) $ 上的正交投影与 $ s \in I $ 无关。因此 $ \Gamma $ 位于与 $ \mathrm{e}_{3}(\alpha) $ 正交的平面中。\nobreak\hspace{0.35em}■

\end{proposition}

\exerciseshead

\begin{exercises}
\item

设 $ r \in C^{2}(J, \mathbb{R}^{+}) $ 满足对 $ t \in J $ 有 $ [r(t)]^{2} + [\dot r(t)]^{2} > 0 $，且 $ \gamma(t) := r(t)(\cos t, \sin t) $。证明由 $ \gamma $ 参数化的曲线具有曲率

\[
\kappa = \frac{2[\dot r]^{2} - r\ddot r + r^{2}}{(r^{2} + [\dot r]^{2})^{3/2}} .
\]
\item

计算下列曲线的曲率

\begin{equation*}
\begin{aligned}
\text{帕斯卡蜗线}\quad &[-\pi,\pi] \to \mathbb{R}^{2} , & t &\mapsto (1+2\cos t)(\cos t,\sin t) ; \\ \text{对数螺线}\quad &\mathbb{R} \to \mathbb{R}^{2} , & t &\mapsto e^{\lambda t}(\cos t,\sin t) ; \\ \text{摆线}\quad &[0,2\pi] \to \mathbb{R}^{2} , & t &\mapsto R(t-\sin t,1-\cos t) .
\end{aligned}
\end{equation*}
\item

设 $ \gamma : J \to \mathbb{R}^{3} $, $ t \mapsto (x(t), y(t)) $ 是平面曲线 $ \Gamma $ 的正则 $ C^{2} $ 参数化，且对 $ t \in J $ 有 $ \kappa(t) \ne 0 $。证明 $ \Gamma $ 的渐屈线由 $ t \mapsto m(t) $ 参数化，其中

\[
m := \left( x - \frac{\dot x^{2} + \dot y^{2}}{\dot x\ddot y - \ddot x\dot y}\dot y ,\ y + \frac{\dot x^{2} + \dot y^{2}}{\dot x\ddot y - \ddot x\dot x}\dot x \right) .
\]
\item

证明

\begin{items}
\item

抛物线 $ y = x^{2}/2 $ 的渐屈线是尼尔抛物线，即半立方抛物线 $ \mathbb{R} \to \mathbb{R}^{2} $, $ t \mapsto (-t^{3}, 1 + 3t^{2}/2) $；

\item

对数螺线 $ \mathbb{R} \to \mathbb{R}^{2} $, $ t \mapsto e^{\lambda t}(\cos t, \sin t) $ 的渐屈线是对数螺线 $ \mathbb{R} \to \mathbb{R}^{2} $, $ t \mapsto \lambda e^{\lambda(t-\pi/2)}(\cos t, \sin t) $；

\bookfigure[0.88]{figures/9d7d0f8533807713.png}

\item

摆线 $ \mathbb{R} \to \mathbb{R}^{2} $, $ t \mapsto (t - \sin t, 1 - \cos t) $ 的渐屈线是摆线 $ \mathbb{R} \to \mathbb{R}^{2} $, $ t \mapsto (t + \sin t, \cos t - 1) $。

\bookfigure[0.52]{figures/9555e65074fd22a8.png}

\end{items}

\item

设 $ \gamma \in C^{3}(J, \mathbb{R}^{2}) $ 是平面曲线 $ \Gamma $ 的正则参数化，并假设对 $ t \in J $ 有 $ \kappa(t)\dot\kappa(t) \ne 0 $。用 $ m $ 表示注 2.13(c) 中给出的 $ \Gamma $ 的渐屈线 $ M $ 的参数化。证明

\begin{items}
\item

$ M $ 是正则 $ C^{1} $ 曲线；

\item

$ \Gamma $ 在 $ \gamma(t) $ 处的切线与 $ M $ 在 $ m(t) $ 处的切线垂直；

\item

$ M $ 的弧长由下式给出

\[
L(M) = \int_{J} |\dot\kappa_{\gamma}|\,\kappa_{\gamma}^{-2}\,dt .
\]

（(b) 与 (c) 的提示：Frenet 导数公式。）

\end{items}

\item

证明对数螺线 $ \Gamma $ 以常角穿过从原点出发的直线，即沿螺线的任一点处，$ \Gamma $ 的切线与从原点出发的直线之间的夹角都相同。

\item

计算曲率与挠率

\begin{items}
\item

对于\textbf{椭圆螺旋线} $ \mathbb{R} \to \mathbb{R}^{3} $, $ t \mapsto (a\cos t, b\sin t, ct) $，其中 $ a, b, c \in \mathbb{R} $ 满足 $ ab \ne 0 $；

\item

对于曲线 $ \mathbb{R} \to \mathbb{R}^{3} $, $ t \mapsto t(\cos t, \sin t, 1) $。

\end{items}

\item

设 $ \Gamma $ 是平面上的正则闭 $ C^{1} $ 曲线。证明\textbf{等周不等式}

\[
4\pi A(\Gamma) \le [L(\Gamma)]^{2}
\]

（参见习题 1.6）。

（提示：

\begin{subitems}
\item

将 $ \mathbb{R}^{2} $ 与 $ \mathbb{C} $ 等同起来，先考虑 $ L(\Gamma) = 2\pi $ 的情形。然后不失一般性，设 $ \gamma \in C^{1}([0,2\pi], \mathbb{C}) $ 按弧长参数化 $ \Gamma $。利用帕塞瓦尔方程（推论 VI.7.16）与注 VI.7.20(b)，证明

\[
1 = \frac{1}{2\pi}\int_{0}^{2\pi}|\dot\gamma|^{2}\,dt = \sum_{k=-\infty}^{\infty} k^{2}|\hat\gamma_{k}|^{2} .
\]

进一步证明

\[
A(\Gamma) = \frac{1}{2}\int_{0}^{2\pi} \operatorname{Im}\bar\gamma\dot\gamma\,dt .
\]

于是由习题 VI.7.11 与注 VI.7.20(b)，可见

\[
A(\Gamma) = \pi \sum_{k=-\infty}^{\infty} k|\hat\gamma_{k}|^{2} ,
\]

从而

\[
A(\Gamma) \le \pi \sum_{k=-\infty}^{\infty} k^{2}|\hat\gamma_{k}|^{2} = \pi = [L(\Gamma)]^{2}/4\pi .
\]
\item

$ L(\Gamma) \ne 2\pi $ 的情形通过变换 $ \gamma \to 2\pi\gamma/L(\Gamma) $ 化归为 (i)。）

\end{subitems}

\item

设 $ I $ 与 $ J $ 是紧致区间，$ U $ 是 $ I \times J $ 在 $ \mathbb{R}^{2} $ 中的开邻域。进一步假设 $ \gamma \in C^{3}(U, \mathbb{R}^{2}) $ 满足条件

\begin{subitems}
\item

对每个 $ \lambda \in I $，函数 $ \gamma_{\lambda} := \gamma(\lambda,\cdot)|J : J \to \mathbb{R}^{2} $ 是 $ C^{3} $ 曲线 $ \Gamma_{\lambda} $ 的正则参数化；

\item

对 $ (\lambda,t) \in I \times J $ 有 $ \partial_{1}\gamma(\lambda,t) = \kappa_{\Gamma_{\lambda}}(t)\mathbf{n}_{\Gamma_{\lambda}}(t) $。
证明函数

\[
\ell : I \to \mathbb{R} ,\quad \lambda \mapsto L(\Gamma_{\lambda})
\]

有连续导数，并且

\[
\ell'(\lambda) = -\int_{0}^{\ell(\lambda)} [\kappa(s(\lambda))]^{2}ds(\lambda) \quad \text{对 } \lambda \in I ,
\]

其中对 $ \lambda \in I $，$ s(\lambda) $ 是 $ \Gamma_{\lambda} $ 的弧长参数。
（提示：定义 $ v : I \times J \to \mathbb{R} $, $ (\lambda,t) \mapsto |\partial_{1}\gamma(\lambda,t)| $。利用 Frenet 公式与 (ii)，推出 $ \partial_{1}v = -\kappa^{2}v $。）

\end{subitems}

\item

证明对每个 $ T \in SO(2) $，存在唯一的 $ \alpha \in [0, 2\pi) $，使得

\begin{equation*}
T = \begin{bmatrix} \cos\alpha & -\sin\alpha \\ \sin\alpha & \cos\alpha \end{bmatrix} .
\end{equation*}

（提示：习题 VII.9.2。）

\end{exercises}

\section{普法夫形式}

在第 VI.5 节中，我们以一种特设的方式引入了微分，以便能够更简单地表示积分法则。现在，我们将为使用这些此前只是形式化的对象发展出一套演算。这套演算会自然地容纳之前那些简单的计算法则，同时也会带来更多的清晰与理解。与此同时，这一发展将把积分的概念从在区间上进行的某种运算，扩展到可以在曲线上进行的某种运算。

在下面，设

• $ X $ 是 $ \mathbb{R}^{n} $ 中非空的开集；并设 $ q \in \mathbb{N} \cup \{\infty\} $。

\needspace{3\baselineskip}\subsection*{向量场与普法夫形式}

满足对 $ p \in X $ 有 $ \boldsymbol{v}(p) \in T_{p}X $ 的映射 $ \boldsymbol{v} : X \to TX $ 称为 $ X $ 上的\textbf{向量场}。对每个向量场 $ \boldsymbol{v} $，对应唯一一个函数 $ v : X \to \mathbb{R}^{n} $，即 $ \boldsymbol{v} $ 的\textbf{切部分}，使得 $ \boldsymbol{v}(p) = (p,v(p)) $。借助自然投影

\[
\mathrm{pr}_{2} : TX = X \times \mathbb{R}^{n} \to \mathbb{R}^{n} , \quad (p,\xi) \mapsto \xi ,
\]

有 $ v = \mathrm{pr}_{2} \circ \boldsymbol{v} $。

若 $ v \in C^{q}(X,\mathbb{R}^{n}) $，则称向量场 $ \boldsymbol{v} $ 属于类 $ C^{q} $。我们用 $ \mathcal{V}^{q}(X) $ 表示 $ X $ 上所有类 $ C^{q} $ 的向量场之集。

设 $ E $ 是 $ \mathbb{K} $ 上的巴拿赫空间，$ E' := \mathcal{L}(E,\mathbb{K}) $ 是其对偶空间。映射

\[
\langle\cdot,\cdot\rangle : E' \times E \to \mathbb{K} , \quad (e',e) \mapsto \langle e',e\rangle := e'(e) ,
\]

把每个二元组 $ (e',e) \in E' \times E $ 映为 $ e' $ 在位置 $ e $ 处的值，称为 $ E $ 与 $ E' $ 之间的\textbf{对偶配对}。

现在我们来考虑与 $ X $ 的切空间对偶的空间。对 $ p \in X $，与 $ T_{p}X $ 对偶的空间是 $ X $ 在点 $ p $ 处的\textbf{余切空间}，记为 $ T_{p}^{*}X $。又

\[
T^{*}X := \bigcup_{p \in X}T_{p}^{*}X
\]

是 $ X $ 的\textbf{余切丛}。因此，$ T_{p}^{*}X $ 的元素，即 $ p $ 处的\textbf{余切向量}，是 $ X $ 在 $ p $ 处的切空间上的线性形式。相应的对偶配对记为

\[
\langle\cdot,\cdot\rangle_{p} : T_{p}^{*}X \times T_{p}X \to \mathbb{R} .
\]

每个满足对 $ p \in X $ 有 $ \boldsymbol{\alpha}(p) \in T_{p}^{*}X $ 的映射 $ \boldsymbol{\alpha} : X \to T^{*}X $ 称为 $ X $ 上的\textbf{普法夫形式}。有时我们也称 $ \boldsymbol{\alpha} $ 为\textbf{ 1 次微分形式}或 \textbf{1-形式}。

\begin{remark}

\begin{items}
\item

若 $ E $ 是 $ \mathbb{K} $ 上的巴拿赫空间，则 $ \langle\cdot,\cdot\rangle \in \mathcal{L}(E',E;\mathbb{K}) $。

\proofhead
显然对偶配对是双线性的。此外，结论 VI.2.4(a) 给出估计
\[
|\langle e',e\rangle| = |e'(e)| \le \|e'\|_{E'}\|e\|_{E} \quad \text{对 } (e',e) \in E' \times E .
\]

因此，命题 VII.4.1 表明 $ \langle\cdot,\cdot\rangle $ 属于 $ \mathcal{L}(E',E;\mathbb{K}) $。\nobreak\hspace{0.35em}■

\item

对 $ (p,e') \in \{p\} \times (\mathbb{R}^{n})' $，通过

\begin{equation}
\langle J(p,e'),(p,e)\rangle_{p} = \langle e',e\rangle \quad \text{对 } e \in \mathbb{R}^{n} ,\tag{3.1}
\end{equation}

定义 $ J(p,e') \in T_{p}^{*}X $，其中 $ \langle\cdot,\cdot\rangle $ 表示 $ \mathbb{R}^{n} $ 与 $ (\mathbb{R}^{n})' $ 之间的对偶配对。则 $ J $ 是 $ \{p\} \times (\mathbb{R}^{n})' $ 到 $ T_{p}^{*}X $ 上的等距同构。利用这个同构，我们把 $ T_{p}^{*}X $ 与 $ \{p\} \times (\mathbb{R}^{n})' $ 等同。

\proofhead
显然 $ J $ 是线性的。若对某个 $ (p,e') \in \{p\} \times (\mathbb{R}^{n})' $ 有 $ J(p,e') = 0 $，则由 (3.1) 得 $ e' = 0 $。因此 $ J $ 是单射。因为
\[
\dim\left(\{p\} \times (\mathbb{R}^{n})'\right) = \dim\left(\{p\} \times \mathbb{R}^{n}\right) = \dim(T_{p}X) = \dim(T_{p}^{*}X) ,
\]

$ J $ 也是满射。于是该断言由定理 VII.1.6 得出。\nobreak\hspace{0.35em}■

\end{items}

因为 $ T_{p}^{*}X = \{p\} \times (\mathbb{R}^{n})' $，所以对每个普法夫形式 $ \boldsymbol{\alpha} $ 存在唯一一个映射 $ \alpha : X \to (\mathbb{R}^{n})' $，即 $ \boldsymbol{\alpha} $ 的\textbf{余切部分}，使得 $ \boldsymbol{\alpha}(p) = (p,\alpha(p)) $。若 $ \alpha \in C^{q}\left(X,(\mathbb{R}^{n})'\right) $，则称 1-形式 $ \boldsymbol{\alpha} $ 属于\textbf{类} $ C^{q} $。我们用 $ \Omega_{(q)}(X) $ 表示 $ X $ 上所有类 $ C^{q} $ 的 1-形式之集。

在 $ \mathcal{V}^{q}(X) $ 上，我们用 $ C^{q}(X) $ 中的函数定义乘法，即

\[
C^{q}(X) \times \mathcal{V}^{q}(X) \to \mathcal{V}^{q}(X) , \quad (a,\boldsymbol{v}) \mapsto a\boldsymbol{v}
\]

其中乘法按逐点方式定义：$ (a\boldsymbol{v})(p) := a(p)\boldsymbol{v}(p) $（$ p \in X $）。同样，在 $ \Omega_{(q)}(X) $ 上，我们用 $ C^{q}(X) $ 中的函数定义乘法，即

\[
C^{q}(X) \times \Omega_{(q)}(X) \to \Omega_{(q)}(X) , \quad (a,\boldsymbol{\alpha}) \mapsto a\boldsymbol{\alpha}
\]

其中乘法按逐点方式定义：$ (a\boldsymbol{\alpha})(p) := a(p)\boldsymbol{\alpha}(p) $（$ p \in X $）。

于是容易验证，$ \mathcal{V}^{q}(X) $ 和 $ \Omega_{(q)}(X) $ 是（交换）环 $ C^{q}(X) $ 上的模\footnote{为完备起见，我们把与模有关的事实放在本节末尾。}。因为 $ C^{q}(X) $ 自然地把 $ \mathbb{R} $ 作为子环包含在内（利用 $ \lambda \in \mathbb{R} $ 与 $ \lambda\mathbf{1} \in C^{q}(X) $ 的等同），所以特别地，$ \mathcal{V}^{q}(X) $ 和 $ \Omega_{(q)}(X) $ 是实向量空间。

\end{remark}

\begin{remark}

\begin{items}
\item

对 $ \boldsymbol{\alpha} \in \Omega_{(q)}(X) $ 和 $ \boldsymbol{v} \in \mathcal{V}^{q}(X) $，

\[
\left(p \mapsto \langle\boldsymbol{\alpha}(p),\boldsymbol{v}(p)\rangle_{p} = \langle\boldsymbol{\alpha}(p),\boldsymbol{v}(p)\rangle\right) \in C^{q}(X) .
\]
\item

在不会引起混淆时，我们把 $ \boldsymbol{\alpha} \in \Omega_{(q)}(X) $ 与其余切部分 $ \alpha = \mathrm{pr}_{2} \circ \boldsymbol{\alpha} $ 等同，并把 $ \boldsymbol{v} \in \mathcal{V}^{q}(X) $ 与其切部分 $ v = \mathrm{pr}_{2} \circ \boldsymbol{v} $ 等同。则

\[
\mathcal{V}^{q}(X) = C^{q}(X,\mathbb{R}^{n}) \quad \text{且} \quad \Omega_{(q)}(X) = C^{q}\left(X,(\mathbb{R}^{n})'\right) .
\]
\item

设 $ f \in C^{q+1}(X) $，且

\[
d_{p}f = \mathrm{pr}_{2} \circ T_{p}f \quad \text{对 } p \in X ,
\]

是 $ f $ 的微分（见第 VII.10 节）。则映射 $ df := (p \mapsto d_{p}f) $ 属于 $ \Omega_{(q)}(X) $，且有 $ d_{p}f = \partial f(p) $。从现在起，我们用 $ df(p) $ 表示 $ d_{p}f $。\nobreak\hspace{0.35em}■

\end{items}

\end{remark}

\needspace{3\baselineskip}\subsection*{典范基}

下面，设 $ (e_{1},\ldots,e_{n}) $ 是 $ \mathbb{R}^{n} $ 的标准基，$ (\varepsilon^{1},\ldots,\varepsilon^{n}) $ 是 $ (\mathbb{R}^{n})' $ 中相应的对偶基\footnote{对偶基的存在性是线性代数的标准结果。}，即\footnote{$ \delta_{k}^{j} := \delta^{jk} := \delta_{jk} $ 是克罗内克符号。}

\[
\langle\varepsilon^{j},e_{k}\rangle = \delta_{k}^{j} \quad \text{对 } j,k \in \{1,\ldots,n\} .
\]

我们还设

\[
dx^{j} := d(\mathrm{pr}_{j}) \in \Omega_{(\infty)}(X) \quad \text{对 } j = 1,\ldots,n ,
\]

其中用到第 $ j $ 个投影 $ \mathrm{pr}_{j} : \mathbb{R}^{n} \to \mathbb{R} $，$ x = (x^{1},\ldots,x^{n}) \mapsto x^{j} $。

\begin{remark}

\begin{items}
\item

由注 3.1(b)，对每个 $ j = 1,\ldots,n $，$ (\varepsilon^{j})_{p} = (p,\varepsilon^{j}) $ 属于 $ T_{p}^{*}X $，且 $ \left((\varepsilon^{1})_{p},\ldots,(\varepsilon^{n})_{p}\right) $ 是与 $ T_{p}X $ 的典范基 $ \left((e_{1})_{p},\ldots,(e_{n})_{p}\right) $ 对偶的基。

\item

有

\[
dx^{j}(p) = (p,\varepsilon^{j}) \quad \text{对 } p \in X \text{ 且 } j = 1,\ldots,n .
\]

\proofhead
该断言由命题 VII.2.8 得出，后者给出
\[
\langle dx^{j}(p),(e_{k})_{p}\rangle_{p} = \langle\partial(\mathrm{pr}_{j})(p),e_{k}\rangle = \partial_{k}\mathrm{pr}_{j}(p) = \delta_{jk} .\quad\blacksquare
\]
\item

对 $ 1 \le j \le n $ 和 $ p \in X $，我们设 $ e_{j}(p) := (p,e_{j}) $。则 $ (e_{1},\ldots,e_{n}) $ 是 $ \mathcal{V}^{q}(X) $ 的模基，$ (dx^{1},\ldots,dx^{n}) $ 是 $ \Omega_{(q)}(X) $ 的模基。此外，有

\begin{equation}
\langle dx^{j}(p),e_{k}(p)\rangle_{p} = \delta_{k}^{j} \quad \text{对 } 1 \le j,k \le n \text{ 且 } p \in X ,\tag{3.2}
\end{equation}

且 $ (e_{1},\ldots,e_{n}) $ 是 $ \mathcal{V}^{q}(X) $ 的\textbf{典范基}。同样，$ (dx^{1},\ldots,dx^{n}) $ 是 $ \Omega_{(q)}(X) $ 的典范基。\footnote{这两组基在模对偶的意义下互为对偶（见 [SS88, § 70]）。}
最后，我们定义

\[
\langle\boldsymbol{\alpha},\boldsymbol{v}\rangle : \Omega_{(q)}(X) \times \mathcal{V}^{q}(X) \to C^{q}(X)
\]

如下

\[
\langle\boldsymbol{\alpha},\boldsymbol{v}\rangle(p) := \langle\boldsymbol{\alpha}(p),\boldsymbol{v}(p)\rangle_{p} \quad \text{对 } p \in X .
\]

则 (3.2) 取如下形式

\[
\langle dx^{j},e_{k}\rangle = \delta_{k}^{j} \quad \text{对 } 1 \le j,k \le n .
\]

\proofhead
设 $ \boldsymbol{\alpha} \in \Omega_{(q)}(X) $。我们设
\[
a_{j}(p) := \langle\boldsymbol{\alpha}(p),e_{j}(p)\rangle_{p} \quad \text{对 } p \in X .
\]

于是，由注 3.2(a)，$ a_{j} $ 属于 $ C^{q}(X) $，因此 $ \boldsymbol{\beta} := \sum a_{j}dx^{j} $ 属于 $ \Omega_{(q)}(X) $。此外，

\[
\langle\boldsymbol{\beta}(p),e_{k}(p)\rangle_{p} = \sum_{j=1}^{n}a_{j}(p)\langle dx^{j}(p),e_{k}(p)\rangle_{p} = a_{k}(p) = \langle\boldsymbol{\alpha}(p),e_{k}(p)\rangle_{p} \quad \text{对 } p \in X
\]

且 $ k = 1,\ldots,n $。因此 $ \boldsymbol{\alpha} = \boldsymbol{\beta} $，这表明 $ (e_{1},\ldots e_{n}) $ 与 $ (dx^{1},\ldots,dx^{n}) $ 是模基，且 (3.2) 成立。\nobreak\hspace{0.35em}■

\item

每个 $ \boldsymbol{\alpha} \in \Omega_{(q)}(X) $ 有\textbf{典范基表示}

\[
\boldsymbol{\alpha} = \sum_{j=1}^{n}\langle\boldsymbol{\alpha},e_{j}\rangle dx^{j} .
\]
\item

对 $ f \in C^{q+1}(X) $，

\[
df = \partial_{1}f\,dx^{1} + \cdots + \partial_{n}f\,dx^{n} .
\]

映射 $ d : C^{q+1}(X) \to \Omega_{(q)}(X) $ 是 $ \mathbb{R} $-线性的。

\proofhead
因为
\[
\langle df(p),e_{j}(p)\rangle_{p} = \partial_{j}f(p) \quad \text{对 } p \in X ,
\]

第一个断言由 (d) 得出。第二个断言是显然的。\nobreak\hspace{0.35em}■

\item

$ \mathcal{V}^{q}(X) $ 和 $ \Omega_{(q)}(X) $ 是 $ \mathbb{R} $ 上的无穷维向量空间。

\proofhead
我们考虑 $ X = \mathbb{R} $，一般情形留给读者。由代数基本定理容易得出，$ \mathcal{V}^{q}(\mathbb{R}) = C^{q}(\mathbb{R}) $ 中的单项式 $ \{X^{m} \;;\; m \in \mathbb{N}\} $ 在 $ \mathbb{R} $ 上线性无关。因此 $ \mathcal{V}^{q}(\mathbb{R}) $ 是无穷维向量空间。我们已经知道 $ \Omega_{(q)}(\mathbb{R}) $ 是 $ \mathbb{R} $-向量空间。与 $ \mathcal{V}^{q}(\mathbb{R}) $ 的情形一样，我们看出 $ \{X^{m}dx \;;\; m \in \mathbb{N}\} \subset \Omega_{(q)}(\mathbb{R}) $ 是 $ \mathbb{R} $ 上的线性无关集。因此 $ \Omega_{(q)}(\mathbb{R}) $ 也是无穷维 $ \mathbb{R} $-向量空间。\nobreak\hspace{0.35em}■
\item

$ C^{q}(X) $-模 $ \mathcal{V}^{q}(X) $ 与 $ \Omega_{(q)}(X) $ 同构。我们如下定义一个模同构，即\textbf{典范同构}：

\[
\Theta : \mathcal{V}^{q}(X) \to \Omega_{(q)}(X) , \quad \sum_{j=1}^{n}a_{j}e_{j} \mapsto \sum_{j=1}^{n}a_{j}dx^{j} .
\]

\proofhead
这由 (c) 和 (d) 得出。\nobreak\hspace{0.35em}■
\item

对 $ f \in C^{q+1}(X) $，有 $ \Theta^{-1}df = \mathrm{grad}\,f $，即，图\footnote{与记号 $ df $ 类似，我们用 $ \mathrm{grad}\,f $ 或 $ \nabla f $ 表示映射 $ p \mapsto \nabla_{p}f $（见第 VII.10 节）。在每一种具体情形中，从上下文可以清楚地看出 $ \mathrm{grad}\,f(p) $ 是指这个函数在点 $ p $ 处的值，还是指 $ \nabla_{p}f $ 的切部分。}

\begin{equation*}
\begin{array}{ccc} & \xrightarrow{\hspace{1.2em} d \hspace{1.2em}} & \Omega_{(q)}(X) \\[0.3ex] C^{q+1}(X) & & \cong\uparrow\ \Theta \\[0.3ex] & \xrightarrow{\hspace{1em}\mathrm{grad}\hspace{1em}} & \mathcal{V}^{q}(X) \end{array}
\end{equation*}

是交换的。

\proofhead
由 (e)，有
\[
\Theta^{-1}df = \Theta^{-1}\left(\sum_{j=1}^{n}\partial_{j}f\,dx^{j}\right) = \sum_{j=1}^{n}\partial_{j}f\,e_{j} = \mathrm{grad}\,f \quad \text{对 } f \in C^{q+1}(X) .\quad\blacksquare
\]
\end{items}

\end{remark}

\needspace{3\baselineskip}\subsection*{恰当形式与梯度场}

我们已经看到，每个函数 $ f \in C^{q+1}(X) $ 都诱导出普法夫形式 $ df $。现在我们要问，是否每个普法夫形式都是这样诱导出来的。

设 $ \alpha \in \Omega_{(q)}(X) $，下面我们使用注 3.2(b) 的等同规则。若存在 $ f \in C^{q+1}(X) $ 使得 $ df = \alpha $，则称 $ \alpha $ 是\textbf{恰当的}，$ f $ 是 $ \alpha $ 的一个\textbf{原函数}。

\begin{remark}

\begin{items}
\item

设 $ X $ 是区域，$ \alpha \in \Omega_{(0)}(X) $ 是恰当的。若 $ f $ 和 $ g $ 是 $ \alpha $ 的原函数，则 $ f - g $ 是常数。

\proofhead
首先，$ \alpha = df = dg $ 蕴含 $ d(f - g) = 0 $。然后应用注 VII.3.11(c)。\nobreak\hspace{0.35em}■
\item

设 $ \alpha = \sum_{j=1}^{n}a_{j}dx^{j} \in \Omega_{(1)}(X) $ 是恰当的。则它满足\textbf{可积性条件}

\[
\partial_{k}a_{j} = \partial_{j}a_{k} \quad \text{对 } 1 \le j,k \le n .
\]

\proofhead
设 $ f \in C^{2}(X) $ 满足 $ df = \alpha $。由注 3.3(d) 和 (e)，有 $ a_{j} = \partial_{j}f $。于是由施瓦茨定理（推论 VII.5.5(ii)）可得
\[
\partial_{k}a_{j} = \partial_{k}\partial_{j}f = \partial_{j}\partial_{k}f = \partial_{j}a_{k} \quad \text{对 } 1 \le j,k \le n .\quad\blacksquare
\]
\item

若 $ X \subset \mathbb{R} $，则每个形式 $ \alpha \in \Omega_{(q)}(X) $ 都是恰当的。

\proofhead
因为 $ X $ 是开区间的无交并（见习题 III.4.6），故只需考虑 $ a < b $ 时 $ X = (a,b) $ 的情形。由注 3.3(c)，对每个 $ \alpha \in \Omega_{(q)}(X) $ 存在 $ a \in C^{q}(X) $ 使得 $ \alpha = a\,dx $。设 $ p_{0} \in (a,b) $，且
\[
f(p) := \int_{p_{0}}^{p}a(x)\,dx \quad \text{对 } p \in (a,b) .
\]

则 $ f $ 属于 $ C^{q+1}(X) $，且 $ f' = a $。因此由注 3.3(e) 得 $ df = a\,dx = \alpha $。\nobreak\hspace{0.35em}■

\end{items}

借助典范同构 $ \Theta : \mathcal{V}^{q}(X) \to \Omega_{(q)}(X) $，普法夫形式的性质可以转移到向量场上。于是我们引入如下记号：若存在 $ f \in C^{q+1}(X) $ 使得 $ v = \nabla f $，则称向量场 $ v \in \mathcal{V}^{q}(X) $ 为\textbf{梯度场}。在此语境下，称 $ f $ 为 $ v $ 的一个\textbf{势}。

注 3.3 和 3.4 直接蕴含下面关于梯度场的事实。

\end{remark}

\begin{remark}

\begin{items}
\item

若 $ X $ 是区域，则其上任何梯度场的两个势至多相差一个常数。

\item

若 $ v = (v_{1},\ldots,v_{n}) \in \mathcal{V}^{1}(X) $ 是梯度场，则可积性条件 $ \partial_{k}v_{j} = \partial_{j}v_{k} $ 对 $ 1 \le j,k \le n $ 成立。

\item

普法夫形式 $ \alpha \in \Omega_{(q)}(X) $ 是恰当的当且仅当 $ \Theta^{-1}\alpha \in \mathcal{V}^{q}(X) $ 是梯度场。\nobreak\hspace{0.35em}■

\end{items}

\end{remark}

\begin{example}

\begin{items}
\item

（中心场）设 $ X := \mathbb{R}^{n}\setminus\{0\} $，且普法夫形式

\[
\alpha = \sum_{j=1}^{n}a_{j}dx^{j} \in \Omega_{(q)}(X)
\]

的分量有表示 $ a_{j}(x) := x^{j}\varphi(|x|) $，其中 $ x = (x^{1},\ldots,x^{n}) $，$ 1 \le j \le n $，且 $ \varphi \in C^{q}((0,\infty),\mathbb{R}) $。则 $ \alpha $ 是恰当的。原函数 $ f $ 由 $ f(x) := \Phi(|x|) $ 给出，其中

\[
\Phi(r) := \int_{r_{0}}^{r}t\varphi(t)\,dt \quad \text{对 } r > 0 ,
\]

其中 $ r_{0} $ 是严格正数。因此向量场 $ v := \sum_{j=1}^{n}a_{j}e_{j} = \Theta^{-1}\alpha $ 是梯度场，其中对每个 $ x \in X $，向量 $ v(x) \in T_{x}(X) $ 都落在过 $ x $ 和 $ 0 $ 的直线上。\textbf{等势面}，即 $ f $ 的水平集，是 $ r > 0 $ 时的球面 $ rS^{n-1} $；这些球面与梯度场垂直（见例 VII.3.6(c)）。

\end{items}

\bookfigure[0.35]{figures/3c35044a8beed72f.png}

\proofhead
显然 $ \Phi $ 属于 $ C^{q+1}((0,\infty),\mathbb{R}) $，且 $ \Phi'(r) = r\varphi(r) $。于是由链式法则可得
\[
\partial_{j}f(x) = \Phi'(|x|)x^{j}/|x| = x^{j}\varphi(|x|) = a_{j}(x) \quad \text{对 } j = 1,\ldots,n \text{ 且 } x \in X .
\]

其余的断言是显然的。\nobreak\hspace{0.35em}■

\begin{items}
\item[（b）]

中心场 $ x \mapsto cx/|x|^{n} $，其中 $ c \in \mathbb{R}^{\times} $，$ n \ge 2 $，是梯度场。势 $ U $ 由下式给出：

\begin{equation*}
U(x) := \left\{\begin{array}{ll}c\log|x| , & n = 2 , \\ c|x|^{2-n}/(2-n) , & n > 2 ,\end{array}\right.
\end{equation*}

对 $ x \ne 0 $。这个势在物理学中起着重要作用。\footnote{“势”这一名称来源于物理学。}视具体的物理情境而定，它可以是\textbf{牛顿势}，也可以是\textbf{库仑势}。

\proofhead
这是 (a) 的特殊情形。\nobreak\hspace{0.35em}■
\end{items}

\end{example}

\needspace{3\baselineskip}\subsection*{庞加莱引理}

恰当 1-形式是一类特别重要的微分形式。我们已经知道，每个恰当 1-形式都满足可积性条件。另一方面，存在一些 1-形式，它们不恰当却仍满足可积性条件。因此，寻找这样的条件是有用的：除了已经必要的可积性条件之外，这些条件还能保证原函数的存在。我们现在将看到——也许令人惊讶——这种存在性依赖于其所在区域的拓扑性质。

若连续可微的普法夫形式满足可积性条件，则称它是\textbf{闭的}。若 $ \mathbb{R}^{n} $ 的子集 $ M $（关于 $ x_{0} \in M $）满足：存在 $ x_{0} \in M $，使得对每个 $ x \in M $，从 $ x_{0} $ 到 $ x $ 的直线道路 $ [[x_{0},x]] $ 都落在 $ M $ 中，则称 $ M $ 是\textbf{星形的}。

\bookfigure[0.92]{figures/205da2d7ba804ea7.png}

\begin{remark}

\begin{items}
\item

每个恰当的、连续可微的 1-形式都是闭的。

\item

若 $ M $ 是星形的，则 $ M $ 连通。

\proofhead
这由命题 III.4.8 得出。\nobreak\hspace{0.35em}■
\item

每个凸集关于其任何一点都是星形的。\nobreak\hspace{0.35em}■

\end{items}

设 $ X $ 关于 $ 0 $ 星形，且 $ f \in C^{1}(X) $ 是 $ \alpha = \sum_{j=1}^{n}a_{j}dx^{j} $ 的原函数。则 $ a_{j} = \partial_{j}f $，且注 3.3(c) 给出

\[
\langle\alpha(tx),e_{j}\rangle = a_{j}(tx) = \partial_{j}f(tx) = (\nabla f(tx)|e_{j}) \quad \text{对 } j = 1,\ldots,n .
\]

因此对 $ x \in X $ 有 $ \langle\alpha(tx),x\rangle = (\nabla f(tx)|x) $，且积分形式的中值定理（定理 VII.3.10）给出表示

\begin{equation}
f(x) - f(0) = \int_{0}^{1}(\nabla f(tx)|x)\,dt = \int_{0}^{1}\langle\alpha(tx),x\rangle\,dt ,\tag{3.3}
\end{equation}

它是下一个定理的关键。

\end{remark}

\begin{theorem}

（庞加莱引理）设 $ X $ 是星形的，$ q \ge 1 $。若 $ \alpha \in \Omega_{(q)}(X) $ 是闭的，则 $ \alpha $ 是恰当的。

\proofhead
设 $ X $ 关于 $ 0 $ 星形，$ \alpha = \sum_{j=1}^{n}a_{j}dx^{j} $。因为对 $ x \in X $，线段 $ [[0,x]] $ 落在 $ X $ 中，所以
\begin{equation}
f(x) := \int_{0}^{1}\langle\alpha(tx),x\rangle\,dt = \sum_{k=1}^{n}\int_{0}^{1}a_{k}(tx)x^{k}\,dt \quad \text{对 } x \in X\tag{3.4}
\end{equation}

定义良好。此外

\[
[(t,x) \mapsto \langle\alpha(tx),x\rangle] \in C^{q}([0,1] \times X) ,
\]

且由可积性条件 $ \partial_{j}a_{k} = \partial_{k}a_{j} $，我们得到

\[
\frac{\partial}{\partial x^{j}}\langle\alpha(tx),x\rangle = a_{j}(tx) + \sum_{k=1}^{n}tx^{k}\partial_{k}a_{j}(tx) .
\]

于是命题 VII.6.7 蕴含 $ f $ 属于 $ C^{q}(X) $，且

\[
\partial_{j}f(x) = \int_{0}^{1}a_{j}(tx)\,dt + \int_{0}^{1}t\left(\sum_{k=1}^{n}x^{k}\partial_{k}a_{j}(tx)\right)dt .
\]

对第二个积分分部积分，取 $ u(t) := t $ 和 $ v_{j}(t) := a_{j}(tx) $，得

\[
\partial_{j}f(x) = \int_{0}^{1}a_{j}(tx)\,dt + ta_{j}(tx)|_{0}^{1} - \int_{0}^{1}a_{j}(tx)\,dt = a_{j}(x) \quad \text{对 } x \in X .
\]

最后，因为对 $ j = 1,\ldots,n $，$ \partial_{j}f = a_{j} $ 属于 $ C^{q}(X) $，所以我们得到 $ f \in C^{q+1}(X) $ 和 $ df = \alpha $（见注 3.3(e)）。

$ X $ 关于点 $ x_{0} \in X $ 星形的情形，显然可以通过平移 $ x \mapsto x - x_{0} $ 化归为上面的情形。\nobreak\hspace{0.35em}■

\end{theorem}

\begin{remark}

\begin{items}
\item

例 3.6(a) 的中心场表明，为使势（或原函数）存在，区域不必是星形的。

\item

定理 3.8 的证明是构造性的，也就是说，公式 (3.4) 给出了一个原函数。\nobreak\hspace{0.35em}■

\end{items}

到目前为止，在 $ \Omega_{(q)}(X) $ 中我们只用由 1-形式 $ dx^{1},\ldots,dx^{n} $ 组成的典范模基。如果我们在 $ \mathcal{V}^{q}(X) $ 中使用典范基 $ (e_{1},\ldots,e_{n}) $（见注 3.3(c)），则这种选择是自然的。换言之，如果我们把 $ X $ 理解为 $ \mathbb{R}^{n} $ 的子流形，则模基 $ (dx^{1},\ldots,dx^{n}) $ 适应于平凡坐标卡 $ \varphi = \mathrm{id}_{X} $。如果我们用另一个坐标卡 $ \psi $ 来描述 $ X $，则可以预期，$ \alpha \in \Omega_{(q)}(X) $ 在适应于 $ \psi $ 的模基中能更好地表示。\footnote{在第 III 卷研究流形上的微分形式时，我们将对这一问题有更好的理解。}为了从一种表示过渡到另一种表示，我们必须研究普法夫形式在坐标卡变换下如何变换。

\end{remark}

\needspace{3\baselineskip}\subsection*{对偶算子}

作为矩阵转置的推广，我们首先引入线性算子的“转置”或“对偶”的概念。

设 $ E $ 和 $ F $ 是巴拿赫空间。则对每个 $ A \in \mathcal{L}(E,F) $，我们如下定义一个\textbf{对偶}算子或\textbf{转置}算子：

\[
A^{\top} : F' \to E' , \quad f' \mapsto f' \circ A .
\]

\begin{remark}

\begin{items}
\item

$ A $ 的对偶算子是线性的且连续的，即 $ A^{\top} \in \mathcal{L}(F',E') $。此外\footnote{在泛函分析中已证明，$ A^{\top} $ 由 (3.5) 刻画，且 (3.6) 中等号成立。}

\begin{equation}
\langle A^{\top}f',e\rangle = \langle f',Ae\rangle \quad \text{对 } e \in E \text{ 且 } f' \in F' ,\tag{3.5}
\end{equation}

且

\begin{equation}
\|A^{\top}\|_{\mathcal{L}(F',E')} \le \|A\|_{\mathcal{L}(E,F)} \quad \text{对 } A \in \mathcal{L}(E,F) .\tag{3.6}
\end{equation}

\proofhead
由 $ A \in \mathcal{L}(E,F) $ 可知，对 $ f' \in F' $，$ A^{\top}f' $ 属于 $ E' $。显然 $ A^{\top} : F' \to E' $ 是线性的，且 (3.5) 成立。后者蕴含
\[
|\langle A^{\top}f',e\rangle| = |\langle f',Ae\rangle| \le \|f'\|\,\|Ae\| \le \|f'\|\,\|A\|\,\|e\| \quad \text{对 } e \in E \text{ 且 } f' \in F' .
\]

由此我们读出（见注 VII.2.13(a)）不等式

\[
\|A^{\top}f'\| \le \|A\|\,\|f'\| \quad \text{对 } f' \in F' .
\]

因此由算子范数的定义（见第 VI.2 节）我们得到 (3.6)。现在定理 VI.2.5 保证 $ A^{\top} $ 连续。\nobreak\hspace{0.35em}■

\item

若 $ \mathcal{E} = (e_{1},\ldots,e_{n}) $ 和 $ \mathcal{F} = (f_{1},\ldots,f_{m}) $ 分别是有限维巴拿赫空间 $ E $ 和 $ F $ 的基，则 $ [A^{\top}]_{\mathcal{E}',\mathcal{F}'} = \left([A]_{\mathcal{E},\mathcal{F}}\right)^{\top} $，其中 $ M^{\top} \in \mathbb{K}^{n \times m} $ 表示与 $ M \in \mathbb{K}^{m \times n} $ 对偶的矩阵，$ \mathcal{E}' $ 和 $ \mathcal{F}' $ 是与 $ \mathcal{E} $ 和 $ \mathcal{F} $ 对偶的基。

\proofhead
这容易由表示矩阵的定义（见第 VII.1 节）以及把 $ f' $ 在与 $ \mathcal{F} $ 对偶的基中展开得出。\nobreak\hspace{0.35em}■
\item

若 $ G $ 是另一个巴拿赫空间，

\[
(BA)^{\top} = A^{\top}B^{\top} \quad \text{对 } A \in \mathcal{L}(E,F) \text{ 且 } B \in \mathcal{L}(F,G) ,
\]

且 $ (1_{E})^{\top} = 1_{E'} $。

\item

对 $ A \in \mathcal{L}\mathrm{is}(E,F) $，$ A^{\top} $ 属于 $ \mathcal{L}\mathrm{is}(F',E') $，且 $ (A^{\top})^{-1} = (A^{-1})^{\top} $。

\proofhead
这直接由 (c) 得出。\nobreak\hspace{0.35em}■
\end{items}

\end{remark}

\needspace{3\baselineskip}\subsection*{变换法则}

在下文中，设 $ X $ 是 $ \mathbb{R}^{n} $ 中的开集，$ Y $ 是 $ \mathbb{R}^{m} $ 中的开集。在特殊情形 $ n = 1 $ 或 $ m = 1 $ 下，我们也允许 $ X $ 或 $ Y $（视情况而定）是完全区间。又设 $ \varphi \in C^{q}(X,Y) $，其中 $ q \in \mathbb{N}^{\times} \cup \{\infty\} $。

利用 $ \varphi $，我们定义映射

\[
\varphi^{*} : \mathbb{R}^{Y} \to \mathbb{R}^{X} ,
\]

即\textbf{拉回}，其定义为

\[
\varphi^{*}f := f\circ\varphi \quad \text{对 } f \in \mathbb{R}^{Y} .
\]

更确切地说，$ \varphi^{*} $ 是\textbf{函数的拉回}。\footnote{这个定义显然对 $ f \in Z^{Y} $ 以及任意集合 $ X $、$ Y $、$ Z $ 都有意义。}

对普法夫形式，我们通过 $ \varphi $ 定义\textbf{拉回}为

\[
\varphi^{*} : \Omega_{(q-1)}(Y) \to \Omega_{(q-1)}(X) , \quad \boldsymbol{\alpha} \mapsto \varphi^{*}\boldsymbol{\alpha}
\]

其中

\[
\varphi^{*}\boldsymbol{\alpha}(p) := (T_{p}\varphi)^{\top}\boldsymbol{\alpha}(\varphi(p)) .
\]

由其对象出发，总可明显地看出 $ \varphi^{*} $ 指的是函数的拉回还是 1-形式的拉回。拉回把“生活”在 $ Y $ 上的函数（或 1-形式）映到 $ X $ 上的函数（或 1-形式），而原来的 $ \varphi $ 方向相反。

\begin{remark}

\begin{items}
\item

$ f \in C^{q}(Y) $ 与 $ \boldsymbol{\alpha} \in \Omega_{(q-1)}(Y) $ 的拉回由下面交换图的交换性来定义：

\begin{equation*}
\begin{array}{ccc} X & \xrightarrow{\varphi} & Y \\ \varphi^{*}f \searrow & & \swarrow f \\ & \mathbb{R} & \end{array} \qquad \begin{array}{ccc} TX & \xrightarrow{T\varphi} & TY \\ \varphi^{*}\boldsymbol{\alpha} \searrow & & \swarrow \boldsymbol{\alpha} \\ & \mathbb{R} & \end{array}
\end{equation*}

其中对 $ \boldsymbol{v} \in TX $，$ (T\varphi)\boldsymbol{v} \in TY $ 由下式确定：

\[
((T\varphi)\boldsymbol{v})(\varphi(p)) := (T_{p}\varphi)\boldsymbol{v}(p) \quad \text{对 } p \in X .
\]
\item

对 $ \boldsymbol{\alpha} \in \Omega_{(q-1)}(Y) $，

\[
\varphi^{*}\boldsymbol{\alpha}(p) = (T_{p}\varphi)^{\top}\boldsymbol{\alpha}(\varphi(p)) = \left( p,(\partial\varphi(p))^{\top}\boldsymbol{\alpha}(\varphi(p)) \right) \in T_{p}^{*}X \quad \text{对 } p \in X .
\]

\proofhead
由 $ \varphi^{*}\boldsymbol{\alpha} $ 的定义，我们有
\begin{equation*}
\begin{aligned}
\langle\varphi^{*}\boldsymbol{\alpha}(p),\boldsymbol{v}(p)\rangle_{p} &= \langle\boldsymbol{\alpha}(\varphi(p)),(T_{p}\varphi)\boldsymbol{v}(p)\rangle_{\varphi(p)} \\ &= \langle\boldsymbol{\alpha}(\varphi(p)),\partial\varphi(p)\boldsymbol{v}(p)\rangle = \langle(\partial\varphi(p))^{\top}\boldsymbol{\alpha}(\varphi(p)),\boldsymbol{v}(p)\rangle
\end{aligned}
\end{equation*}

对 $ p \in X $ 与 $ \boldsymbol{v} \in \mathcal{V}^{q}(X) $ 成立，于是断言得证。\nobreak\hspace{0.35em}■

\item

映射

\begin{equation}
\varphi^{*} : C^{q}(Y) \to C^{q}(X)\tag{3.7}
\end{equation}

与

\begin{equation}
\varphi^{*} : \Omega_{(q-1)}(Y) \to \Omega_{(q-1)}(X)\tag{3.8}
\end{equation}

是良定义的且为 $ \mathbb{R} $-线性的。应特别指出，在 (3.8) 中所考虑的普法夫形式的“正则性的阶”是 $ q - 1 $ 而不是 $ q $，而在 (3.7) 中 $ q $ 才是正则性的正确阶数。由于定义涉及切映射，正则性被“损失”了一阶。

\proofhead
关于映射 (3.7) 的断言直接由链式法则得出。(3.8) 的性质同样由链式法则与 (b) 推出。\nobreak\hspace{0.35em}■
\end{items}

在下一个定理中，我们汇总拉回的法则。

\end{remark}

\begin{proposition}

设 $ Z $ 是 $ \mathbb{R}^{\ell} $ 中的开集，或（在 $ \ell = 1 $ 的情形）是完全区间。则

\begin{subitems}
\item

对 $ \mathcal{F}(X) := C^{q}(X) $ 或 $ \mathcal{F}(X) := \Omega_{(q-1)}(X) $，有 $ (\operatorname{id}_{X})^{*} = \operatorname{id}_{\mathcal{F}(X)} $；且对 $ \varphi \in C^{q}(X,Y) $ 与 $ \psi \in C^{q}(Y,Z) $，有 $ (\psi\circ\varphi)^{*} = \varphi^{*}\psi^{*} $；

\item

对 $ f,g \in C^{q}(Y) $，有 $ \varphi^{*}(fg) = (\varphi^{*}f)(\varphi^{*}g) $；且对 $ h \in C^{(q-1)}(Y) $ 与 $ \boldsymbol{\alpha} \in \Omega_{(q-1)}(Y) $，有 $ \varphi^{*}(h\boldsymbol{\alpha}) = (\varphi^{*}h)\varphi^{*}\boldsymbol{\alpha} $；

\item

对 $ f \in C^{q}(Y) $，有 $ \varphi^{*}(df) = d(\varphi^{*}f) $，也就是说，下图

\begin{equation*}
\begin{array}{ccc} C^{q}(Y) & \xrightarrow{d} & \Omega_{(q-1)}(Y) \\ \varphi^{*}\downarrow & & \downarrow\varphi^{*} \\ C^{q}(X) & \xrightarrow{d} & \Omega_{(q-1)}(X) \end{array}
\end{equation*}

交换。

\proofhead
(i) 第一个断言是显然的。为证第二个公式，我们需证明
\[
(\psi\circ\varphi)^{*}f = (\varphi^{*}\circ\psi^{*})f \quad \text{对 } f \in C^{q}(Z) ,
\]

以及

\[
(\psi\circ\varphi)^{*}\boldsymbol{\alpha} = (\varphi^{*}\circ\psi^{*})\boldsymbol{\alpha} \quad \text{对 } \boldsymbol{\alpha} \in \Omega_{(q-1)}(Z) .
\]

对 $ f \in C^{q}(Z) $，我们有

\[
(\psi\circ\varphi)^{*}f = f\circ(\psi\circ\varphi) = (f\circ\psi)\circ\varphi = (\psi^{*}f)\circ\varphi = (\varphi^{*}\circ\psi^{*})f .
\]

对 $ \boldsymbol{\alpha} \in \Omega_{(q-1)}(Y) $，由注 VII.10.2(b) 的链式法则与注 3.10(c) 可得，对 $ p \in X $，

\begin{equation*}
\begin{aligned}
(\psi\circ\varphi)^{*}\boldsymbol{\alpha}(p) &= (T_{p}(\psi\circ\varphi))^{\top}\boldsymbol{\alpha}(\psi\circ\varphi(p)) = ((T_{\varphi(p)}\psi)\circ T_{p}\varphi)^{\top}\boldsymbol{\alpha}(\psi(\varphi(p))) \\ &= (T_{p}\varphi)^{\top}(T_{\varphi(p)}\psi)^{\top}\boldsymbol{\alpha}(\psi(\varphi(p))) = (T_{p}\varphi)^{\top}\psi^{*}\boldsymbol{\alpha}(\varphi(p)) \\ &= \varphi^{*}\psi^{*}\boldsymbol{\alpha}(p) .
\end{aligned}
\end{equation*}
\end{subitems}

\begin{subitems}
\item[（ii）]

第一个陈述是显然的。对 $ p \in X $，第二个陈述由下式得出：

\begin{equation*}
\begin{aligned}
\varphi^{*}(h\boldsymbol{\alpha})(p) &= (T_{p}\varphi)^{\top}(h\boldsymbol{\alpha})(\varphi(p)) = (T_{p}\varphi)^{\top}h(\varphi(p))\boldsymbol{\alpha}(\varphi(p)) \\ &= h(\varphi(p))(T_{p}\varphi)^{\top}\boldsymbol{\alpha}(\varphi(p)) = ((\varphi^{*}h)\varphi^{*}\boldsymbol{\alpha})(p) .
\end{aligned}
\end{equation*}
\item[（iii）]

由切映射的链式法则与微分的定义，我们得到

\begin{equation*}
\begin{aligned}
\varphi^{*}(df)(p) &= (T_{p}\varphi)^{\top}df(\varphi(p)) = df(\varphi(p))\circ T_{p}\varphi \\ &= \mathrm{pr}_{2}\circ T_{\varphi(p)}f\circ T_{p}\varphi = \mathrm{pr}_{2}\circ T_{p}(f\circ\varphi) = d(\varphi^{*}f)(p) ,
\end{aligned}
\end{equation*}

因此断言成立。\nobreak\hspace{0.35em}■

\end{subitems}

\end{proposition}

\begin{corollary}

对 $ \varphi \in \operatorname{Diff}^{q}(X,Y) $，映射

\[
\varphi^{*} : C^{q}(Y) \to C^{q}(X)
\]

与

\[
\varphi^{*} : \Omega_{(q-1)}(Y) \to \Omega_{(q-1)}(X)
\]

是双射，且 $ (\varphi^{*})^{-1} = (\varphi^{-1})^{*} $。

\proofhead
这由命题 3.12(i) 推出。\nobreak\hspace{0.35em}■
\end{corollary}

\begin{example}

\begin{items}
\item

分别以 $ x = (x^{1},\ldots,x^{n}) \in X $ 与 $ y = (y^{1},\ldots,y^{m}) \in Y $ 记 $ X $ 与 $ Y $ 中的点。则

\[
\varphi^{*}dy^{j} = d\varphi^{j} = \sum_{k=1}^{n}\partial_{k}\varphi^{j}\,dx^{k} \quad \text{对 } 1 \le j \le m .
\]

\proofhead
由命题 3.12(iii)，我们有
\[
\varphi^{*}dy^{j} = \varphi^{*}d(\mathrm{pr}_{j}) = d(\varphi^{*}\mathrm{pr}_{j}) = d\varphi^{j} = \sum_{k=1}^{n}\partial_{k}\varphi^{j}\,dx^{k} ,
\]

其中最后一个等式由注 3.3(e) 给出。\nobreak\hspace{0.35em}■

\item

对 $ \alpha = \sum_{j=1}^{m}a_{j}\,dy^{j} \in \Omega_{(0)}(Y) $，我们有

\[
\varphi^{*}\alpha = \sum_{j=1}^{m}\sum_{k=1}^{n}(a_{j}\circ\varphi)\partial_{k}\varphi^{j}\,dx^{k} .
\]

\proofhead
由命题 3.12(ii) 与 $ \varphi^{*} $ 的线性性可得
\[
\varphi^{*}\left( \sum_{j=1}^{m}a_{j}\,dy^{j} \right) = \sum_{j=1}^{m}\varphi^{*}(a_{j}\,dy^{j}) = \sum_{j=1}^{m}(\varphi^{*}a_{j})(\varphi^{*}dy^{j}) .
\]

于是断言由 (a) 推出。\nobreak\hspace{0.35em}■

\item

设 $ X $ 是 $ \mathbb{R} $ 中的开集且 $ \alpha \in \Omega_{(0)}(Y) $。则

\[
\varphi^{*}\alpha(t) = \langle\alpha(\varphi(t)),\dot{\varphi}(t)\rangle\,dt \quad \text{对 } t \in X .
\]

\proofhead
这是 (b) 的特殊情形。\nobreak\hspace{0.35em}■
\item

设 $ Y \subset \mathbb{R} $ 是紧致区间，$ f \in C(Y) $，且 $ \alpha := f\,dy \in \Omega_{(0)}(Y) $。再定义 $ X := [a,b] $，并设 $ \varphi \in C^{1}(X,Y) $。则由 (b)，

\[
\varphi^{*}\alpha = \varphi^{*}(f\,dy) = (f\circ\varphi)\,d\varphi = (f\circ\varphi)\varphi'\,dx .
\]

于是，积分的换元法则（定理 VI.5.1）

\[
\int_{a}^{b}(f\circ\varphi)\varphi'\,dx = \int_{\varphi(a)}^{\varphi(b)}f\,dy
\]

可借助普法夫形式 $ \alpha = f\,dy $ 与 $ \varphi^{*}\alpha $ 形式地写成

\begin{equation}
\int_{X}\varphi^{*}\alpha = \int_{\varphi(X)}\alpha .\tag{3.9}
\end{equation}

因为 $ \Omega_{(0)}(X) $ 是一维的 $ C(X) $-模，$ X $ 上的每个连续 1-形式 $ \beta $ 都可唯一地表示成 $ \beta = b\,dx $ 的形式，其中 $ b \in C(X) $。受 (3.9) 的启发，我们如下定义 $ \beta = b\,dx \in \Omega_{(0)}(X) $ 在 $ X $ 上的\textbf{积分}：

\begin{equation}
\int_{X}\beta := \int_{X}b ,\tag{3.10}
\end{equation}

其中右边理解为柯西–黎曼积分。于是，我们把先前在注 VI.5.2 中给出的关于“微分”的形式定义替换成了数学上精确的定义（至少对实值函数的情形如此）。\nobreak\hspace{0.35em}■

\end{items}

\end{example}

\needspace{3\baselineskip}\subsection*{模}

为完整起见，我们在此汇集模论中最相关的事实。在第 III 卷中，我们将通过一些例子进一步研究模 $ \mathcal{V}^{q}(X) $ 与 $ \Omega_{(q)}(X) $。

设 $ R = (R,+,\cdot) $ 是含单位元的交换环。\textbf{环 $ R $ 上的模}，即 \textbf{$ R $-模}，是一个三元组 $ (M,+,\cdot) $，它由一个非空集合 $ M $ 和两种运算构成：一种“内”运算 $ + $，称为\textbf{加法}，和一种“外”运算 $ \cdot $，称为\textbf{乘法}（用 $ R $ 的元素去乘）。它们必须满足

$ (\mathrm{M}_{1}) $ $ (M,+) $ 是阿贝尔群；

$ (\mathrm{M}_{2}) $ 分配性质，即

\[
\lambda\cdot(v+w) = \lambda\cdot v + \lambda\cdot w , \quad (\lambda+\mu)\cdot v = \lambda\cdot v + \mu\cdot v \quad \text{对 } \lambda,\mu \in R \text{ 且 } v,w \in M ;
\]

$ (\mathrm{M}_{3}) $ $ \lambda\cdot(\mu\cdot v) = (\lambda\cdot\mu)\cdot v , \quad 1\cdot v = v \quad \text{对 } \lambda,\mu \in R \text{ 且 } v \in M $。

\begin{remark}

设 $ M = (M,+,\cdot) $ 是 $ R $-模。

\begin{items}
\item

与向量空间的情形一样，我们规定乘法优先于加法结合，并通常用 $ \lambda v $ 记 $ \lambda\cdot v $。

\item

公理 $ (\mathrm{M}_{3}) $ 表明环 $ R $ 从左作用于集合 $ M $（见习题 I.7.6）。

\item

$ R $ 与 $ M $ 中的加法运算都记作 $ + $。类似地，$ R $ 中的环乘法与 $ R $ 在 $ M $ 上的乘法都用同一个 $ \cdot $ 表示。最后，我们用 0 表示 $ (R,+) $ 与 $ (M,+) $ 的零元。经验表明，这种简化的记号不会引起误解。

\item

若 $ R $ 是域，则 $ M $ 是 $ R $ 上的向量空间。\nobreak\hspace{0.35em}■

\end{items}

\end{remark}

\begin{example}

\begin{items}
\item

每个含单位元的交换环 $ R $ 都是 $ R $-模。

\item

设 $ G = (G,+) $ 是阿贝尔群。对 $ (z,g) \in \mathbb{Z}\times G $，令\footnote{见例 I.5.12。}

\begin{equation*}
z\cdot g := \begin{cases} \sum_{k=1}^{z}g , & z > 0 , \\ 0 , & z = 0 , \\ -\sum_{k=1}^{-z}g , & z < 0 . \end{cases}
\end{equation*}

则 $ G = (G,+,\cdot) $ 是 $ \mathbb{Z} $-模。

\item

设 $ G $ 是阿贝尔群，$ R $ 是 $ G $ 的自同态环的交换子环。则 $ G $ 连同运算

\[
R\times G \to G , \quad (T,g) \mapsto Tg ,
\]

是 $ R $-模。

\item

设 $ M $ 是 $ R $-模。$ M $ 的非空子集 $ U $ 是 $ M $ 的\textbf{子模}，如果
$ (\mathrm{UM}_{1}) $ $ U $ 是 $ M $ 的子群，
$ (\mathrm{UM}_{2}) $ $ R\cdot U \subseteq U $。
不难验证，$ M $ 的非空子集 $ U $ 是 $ M $ 的子模当且仅当 $ U $ 关于由 $ M $ 诱导的运算构成模。
若 $ \{U_{\alpha} ; \ \alpha \in A\} $ 是 $ M $ 的子模族，则 $ \bigcap_{\alpha\in A}U_{\alpha} $ 也是 $ M $ 的子模。\nobreak\hspace{0.35em}■

\end{items}

设 $ M $ 与 $ N $ 是 $ R $ 上的模。映射 $ T : M \to N $ 是\textbf{模同态}，如果

\[
T(\lambda v + \mu w) = \lambda T(v) + \mu T(w) \quad \text{对 } \lambda,\mu \in R \text{ 且 } v,w \in M .
\]

双射的模同态称为\textbf{模同构}。此时逆映射 $ T^{-1} : N \to M $ 也是模同态。当存在从 $ M $ 到 $ N $ 的模同构时，我们说 $ M $ 与 $ N $ 同构，并记作 $ M \cong N $。

设 $ M $ 是 $ R $-模且 $ A \subset M $。则

\[
\operatorname{span}(A) := \bigcap\{U ; \ U \text{ 是 } M \text{ 的子模，且 } U \supseteq A\}
\]

称为 $ A $ 在 $ M $ 中的\textbf{张成}。显然 $ \operatorname{span}(A) $ 是 $ M $ 中包含 $ A $ 的最小子模。当 $ \operatorname{span}(A) $ 等于 $ M $ 时，我们称 $ A $ 是 $ M $ 的\textbf{生成系}。

元素 $ v_{0},\ldots,v_{k} \in M $ 称为在 $ R $ 上\textbf{线性无关}的，如果 $ \lambda_{j} \in R $ 与 $ \sum_{j=0}^{k}\lambda_{j}v_{j} = 0 $ 蕴含 $ \lambda_{0} = \cdots = \lambda_{k} = 0 $。$ M $ 的子集 $ B $ 称为在 $ R $ 上\textbf{自由}的，如果 $ B $ 的每个有限子集都线性无关。若对 $ R $ 上的某个子集 $ B $ 有 $ \operatorname{span}(B) = M $，则 $ B $ 构成 $ M $ 的一个\textbf{基}。具有基的 $ R $-模称为\textbf{自由 $ R $-模}。

\end{example}

\begin{remark}

\begin{items}
\item

以下陈述等价：

\begin{subitems}
\item

$ M $ 是以 $ B $ 为基的自由 $ R $-模。

\item

每个 $ v \in M $ 都具有 $ v = \sum_{j=0}^{k}\lambda_{j}v_{j} $ 的形式，其中 $ \lambda_{j} \in R^{\times} $ 与 $ v_{j} \in B $ 是唯一的。

\proofhead
参见代数教材（例如 [SS88, §§ 19 与 22]）。\nobreak\hspace{0.35em}■
\end{subitems}

\item

设 $ M $ 与 $ N $ 是 $ R $ 上的模，且 $ M $ 是以 $ B $ 为基的自由模。再设 $ S $ 与 $ T $ 是从 $ M $ 到 $ N $ 的模同态，且 $ S|B = T|B $。则 $ S = T $，也就是说，自由模上的同态由它对其基的作用唯一确定。

\item

设 $ B $ 与 $ B' $ 是自由 $ R $-模 $ M $ 的两个基。则可以证明，若 $ R $ 不是零环，则 $ \operatorname{Num}(B) = \operatorname{Num}(B') $（例如 [Art93, § XII.2]）。于是 $ M $ 的\textbf{维数}定义为 $ \dim(M) := \operatorname{Num}(B) $。

\item

$ \mathbb{Z}_{2} $ 不是自由 $ \mathbb{Z} $-模。

\proofhead
由习题 I.9.2 知 $ \mathbb{Z}_{2} $ 是环。因此 $ (\mathbb{Z}_{2},+) $ 是阿贝尔群，例 3.16(b) 表明 $ \mathbb{Z}_{2} $ 是 $ \mathbb{Z} $-模。设 $ B $ 是 $ \mathbb{Z}_{2} $ 在 $ \mathbb{Z} $ 上的基。则陪集 $ [1] $ 属于 $ B $。另一方面，$ [0] = 2\cdot[1] $ 表明 $ [1] $ 在 $ \mathbb{Z} $ 上不是线性无关的，即 $ [1] \notin B $。\nobreak\hspace{0.35em}■
\item

设 $ M $ 是 $ R $-模，且对 $ B \subset M $ 有 $ M = \operatorname{span}(B) $。则\footnote{设 $ V $ 是 $ K $ 向量空间且 $ V = \operatorname{span}(B) $。则在线性代数中已证明 $ B $ 包含 $ V $ 的一个基。} $ B $（一般地）不包含 $ M $ 的基。

\proofhead
我们在 $ \mathbb{Z} $-模\footnote{注意 $ \mathbb{Z} $ 是自由 $ \mathbb{Z} $-模。} $ \mathbb{Z} $ 中考虑集合 $ B = \{2,3\} $。由 $ \mathbb{Z} = \mathbb{Z}(3-2) $ 得 $ \operatorname{span}(B) = \mathbb{Z} $。此外，由
\[
3\cdot 2 - 2\cdot 3 = 0 \quad \text{且} \quad z\cdot 2 \ne 5 \ne z\cdot 3 \quad \text{对 } z \in \mathbb{Z}
\]

可知 $ \{2,3\} $、$ \{2\} $、$ \{3\} $ 都不是 $ \mathbb{Z} $ 的基。\nobreak\hspace{0.35em}■

\item

设 $ M $ 是自由 $ R $-模，$ U $ 是 $ M $ 的以 $ B' $ 为基的自由子模。则 $ B' $（一般地）不生成 $ M $ 的基。

\proofhead
$ \mathbb{Z} $ 是 $ \mathbb{Z} $ 上的自由模。偶数 $ 2\mathbb{Z} $ 构成 $ \mathbb{Z} $ 的自由子模。然而，$ \{2\} $ 与 $ \{-2\} $ 都不生成 $ \mathbb{Z} $ 的基。并且 $ \dim(\mathbb{Z}) = \dim(2\mathbb{Z}) = 1 $。\nobreak\hspace{0.35em}■
\item

若 $ R $ 是域，则上述“张成”、“线性无关”、“基”与“维数”等概念与第 I.12 节中的概念一致。\nobreak\hspace{0.35em}■

\end{items}

\end{remark}

\exerciseshead

\begin{exercises}
\item

设 $ E, F $ 与 $ G $ 为巴拿赫空间；又令 $ A \in \mathcal{L}(E, F) $ 与 $ B \in \mathcal{L}(G, E) $。证明

\begin{items}
\item

$ (AB)^{\top} = B^{\top}A^{\top} $；

\item

对 $ A \in \mathcal{L}\mathrm{is}(E, F) $，有 $ A^{\top} \in \mathcal{L}\mathrm{is}(F', E') $ 且 $ [A^{\top}]^{-1} = [A^{-1}]^{\top} $。

\end{items}

\item

这些 1-形式中哪些是闭的？

\begin{items}
\item

$ 2xy^{3}\,dx + 3x^{2}y^{2}\,dy \in \Omega_{(\infty)}(\mathbb{R}^{2}) $。

\item

$ \frac{2(x^{2} - y^{2} - 1)\,dy - 4xy\,dx}{(x^{2} + y^{2} - 1)^{2} + 4y^{2}} \in \Omega_{(\infty)}(X) $，其中 $ X := \mathbb{R}^{2}\setminus\{(-1,0),(1,0)\} $。

\end{items}

\item

设 $ X := \mathbb{R}^{2}\setminus\{(-1,0),(1,0)\} $，

\[
\varphi : X \to \mathbb{R}^{2}\setminus\{(0,0)\} , \quad (x, y) \mapsto (x^{2} + y^{2} - 1, 2y)
\]

并且

\[
\alpha := \frac{u\,dv - v\,du}{u^{2} + v^{2}} \in \Omega(\mathbb{R}^{2}\setminus\{(0,0)\}) .
\]

证明 $ \varphi^{*}\alpha $ 与习题 2(b) 的 1-形式相同。

\item

设 $ X $ 在 $ \mathbb{R}^{n} $ 中为开集且 $ \alpha \in \Omega_{(1)}(X) $。用 $ a \in C^{1}(X, \mathbb{R}^{n}) $ 表示 $ \Theta^{-1}\alpha $ 的余切部分。证明 $ \alpha $ 是闭的当且仅当对每个 $ x \in X $，$ \partial a(x) $ 对称。

\item

设

\[
\alpha := x\,dy - y\,dx \in \Omega_{(\infty)}(\mathbb{R}^{2}) \quad \text{且} \quad \beta := \alpha/(x^{2} + y^{2}) \in \Omega_{(\infty)}(\mathbb{R}^{2}\setminus\{(0,0)\})
\]

并且

\[
\varphi : (0,\infty)\times(0,2\pi) \to \mathbb{R}^{2}\setminus\{(0,0)\} , \quad (r, \theta) \mapsto (r\cos\theta, r\sin\theta) .
\]

计算 $ \varphi^{*}\alpha $ 与 $ \varphi^{*}\beta $。

\item

对 $ \alpha := x\,dy - y\,dx - (x + y)\,dz \in \Omega_{(\infty)}(\mathbb{R}^{3}) $ 与

\[
\varphi : (0,2\pi)\times(0,\pi) \to \mathbb{R}^{3} , \quad (\theta_{1}, \theta_{2}) \mapsto (\cos\theta_{1}\sin\theta_{2}, \sin\theta_{1}\sin\theta_{2}, \cos\theta_{2})
\]

确定 $ \varphi^{*}\alpha $。

\item

设 $ X $ 与 $ Y $ 在 $ \mathbb{R}^{n} $ 中为开集，$ \varphi \in \operatorname{Diff}^{q}(X, Y) $，且对某 $ q \ge 1 $，$ \alpha \in \Omega_{(q-1)}(Y) $。证明

\begin{items}
\item

$ \alpha $ 是恰当的当且仅当 $ \varphi^{*}\alpha $ 是恰当的；

\item

对 $ q \ge 2 $，$ \alpha $ 是闭的当且仅当 $ \varphi^{*}\alpha $ 是闭的。

\end{items}

\item

设 $ X $ 在 $ \mathbb{R}^{2} $ 中为开集且 $ \alpha \in \Omega_{(1)}(X) $。函数 $ h \in C^{1}(X) $ 称为 $ \alpha $ 的\textbf{欧拉乘子}（或\textbf{积分因子}），如果 $ h\alpha $ 是闭的且对 $ (x, y) \in X $ 有 $ h(x, y) \ne 0 $。

\begin{items}
\item

证明若 $ h \in C^{1}(X) $ 满足对 $ (x, y) \in X $ 有 $ h(x, y) \ne 0 $，则 $ h $ 是 $ \alpha = a\,dx + b\,dy $ 的欧拉乘子当且仅当

\begin{equation}
a h_{y} - b h_{x} + (a_{y} - b_{x})h = 0 .\tag{3.11}
\end{equation}
\item

设 $ a, b, c, e > 0 $，并令

\[
\beta := (c - ex)y\,dx + (a - by)x\,dy .
\]

\proofhead
$ \beta $ 在 $ X = (0,\infty)^{2} $ 上具有形如 $ h(x, y) = m(xy) $ 的积分因子。（提示：应用 (3.11) 求出 $ m $ 的微分方程，并猜测其解。）
\end{items}

\item

设 $ J $ 与 $ D $ 为开区间，$ f \in C(J\times D, \mathbb{R}) $ 且 $ \alpha := dx - f\,dt \in \Omega_{(0)}(J\times D) $。进一步设 $ u \in C^{1}(J, D) $，并定义 $ \varphi_{u} : J \to \mathbb{R}^{2} $，$ t \mapsto (t, u(t)) $。证明 $ \varphi^{*}_{u}\alpha = 0 $ 当且仅当 $ u $ 满足 $ \dot x = f(t, x) $。

\item

验证 $ \mathcal{V}^{q}(X) $ 与 $ \Omega_{(q)}(X) $ 是 $ C^{q}(X) $-模，且 $ \mathcal{V}^{q+1}(X) $ 与 $ \Omega_{(q+1)}(X) $ 分别是 $ \mathcal{V}^{q}(X) $ 与 $ \Omega_{(q)}(X) $ 的子模。

\item

设 $ R $ 是带单位元的交换环且 $ a \in R $。证明

\begin{items}
\item

$ aR $ 是 $ R $ 的子模；

\item

$ aR \ne 0 \iff a \ne 0 $；

\item

$ aR \ne R \iff a $ 没有逆元。
$ R $-模 $ M $ 的子模 $ U $ 称为\textbf{非平凡的}，如果 $ U \ne \{0\} $ 且 $ U \ne M $。$ \mathbb{Z}_{2} $ 与 $ \mathbb{Z}_{6} $ 的非平凡子模是什么？

\end{items}

\item

证明例 3.16(d) 的陈述。

\item

设 $ M $ 与 $ N $ 是 $ R $-模且 $ T : M \to N $ 是模同态。验证 $ \ker(T) := \{v \in M \ ; \ Tv = 0\} $ 与 $ \operatorname{im}(T) $ 分别是 $ M $ 与 $ N $ 的子模。

\end{exercises}

\section{线积分}

在第 VI 章，我们建立了在区间上对实变量函数积分的理论。若把这些区间看作特别简单的曲线，我们也许会问，积分这一概念能否推广到任意曲线。本节正是要这样做。当然，我们已经能对道路的函数积分，所以要建立一种只依赖于曲线的新积分方式，就必须保证它与参数化无关。结果会是，我们必须对普法夫形式而非函数积分。

在这些节中，设

• $ X $ 是 $ \mathbb{R}^{n} $ 中的开集；$ I $ 与 $ I_1 $ 是紧致完全区间。

此外，我们总是分别把向量场和普法夫形式与其切部分和余切部分等同起来。

\needspace{3\baselineskip}\subsection*{定义}

我们现在来处理积分的换元法则。设 $ I $ 与 $ J $ 是紧致区间且 $ \varphi \in C^1(I, J) $。再设 $ a \in C(J) $ 且 $ \alpha = a\,dy $ 是 $ J $ 上的连续 1-形式。则根据 3.14(d) 与定义 (3.10)，换元法则具有如下形式

\begin{equation}
\int_{\varphi(I)}\alpha=\int_I\varphi^{*}\alpha=\int_I(a\circ\varphi)\dot{\varphi}\,dt .\tag{4.1}
\end{equation}

因为 $ \Omega_{(0)}(J) $ 是一维 $ C(J) $-模，所以每个 $ \alpha \in \Omega_{(0)}(J) $ 有唯一表示 $ \alpha = a\,dy $ 使得 $ a \in C(J) $。因此 (4.1) 对每个 $ \alpha \in \Omega_{(0)}(J) $ 都有定义。最后一个被积函数也可写成如下形式

\[
\langle\alpha(\varphi(t)),\dot{\varphi}(t)\rangle\,dt=\langle\alpha(\varphi(t)),\dot{\varphi}(t)1\rangle\,dt .
\]

这一观察是定义 1-形式的线积分所需的基本结果。

对 $ \alpha \in \Omega_{(0)}(X) $ 与 $ \gamma \in C^1(I, X) $

\[
\int_{\gamma}\alpha:=\int_I\gamma^{*}\alpha=\int_I\langle\alpha(\gamma(t)),\dot{\gamma}(t)\rangle\,dt
\]

称为 \textbf{$ \alpha $ 沿道路 $ \gamma $ 的积分}。

\begin{remark}

\begin{items}
\item

在基表示 $ \alpha=\sum_{j=1}^{n}a_j\,dx^{j} $ 中，我们有

\[
\int_{\gamma}\alpha=\sum_{j=1}^{n}\int_I(a_j\circ\gamma)\dot{\gamma}^{j}\,dt .
\]

\proofhead
这直接由例 3.14(b) 推出。\nobreak\hspace{0.35em}■
\item

设 $ \varphi \in C^1(I_1, I) $ 且 $ \gamma_1 := \gamma\circ\varphi $。则

\[
\int_{\gamma}\alpha=\int_{\gamma_1}\alpha\quad\text{对 }\alpha\in\Omega_{(0)}(X) .
\]

\proofhead
命题 3.12(i) 与换元法则 (3.9) 给出
\[
\int_{\gamma_1}\alpha=\int_{I_1}(\gamma\circ\varphi)^{*}\alpha=\int_{I_1}\varphi^{*}(\gamma^{*}\alpha)=\int_I\gamma^{*}\alpha=\int_{\gamma}\alpha\quad\text{对 }\alpha\in\Omega_{(0)}(X) .\quad\blacksquare
\]

这一结果表明 $ \alpha $ 沿 $ \gamma $ 的积分在参数化改变下不变。于是，对 $ X $ 中每条紧致 $ C^1 $ 曲线 $ \Gamma $ 及每个 $ \alpha \in \Omega_{(0)}(X) $，\textbf{$ \alpha $ 沿 $ \Gamma $ 的线积分}

\[
\int_{\Gamma}\alpha:=\int_{\gamma}\alpha
\]

是良定义的，其中 $ \gamma $ 是 $ \Gamma $ 的任意 $ C^1 $ 参数化。

\end{items}

\end{remark}

\begin{example}

\begin{items}
\item

设 $ \Gamma $ 是由 $ \gamma : [0, 2\pi]\to\mathbb{R}^{2} $，$ t\mapsto R(\cos t, \sin t) $ 参数化的圆。则

\[
\int_{\Gamma}X\,dy-Y\,dx=2\pi R^{2} .
\]

\proofhead
由
\[
\gamma^{*}dx=d\gamma^{1}=-(R\sin)\,dt\quad\text{且}\quad\gamma^{*}dy=d\gamma^{2}=(R\cos)\,dt
\]

可得 $ \gamma^{*}(X\,dy-Y\,dx)=R^{2}\,dt $，断言得证。\nobreak\hspace{0.35em}■

\item

设 $ \alpha \in \Omega_{(0)}(X) $ 是恰当的。若 $ f \in C^1(X) $ 是一个原函数，\footnote{这一陈述是推论 VI.4.14 的推广。}

\[
\int_{\Gamma}\alpha=\int_{\Gamma}df=f(E_{\Gamma})-f(A_{\Gamma}) ,
\]

其中 $ A_{\Gamma} $ 与 $ E_{\Gamma} $ 分别是 $ \Gamma $ 的起点与终点。

\proofhead
我们取定 $ \Gamma $ 的一个 $ C^1 $ 参数化 $ \gamma : [a, b]\to X $。则由命题 3.12(iii) 可得
\[
\gamma^{*}df=d(\gamma^{*}f)=d(f\circ\gamma)=(f\circ\gamma)^{\cdot}\,dt .
\]

因此我们得到

\[
\int_{\Gamma}\alpha=\int_{\Gamma}df=\int_a^b\gamma^{*}df=\int_a^b(f\circ\gamma)^{\cdot}\,dt=f\circ\gamma\big|_a^b=f(E_{\Gamma})-f(A_{\Gamma})
\]

因为 $ E_{\Gamma}=\gamma(b) $ 且 $ A_{\Gamma}=\gamma(a) $。\nobreak\hspace{0.35em}■

\item

恰当普法夫形式沿 $ C^1 $ 曲线 $ \Gamma $ 的积分只依赖于起点与终点，而不依赖于它们之间的道路。若 $ \Gamma $ 是闭的，则该积分等于 0。

\proofhead
这直接由 (b) 推出。\nobreak\hspace{0.35em}■
\item

存在不是恰当的闭 1-形式。

\proofhead
设 $ X:=\mathbb{R}^{2}\setminus\{(0, 0)\} $ 且 $ \alpha \in \Omega_{(\infty)}(X) $ 满足
\[
\alpha(x, y):=\frac{x\,dy-y\,dx}{x^{2}+y^{2}}\quad\text{对 }(x, y)\in X .
\]

容易验证 $ \alpha $ 是闭的。圆 $ \Gamma $ 的参数化 $ \gamma : [0, 2\pi]\to X $，$ t\mapsto(\cos t, \sin t) $ 给出 $ \gamma^{*}\alpha=dt $，因此

\[
\int_{\Gamma}\alpha=\int_0^{2\pi}dt=2\pi\ne 0 .
\]

于是 (c) 表明 $ \alpha $ 不是恰当的。\nobreak\hspace{0.35em}■

\item

设 $ x_0 \in X $ 且 $ \gamma_{x_0} : I\to X $，$ t\mapsto x_0 $。这意味着 $ \gamma_{x_0} $ 参数化\textbf{点曲线} $ \Gamma $，其像恰为 $ \{x_0\} $。则

\[
\int_{\gamma_{x_0}}\alpha=0\quad\text{对 }\alpha\in\Omega_{(0)}(X) .
\]

\proofhead
这是显然的，因为 $ \dot{\gamma}_{x_0}=0 $。\nobreak\hspace{0.35em}■
\end{items}

\end{example}

\needspace{3\baselineskip}\subsection*{基本性质}

设 $ I = [a, b] $ 且 $ \gamma \in C(I, X) $。则

\[
\gamma^{-}:I\to X ,\quad t\mapsto\gamma(a+b-t)
\]

称为 $ \gamma $ 的\textbf{道路逆}，而 $ -\Gamma:=[\gamma^{-}] $ 是 $ \Gamma:=[\gamma] $ 的\textbf{曲线逆}。（注意 $ \Gamma $ 与 $ -\Gamma $ 有相同的像但定向相反）。

\bookfigure[0.95]{figures/e281c5911fc0b710.png}

现在设 $ q \in \mathbb{N}^{\times}\cup\{\infty\} $。再设 $ \gamma \in C(I, X) $ 且 $ (t_0,\ldots,t_m) $ 是 $ I $ 的一个划分。若\footnote{见习题 1.5 及第 VI.7 节中分段连续可微函数的定义。这里我们还假设 $ \gamma $ 是连续的。}

\[
\gamma_j:=\gamma|_{[t_{j-1},t_j]}\in C^{q}\left([t_{j-1},t_j], X\right)\quad\text{对 }1\le j\le m ,
\]

则称 $ \gamma $ 是 $ X $ 中的\textbf{分段 $ C^{q} $ 道路}，或称其为 $ C^{q} $ 道路 $ \gamma_j $ 的\textbf{和}。由 $ \gamma $ 参数化的曲线 $ \Gamma=[\gamma] $ 称为 $ X $ 中的\textbf{分段 $ C^{q} $ 曲线}，我们记 $ \Gamma:=\Gamma_1+\cdots+\Gamma_m $，其中 $ \Gamma_j:=[\gamma_j] $。给定 $ X $ 中的分段 $ C^1 $ 曲线 $ \Gamma=\Gamma_1+\cdots+\Gamma_m $ 与 $ \alpha \in \Omega_{(0)}(X) $，我们定义 \textbf{$ \alpha $ 沿 $ \Gamma $ 的线积分}为

\[
\int_{\Gamma}\alpha:=\sum_{j=1}^{m}\int_{\Gamma_j}\alpha .
\]
\bookfigure[0.36]{figures/5029755606ddd965.png}

最后，我们也可以把分段 $ C^{q} $ 曲线连接起来。于是设 $ \Gamma:=\sum_{j=1}^{m}\Gamma_j $ 与 $ \tilde{\Gamma}:=\sum_{j=1}^{\tilde{m}}\tilde{\Gamma}_j $ 是满足 $ E_{\Gamma}=A_{\tilde{\Gamma}} $ 的分段 $ C^{q} $ 曲线。则 $ \Sigma=\Sigma_1+\cdots+\Sigma_{m+\tilde{m}} $，其中

\begin{equation*}
\Sigma_j:=\begin{cases}\Gamma_j , & 1\le j\le m ,\\ \tilde{\Gamma}_{j-m} , & m+1\le j\le m+\tilde{m} ,\end{cases}
\end{equation*}

是一条分段 $ C^{q} $ 曲线，即 $ \Gamma $ 与 $ \tilde{\Gamma} $ 的\textbf{和}\footnote{注意，对两条分段 $ C^{q} $ 曲线 $ \Gamma $ 与 $ \tilde{\Gamma} $，和 $ \Gamma+\tilde{\Gamma} $ 仅在 $ \Gamma $ 的终点是 $ \tilde{\Gamma} $ 的起点时才有定义。}。此时我们记 $ \Gamma+\tilde{\Gamma}:=\Sigma $ 并令

\[
\int_{\Gamma+\tilde{\Gamma}}\alpha:=\int_{\Sigma}\alpha\quad\text{对 }\alpha\in\Omega_{(0)}(X) .
\]

显然，这些记号与此前关于分段 $ C^{q} $ 曲线的记号 $ \Gamma=\Gamma_1+\cdots+\Gamma_m $ 以及线积分的记号 $ \int_{\Gamma}\alpha=\int_{\Gamma_1+\cdots+\Gamma_m}\alpha $ 是一致的。

下一个定理列举了线积分的一些基本性质。

\qyxffnnumdef{1}
\begin{proposition}

设 $ \Gamma $、$ \Gamma_1 $ 与 $ \Gamma_2 $ 是分段 $ C^1 $ 曲线，且 $ \alpha, \alpha_1, \alpha_2 \in \Omega_{(0)}(X) $。则

\begin{subitems}
\item

$ \int_{\Gamma}(\lambda_1\alpha_1+\lambda_2\alpha_2)=\lambda_1\int_{\Gamma}\alpha_1+\lambda_2\int_{\Gamma}\alpha_2 $，$ \lambda_1, \lambda_2 \in \mathbb{R} $，即 $ \int_{\Gamma} : \Omega_{(0)}(X)\to\mathbb{R} $ 是向量空间同态；

\item

$ \int_{-\Gamma}\alpha=-\int_{\Gamma}\alpha $，即线积分是定向的；

\item

若 $ \Gamma_1+\Gamma_2 $ 有定义，则我们有

\[
\int_{\Gamma_1+\Gamma_2}\alpha=\int_{\Gamma_1}\alpha+\int_{\Gamma_2}\alpha ,
\]

即线积分关于其所积分的曲线是可加的；

\item

对 $ \alpha=\sum_{j=1}^{n}a_j\,dx^{j} $ 与 $ a:=\Theta^{-1}\alpha=\sum_{j=1}^{n}a_j e_j=(a_1,\ldots,a_n)\in\mathcal{V}^{0}(X) $，我们有

\[
\left|\int_{\Gamma}\alpha\right|\le\max_{x\in\Gamma}|a(x)|\,L(\Gamma) .
\]

\proofhead
断言 (i) 与 (iii) 是显然的。
\end{subitems}

(ii) 由 (iii)，只需考虑 $ \Gamma $ 是 $ C^{q} $ 曲线的情形。

设 $ \gamma \in C^1([a, b], X) $ 参数化 $ \Gamma $。则我们把 $ \gamma^{-} $ 写成 $ \gamma^{-}=\gamma\circ\varphi $，其中对 $ t \in [a, b] $ 有 $ \varphi(t):=a+b-t $。因为 $ \varphi(a)=b $ 且 $ \varphi(b)=a $，由命题 3.12(i) 与 (3.9) 可得

\[
\int_{-\Gamma}\alpha=\int_a^b(\gamma^{-})^{*}\alpha=\int_a^b(\gamma\circ\varphi)^{*}\alpha=\int_a^b\varphi^{*}(\gamma^{*}\alpha)=\int_b^a\gamma^{*}\alpha=-\int_a^b\gamma^{*}\alpha
\]
\[
=-\int_{\Gamma}\alpha .
\]

(iv) 分段 $ C^1 $ 曲线 $ \Gamma=\Gamma_1+\cdots+\Gamma_m $ 的长度显然等于其各段\qyxffnmark{1} $ \Gamma_j $ 的长度之和。因此只需对一条 $ C^1 $ 曲线 $ \Gamma:=[\gamma] $ 证明该断言。借助柯西-施瓦茨不等式与命题 VI.4.3，可得

\[
\left|\int_{\Gamma}\alpha\right|=\left|\int_I\gamma^{*}\alpha\right|=\left|\int_I\langle\alpha(\gamma(t)),\dot{\gamma}(t)\rangle\,dt\right|=\left|\int_I(a(\gamma(t))\,|\dot{\gamma}(t)|)\,dt\right|
\]
\[
\le\int_I|a(\gamma(t))|\,|\dot{\gamma}(t)|\,dt\le\max_{x\in\Gamma}|a(x)|\int_I|\dot{\gamma}(t)|\,dt
\]
\[
=\max_{x\in\Gamma}|a(x)|\,L(\Gamma) ,
\]

其中我们把 $ (\mathbb{R}^{n})' $ 与 $ \mathbb{R}^{n} $ 等同，且在最后一步应用了定理 1.3。\nobreak\hspace{0.35em}■

\end{proposition}
\qyxffntextof{1}{见习题 1.5。}

\needspace{3\baselineskip}\subsection*{线积分基本定理}

\begin{theorem}

设 $ X\subset\mathbb{R}^{n} $ 是一个区域且 $ \alpha \in \Omega_{(0)}(X) $。则下列断言等价：

\begin{subitems}
\item

$ \alpha $ 是恰当的；

\item

对 $ X $ 中每条闭的分段 $ C^1 $ 曲线都有 $ \int_{\Gamma}\alpha=0 $。

\proofhead
“(i)$ \Rightarrow $(ii)”这一情形由例 4.2(b) 与命题 4.3(iii) 推出。
“(ii)$ \Rightarrow $(i)” 设 $ x_0 \in X $。根据定理 III.4.10，对每个 $ x \in X $，存在 $ X $ 中一条从 $ x_0 $ 通向 $ x $ 的连续、分段直线道路。于是对每个 $ x \in X $，存在 $ X $ 中以 $ x_0 $ 为起点、以 $ x $ 为终点的分段 $ C^1 $ 曲线 $ \Gamma_x $。我们令
\[
f : X\to\mathbb{R} ,\quad x\mapsto\int_{\Gamma_x}\alpha .
\]

为验证 $ f $ 是良定义的，即不依赖于特殊的曲线 $ \Gamma_x $，我们取 $ X $ 中另一条有相同端点的分段 $ C^1 $ 曲线 $ \tilde{\Gamma}_x $。则 $ \Sigma:=\Gamma_x+(-\tilde{\Gamma}_x) $ 是 $ X $ 中的一条闭的分段 $ C^1 $ 曲线。由假设我们有 $ \int_{\Sigma}\alpha=0 $，并且由命题 4.3 我们推出 $ 0=\int_{\Gamma_x}\alpha-\int_{\tilde{\Gamma}_x}\alpha $。因此 $ f $ 是良定义的。

\end{subitems}

\bookfigure[0.44]{figures/456e3dec243c7944.png}

现在设 $ h \in \mathbb{R}^{+} $ 满足 $ \overline{\mathbb{B}}(x, h)\subset X $，且 $ \Pi_j:=[\pi_j] $，其中

\[
\pi_j : [0, 1]\to X ,\quad t\mapsto x+t\,h\,e_j\quad\text{对 }j=1,\ldots,n .
\]

则 $ \Gamma_x+\Pi_j $ 与 $ \Pi_j=(-\Gamma_x)+(\Gamma_x+\Pi_j) $ 是 $ X $ 中的曲线。因为 $ \Gamma_x $ 与 $ \Gamma_x+\Pi_j $ 有相同的起点 $ x_0 $，

\[
f(x+he_j)-f(x)=\int_{\Gamma_x+\Pi_j}\alpha-\int_{\Gamma_x}\alpha=\int_{\Pi_j}\alpha .
\]

令 $ a_j:=\langle\alpha, e_j\rangle $，我们得到

\[
\int_{\Pi_j}\alpha=\int_0^1\pi_j^{*}\alpha=\int_0^1\langle\alpha(x+the_j), he_j\rangle\,dt=h\int_0^1 a_j(x+the_j)\,dt .
\]

因此

\[
f(x+he_j)-f(x)-a_j(x)h=h\int_0^1\left[a_j(x+the_j)-a_j(x)\right]dt=o(h)
\]

当 $ h\to 0 $ 时。于是 $ f $ 对 $ j=1,\ldots,n $ 有连续偏导数 $ \partial_j f=a_j $。现在推论 VII.2.11 表明 $ f $ 属于 $ C^1(X) $，而注 3.3(e) 给出 $ df=\alpha $。\nobreak\hspace{0.35em}■

\end{theorem}

\begin{corollary}

设 $ X $ 是 $ \mathbb{R}^{n} $ 中的开集且是星形的，并设 $ x_0 \in X $。再设 $ q \in \mathbb{N}^{\times}\cup\{\infty\} $ 且 $ \alpha \in \Omega_{(q)}(X) $ 是闭的。令

\[
f(x):=\int_{\Gamma_x}\alpha\quad\text{对 }x\in X ,
\]

其中 $ \Gamma_x $ 是 $ X $ 中以 $ x_0 $ 为起点、以 $ x $ 为终点的分段 $ C^1 $ 曲线。这个函数满足 $ f \in C^{q+1}(X) $ 且 $ df=\alpha $。

\proofhead
庞加莱引理（定理 3.8）与定理 4.4 的证明保证 $ f \in C^1(X) $ 且 $ df=\alpha $。因为对 $ j=1,\ldots,n $ 有 $ \partial_j f=a_j \in C^{q}(X) $，由定理 VII.5.4 可得 $ f $ 属于 $ C^{q+1}(X) $。\nobreak\hspace{0.35em}■
\end{corollary}

这个推论给出了在星形区域上构造闭普法夫形式原函数的办法。在具体计算中，人们选择使所得积分最容易计算的曲线 $ \Gamma_x $（见习题 7）。

\needspace{3\baselineskip}\subsection*{单连通集}

下文中，$ M \subset \mathbb{R}^{n} $ 表示一个非空的道路连通集。$ M $ 中的每条闭连续道路都称为 $ M $ 中的一条\textbf{闭路}。两条闭路\footnote{容易看出，不失一般性，我们可以把 $ \gamma_{0} $ 与 $ \gamma_{1} $ 定义在同一个参数区间上。特别地，可以取 $ I = [0,1] $。} $ \gamma_{0},\gamma_{1} \in C(I,M) $ 称为\textbf{同伦的}，若存在 $ H \in C(I\times[0,1],M) $（一个\textbf{（闭路）同伦}），使得

\begin{subitems}
\item

$ H(\cdot,0) = \gamma_{0} $ 且 $ H(\cdot,1) = \gamma_{1} $；

\item

$ \gamma_{s} := H(\cdot,s) $ 对每个 $ s \in (0,1) $ 都是 $ M $ 中的闭路。

\end{subitems}

\bookfigure[0.33]{figures/994b00cc0d119532.png}

当两条闭路 $ \gamma_{0} $ 与 $ \gamma_{1} $ 同伦时，我们写作 $ \gamma_{0} \sim \gamma_{1} $。对 $ x_{0} \in M $，我们用 $ \gamma_{x_{0}} $ 表示\textbf{常值闭路} $ [t \mapsto x_{0}] $。$ M $ 中与常值闭路同伦的每条闭路都是\textbf{零伦的}。最后，若 $ M $ 中的每条闭路都是零伦的，则称 $ M $ 是\textbf{单连通的}。

\begin{remark}

\begin{items}
\item

在 $ M $ 中所有闭路之集上，$ \sim $ 是等价关系。

\proofhead
显然每条闭路都与自身同伦，即关系 $ \sim $ 是自反的。

\begin{subitems}
\item

设 $ H $ 是从 $ \gamma_{0} $ 到 $ \gamma_{1} $ 的同伦，且

\[
H^{-}(t,s) := H(t,1-s) \quad \text{对 } (t,s) \in I\times[0,1] .
\]

则 $ H^{-} $ 是从 $ \gamma_{1} $ 到 $ \gamma_{0} $ 的同伦。因此关系 $ \sim $ 是对称的。

\item

最后，设 $ H_{0} $ 是从 $ \gamma_{0} $ 到 $ \gamma_{1} $ 的同伦，设 $ H_{1} $ 是从 $ \gamma_{1} $ 到 $ \gamma_{2} $ 的同伦。我们令

\begin{equation*}
H(t,s) := \begin{cases} H_{0}(t,2s) , & (t,s) \in I\times[0,1/2] ,\\ H_{1}(t,2s-1) , & (t,s) \in I\times[1/2,1] . \end{cases}
\end{equation*}

不难验证 $ H $ 属于 $ C(I\times[0,1],M) $；见习题 III.2.13。又，$ H(\cdot,0) = \gamma_{0} $ 且 $ H(\cdot,1) = \gamma_{2} $，且每个 $ H(\cdot,s) $ 都是 $ M $ 中的闭路。\nobreak\hspace{0.35em}■

\end{subitems}

\item

设 $ \gamma $ 是 $ M $ 中的闭路。下列陈述等价：

\begin{subitems}
\item

$ \gamma $ 是零伦的；

\item

对某个 $ x_{0} \in M $，$ \gamma \sim \gamma_{x_{0}} $；

\item

对每个 $ x \in M $，$ \gamma \sim \gamma_{x} $。

\proofhead
只需验证蕴含“(ii)$ \Rightarrow $(iii)”。为此，设 $ \gamma \in C(I,M) $ 是满足对某个 $ x_{0} \in M $ 有 $ \gamma \sim \gamma_{x_{0}} $ 的闭路。又设 $ x \in M $。因 $ M $ 是道路连通的，存在连接 $ x_{0} $ 与 $ x $ 的连续道路 $ w \in C(I,M) $。我们令
\[
H(t,s) := w(s) \quad \text{对 } (t,s) \in I\times[0,1] ,
\]

由此我们看到 $ \gamma_{x_{0}} $ 与 $ \gamma_{x} $ 是同伦的。于是由 (a) 即得结论。\nobreak\hspace{0.35em}■

\end{subitems}

\item

每个星形集都是单连通的。

\proofhead
设 $ M $ 关于 $ x_{0} $ 星形，且 $ \gamma \in C(I,M) $ 是 $ M $ 中的闭路。我们令
\[
H(t,s) := x_{0} + s\left(\gamma(t) - x_{0}\right) \quad \text{对 } (t,s) \in I\times[0,1] .
\]

则 $ H $ 是从 $ \gamma_{x_{0}} $ 到 $ \gamma $ 的同伦。\nobreak\hspace{0.35em}■

\item

集合 $ Q := \mathbb{B}_{2}\setminus T $，其中 $ T := \left((-1/2,1/2)\times\{0\}\right)\cup\left(\{0\}\times(-1,0)\right) $，是单连通的，但不是星形的。

\end{items}

\bookfigure[0.75]{figures/c4b13ae6ec936e9d.png}
\end{remark}

\needspace{3\baselineskip}\subsection*{线积分的同伦不变性}

闭普法夫形式的线积分在闭路同伦下不变：

\begin{proposition}

设 $ \alpha \in \Omega_{(1)}(X) $ 是闭的。设 $ \gamma_{0} $ 与 $ \gamma_{1} $ 是 $ X $ 中同伦的分段 $ C^{1} $ 闭路。则 $ \int_{\gamma_{0}}\alpha = \int_{\gamma_{1}}\alpha $。

证明

\begin{subitems}
\item

设 $ H \in C(I\times[0,1],X) $ 是从 $ \gamma_{0} $ 到 $ \gamma_{1} $ 的同伦。因 $ I\times[0,1] $ 是紧致的，由定理 III.3.6 知 $ K := H(I\times[0,1]) $ 是紧致的。因 $ X^{c} $ 是闭集且 $ X^{c}\cap K = \varnothing $，根据例 III.3.9(c)，存在 $ \varepsilon > 0 $，使得

\[
|H(t,s) - y| \ge \varepsilon \quad \text{对 } (t,s) \in I\times[0,1] \text{ 且 } y \in X^{c} .
\]

定理 III.3.13 保证 $ H $ 是一致连续的。因此存在 $ \delta > 0 $，使得

\[
|H(t,s) - H(\tau,\sigma)| < \varepsilon \quad \text{对 } |t - \tau| < \delta , \quad |s - \sigma| < \delta .
\]
\item

现在取 $ I $ 的一个划分 $ (t_{0},\ldots,t_{m}) $ 与 $ [0,1] $ 的一个划分 $ (s_{0},\ldots,s_{\ell}) $，二者的细度都 $ < \delta $。令 $ A_{j,k} := H(t_{j},s_{k}) $，我们设

\[
\tilde{\gamma}_{k}(t) := A_{j-1,k} + \frac{t - t_{j-1}}{t_{j} - t_{j-1}}\left(A_{j,k} - A_{j-1,k}\right) \quad \text{对 } t_{j-1} \le t \le t_{j} ,
\]

$ 1 \le j \le m $，$ 0 \le k \le \ell $。

\end{subitems}

\bookfigure[0.94]{figures/7764bce3e919a970.png}

显然每个 $ \tilde{\gamma}_{k} $ 都是 $ X $ 中的分段 $ C^{1} $ 闭路。$ \delta $ 的选取表明，我们可以在各点 $ A_{j-1,k-1} $ 的凸邻域 $ \mathbb{B}(A_{j-1,k-1},\varepsilon) $ 中应用庞加莱引理。于是由定理 4.4 得

\[
\int_{\partial V_{j,k}}\alpha = 0 \quad \text{对 } 1 \le j \le m \text{ 且 } 1 \le k \le \ell ,
\]

其中 $ \partial V_{j,k} $ 表示从 $ A_{j-1,k-1} $ 经 $ A_{j,k-1} $、$ A_{j,k} $、$ A_{j-1,k} $ 再回到 $ A_{j-1,k-1} $ 的闭分段直线曲线。因此

\[
\int_{\tilde{\gamma}_{k-1}}\alpha = \int_{\tilde{\gamma}_{k}}\alpha \quad \text{对 } 1 \le k \le \ell ,
\]

因为积分在 $ \tilde{\gamma}_{k-1} $ 与 $ \tilde{\gamma}_{k} $ 之间的“连接段”上相互抵消。同理，用庞加莱引理我们得到 $ \int_{\tilde{\gamma}_{0}}\alpha = \int_{\gamma_{0}}\alpha $ 与 $ \int_{\tilde{\gamma}_{\ell}}\alpha = \int_{\gamma_{\ell}}\alpha $，这正是命题所要求的。\nobreak\hspace{0.35em}■

\end{proposition}

作为同伦不变性定理的一个应用，我们现在得到庞加莱引理的一个影响深远的推广。

\begin{theorem}

设 $ X $ 是 $ \mathbb{R}^{n} $ 中的开集且单连通。若 $ \alpha \in \Omega_{(1)}(X) $ 是闭的，则 $ \alpha $ 是恰当的。

\proofhead
设 $ \gamma $ 是 $ X $ 中的分段 $ C^{1} $ 闭路，且 $ x_{0} \in X $。根据注 4.6(b)，$ \gamma $ 与 $ \gamma_{x_{0}} $ 同伦。于是命题 4.7 与例 4.2(e) 给出 $ \int_{\gamma}\alpha = \int_{\gamma_{x_{0}}}\alpha = 0 $，再由定理 4.4 即得结论。\nobreak\hspace{0.35em}■
\end{theorem}

\begin{example}

\begin{items}
\item

“穿孔平面” $ \mathbb{R}^{2}\setminus\{(0,0)\} $ 不是单连通的。

\proofhead
这由定理 4.8 与例 4.2(d) 推出。\nobreak\hspace{0.35em}■
\item

对 $ n \ge 3 $，集合 $ \mathbb{R}^{n}\setminus\{0\} $ 是单连通的。

\proofhead
见习题 12。\nobreak\hspace{0.35em}■
\end{items}

\end{example}

\begin{remark}

\begin{items}
\item

设 $ X $ 是 $ \mathbb{R}^{n} $ 中的开集且单连通。又设向量场 $ v = (v_{1},\ldots,v_{n}) \in \mathcal{V}^{1}(X) $ 满足可积性条件 $ \partial_{j}v_{k} = \partial_{k}v_{j} $ 对 $ 1 \le j,k \le n $ 成立。则 $ v $ 有势 $ U $，即 $ v $ 是梯度场。若 $ x_{0} $ 是 $ X $ 中任一点，则 $ U $ 可通过下式计算

\[
U(x) := \int_{0}^{1}\left(v(\gamma_{x}(t))\mid\dot{\gamma}_{x}(t)\right)dt \quad \text{对 } x \in X ,
\]

其中 $ \gamma_{x} : [0,1] \to X $ 是 $ X $ 中满足 $ \gamma_{x}(0) = x_{0} $ 与 $ \gamma_{x}(1) = x $ 的分段 $ C^{1} $ 道路。

\proofhead
注 3.3(g) 中定义的典范同构 $ \Theta $ 把普法夫形式 $ \alpha := \Theta v \in \Omega_{(1)}(X) $ 赋给 $ v $。因 $ v $ 满足可积性条件，$ \alpha $ 是闭的，从而由定理 4.8 是恰当的。因在推论 4.5 的证明中，我们可以用定理 4.8 取代庞加莱引理，于是，令 $ \Gamma_{x} := [\gamma_{x}] $，
\[
U(x) := \int_{\Gamma_{x}}\alpha \quad \text{对 } x \in X ,
\]

为 $ \alpha $ 定义了一个势。\nobreak\hspace{0.35em}■

\item

设 $ \boldsymbol{\alpha} = \sum_{j=1}^{n}a_{j}\,dx^{j} \in \Omega_{(0)}(X) $，令 $ \boldsymbol{a} := \Theta^{-1}\boldsymbol{\alpha} = (\alpha_{1},\ldots,\alpha_{n}) \in \mathcal{V}^{0}(X) $ 为相应的向量场。习惯上把 $ \boldsymbol{\alpha} $ 象征性地写成内积

\[
\boldsymbol{\alpha} = \boldsymbol{a}\cdot ds
\]

其中用到了\textbf{向量线元} $ ds := (dx^{1},\ldots,dx^{n}) $。于是，根据注 3.3(g)，线积分

\[
\int_{\Gamma}\boldsymbol{a}\cdot ds := \int_{\Gamma}\Theta\boldsymbol{a}
\]

对每个向量场 $ \boldsymbol{a} \in \mathcal{V}(X) $ 与 $ X $ 中每条分段 $ C^{1} $ 曲线 $ \Gamma $ 都是良定义的。

\item

设 $ \boldsymbol{F} $ 是（三维欧几里得）空间的开子集 $ X $ 中的连续力场（例如静电场或引力场）。若一个（小的）试探粒子沿分段 $ C^{1} $ 曲线 $ \Gamma = [\gamma] $ 运动，则 $ \boldsymbol{F} $ 沿 $ \Gamma $ 所做的\textbf{功}由线积分给出

\[
\int_{\Gamma}\boldsymbol{F}\cdot ds .
\]

特别地，若 $ \gamma $ 是 $ \Gamma $ 的一个分段 $ C^{1} $ 参数化，则

\[
A := \int_{\Gamma}\boldsymbol{F}\cdot ds = \int_{I}\left(F(\gamma(t))\mid\dot{\gamma}(t)\right)dt
\]

是柯西–黎曼积分。因此它由黎曼和近似

\[
\sum_{j=1}^{n}F(\gamma(t_{j-1}))\cdot\dot{\gamma}(t_{j-1})(t_{j} - t_{j-1}) ,
\]

其中 $ \mathfrak{Z} := (t_{0},\ldots,t_{n}) $ 是 $ I $ 的一个适当的划分。因为

\[
\gamma(t_{j}) - \gamma(t_{j-1}) = \dot{\gamma}(t_{j-1})(t_{j} - t_{j-1}) + o(\triangle_{\mathfrak{Z}}) ,
\]

其中 $ \triangle_{\mathfrak{Z}} $ 是 $ \mathfrak{Z} $ 的细度，而和

\[
\sum_{j=1}^{n}F(\gamma(t_{j-1}))\cdot\left(\gamma(t_{j}) - \gamma(t_{j-1})\right)
\]

对细度足够小的划分而言，表示 $ A $ 的一个好的近似，它通过取极限 $ \triangle_{\mathfrak{Z}} \to 0 $ 收敛到 $ A $（见定理 VI.3.4）。内积

\[
F(\gamma(t_{j-1}))\cdot\left(\gamma(t_{j}) - \gamma(t_{j-1})\right)
\]

是（常）力 $ F(\gamma(t_{j-1})) $ 对沿从 $ \gamma(t_{j-1}) $ 到 $ \gamma(t_{j}) $ 的直线段移动的粒子所做的功（“功 = 力 $ \times $ 距离”，或更精确地说，“功 = 力 $ \times $ 沿力的方向的位移”）。
若一个力场 $ \boldsymbol{F} $ 有势，则称它是\textbf{保守的}。此时，根据例 4.2(c)，功是“与路径无关的”，只依赖于初末位置。若 $ X $ 是单连通的（例如 $ X $ 是全空间）且 $ F $ 有连续导数，则根据 (a) 与定理 4.8，$ \boldsymbol{F} $ 保守当且仅当可积性条件成立。若后者成立，则称该力场是\textbf{无旋的}。\footnote{我们将在第 III 卷中澄清这一用语。} \nobreak\hspace{0.35em}■

\end{items}

\end{remark}

\exerciseshead

\begin{exercises}
\item

计算

\[
\int_{\Gamma}(x^{2} - y^{2})\,dx + 3z\,dy + 4xy\,dz
\]

沿螺旋线的一圈。

\item

设 $ \varphi \in C^{1}((0,\infty),\mathbb{R}) $ 且

\[
\alpha(x,y) := x\varphi\left(\sqrt{x^{2} + y^{2}}\right)dy - y\varphi\left(\sqrt{x^{2} + y^{2}}\right)dx \qquad \text{对 } (x,y) \in \mathbb{R}^{2} .
\]

当 $ \Gamma $ 是以 $ 0 $ 为心、半径为 $ R > 0 $ 的正向\footnote{按约定，当圆沿逆时针方向被遍历时，它是正向的（见注 5.8）。}圆时，$ \int_{\Gamma}\alpha $ 的值是多少？

\item

计算

\[
\int_{\Gamma}2xy^{3}\,dx + 3x^{2}y^{2}\,dy ,
\]

其中 $ \Gamma $ 由 $ [\alpha, \beta] \to \mathbb{R}^{2} $，$ t \mapsto (t, t^{2}/2) $ 参数化。

\item

设 $ \Gamma_{+} $ 是 $ \Gamma_{-} $ 是半径为 1 的正向圆，分别以 $ (1,0) $ 和 $ (-1,0) $ 为心。又设

\[
\alpha := \frac{2(x^{2} - y^{2} - 1)\,dy - 4xy\,dx}{(x^{2} + y^{2} - 1)^{2} + 4y^{2}} \in \Omega_{(\infty)}\left(\mathbb{R}^{2}\setminus\{(-1,0),(1,0)\}\right) .
\]

证明

\[
\frac{1}{2\pi}\int_{\Gamma_{\pm}}\alpha = \pm 1 .
\]

（提示：习题 3.3。）

\item

设 $ X $ 是 $ \mathbb{R}^{n} $ 中的开集，$ \Gamma $ 是 $ X $ 中的 $ C^{1} $ 曲线。对 $ q \ge 0 $，证明或否定

\[
\int_{\Gamma} : \Omega_{(q)}(X) \to \mathbb{R} , \quad \alpha \mapsto \int_{\Gamma}\alpha
\]

是一个 $ C^{q}(X) $ 模同态。

\item

设 $ X $ 与 $ Y $ 分别在 $ \mathbb{R}^{n} $ 与 $ \mathbb{R}^{m} $ 中是开集，且 $ \alpha \in \Omega_{(q)}(X\times Y) $。进一步设 $ \Gamma $ 是 $ X $ 中的 $ C^{1} $ 曲线。证明

\[
f : Y \to \mathbb{R} , \quad y \mapsto \int_{\Gamma}\alpha(\cdot,y)
\]

属于 $ C^{q}(Y,\mathbb{R}) $。对 $ q \ge 1 $ 计算 $ \nabla f $。

\item

设 $ -\infty \le \alpha_{j} < \beta_{j} \le \infty $ 对 $ j = 1,\ldots,n $，且 $ X := \prod_{j=1}^{n}(\alpha_{j},\beta_{j}) $。又设 $ \alpha = \sum_{j=1}^{n}a_{j}dx^{j} \in \Omega_{(1)}(X) $ 是闭的。确定（以简单方式）$ \alpha $ 的所有原函数。

\item

设 $ X := (0,\infty)^{2} $ 且 $ a,b,c,e > 0 $。根据习题 3.8(b)，

\[
\alpha := (c - ex)y\,dx + (a - by)x\,dy
\]

具有形如 $ h(x,y) = m(xy) $ 的积分因子 $ h $。确定 $ h\alpha $ 的所有原函数。（提示：习题 7。）

\item

设 $ X $ 是 $ \mathbb{R}^{n} $ 中的开集，$ \Gamma $ 是 $ X $ 中具有参数化 $ \gamma $ 的紧致分段 $ C^{1} $ 曲线。又设 $ f, f_{k} \in C(X,\mathbb{R}^{n}) $ 对 $ k \in \mathbb{N} $。证明

\begin{items}
\item

若 $ (\gamma^{*}f_{k}) $ 一致收敛到 $ \gamma^{*}f $，

\[
\lim_{k}\int_{\Gamma}\sum_{j=1}^{n}f_{k}^{j}\,dx^{j} = \int_{\Gamma}\sum_{j=1}^{n}f^{j}\,dx^{j} ;
\]
\item

若 $ \sum_{k}\gamma^{*}f_{k} $ 一致收敛到 $ \gamma^{*}f $，

\[
\sum_{k}\int_{\Gamma}\sum_{j=1}^{n}f_{k}^{j}\,dx^{j} = \int_{\Gamma}\sum_{j=1}^{n}f^{j}\,dx^{j} .
\]
\end{items}

\item

设 $ a_{j} \in C^{1}(\mathbb{R}^{n},\mathbb{R}) $ 是 $ \lambda \ne -1 $ 次正齐次的，且 1-形式 $ \alpha := \sum_{j=1}^{n}a_{j}dx^{j} $ 是闭的。证明

\[
f(x) := \frac{1}{\lambda + 1}\sum_{j=1}^{n}x^{j}a_{j}(x) \quad \text{对 } x = (x^{1},\ldots,x^{n}) \in \mathbb{R}^{n}
\]

定义了 $ \alpha $ 的一个原函数。（提示：应用欧拉齐次性关系（见习题 VII.3.2）。）

\item

证明：若 $ \gamma \in C(I,X) $ 是一条闭路，则存在一个参数变换 $ \varphi : [0,1] \to I $ 使得 $ \gamma \circ \varphi $ 是一条闭路。

\item

\begin{items}
\item

设 $ X $ 是 $ \mathbb{R}^{n} $ 中的开集，$ \gamma $ 是 $ X $ 中的一条闭路。证明 $ \gamma $ 同伦于 $ X $ 中的一条折线闭路。

\item

证明对 $ a \in \mathbb{R}^{n}\setminus\{0\} $，$ \mathbb{R}^{n}\setminus\mathbb{R}^{+}a $ 是单连通的。

\item

设 $ \gamma $ 是 $ \mathbb{R}^{n}\setminus\{0\} $（$ n \ge 3 $）中的一条折线闭路。证明存在一条半射线 $ \mathbb{R}^{+}a $，其中 $ a \in \mathbb{R}^{n}\setminus\{0\} $，它与 $ \gamma $ 的像不相交。

\item

\proofhead
$ \mathbb{R}^{n}\setminus\{0\} $ 对 $ n \ge 3 $ 是单连通的。
\end{items}

\item

证明或否定：对 $ n \ge 3 $，$ \mathbb{R}^{n}\setminus\{0\} $ 中的每个闭 1-形式都是恰当的。

\item

设 $ X \subset \mathbb{R}^{n} $ 非空且道路连通。对满足 $ \gamma_{1}(1) = \gamma_{2}(0) $ 的 $ \gamma_{1},\gamma_{2} \in C([0,1],X) $，

\begin{equation*}
\gamma_{2}\oplus\gamma_{1} : [0,1] \to X , \quad t \mapsto \begin{cases} \gamma_{1}(2t) , & 0 \le t \le 1/2 , \\ \gamma_{2}(2t - 1) , & 1/2 \le t \le 1 \end{cases}
\end{equation*}

是将 $ \gamma_{2} $ 与 $ \gamma_{1} $ \textbf{连接起来的道路}。对 $ x_{0} \in X $，定义

\[
S_{x_{0}} := \{\, \gamma \in C([0,1],X) ;\ \gamma(0) = \gamma(1) = x_{0} \,\} .
\]

又用 $ \sim $ 表示由 $ S_{x_{0}} $ 的闭路同伦诱导的等价关系（见例 I.4.2(d)）。

验证：

\begin{items}
\item

映射

\[
(S_{x_{0}}/\!\sim)\times(S_{x_{0}}/\!\sim) \to S_{x_{0}}/\!\sim , \quad ([\gamma_{1}],[\gamma_{2}]) \mapsto [\gamma_{2}\oplus\gamma_{1}]
\]

是 $ S_{x_{0}}/\!\sim $ 上的一个良定义的运算，且 $ \Pi_{1}(X,x_{0}) := (S_{x_{0}}/\!\sim,\oplus) $ 是一个群，即 $ X $ 关于 $ x_{0} $ 的\textbf{基本群}，或\textbf{第一同伦群}；

\item

对 $ x_{0},x_{1} \in X $，群 $ \Pi_{1}(X,x_{0}) $ 与 $ \Pi_{1}(X,x_{1}) $ 同构。
\textbf{注} (b) 部分使得我们可以说 $ X $ 的“基本群” $ \Pi_{1}(X) $；见例 I.7.8(g) 中的构造。集合 $ X $ 是单连通的当且仅当 $ \Pi_{1}(X) $ 是平凡的，即它只由单位元组成。
（(b) 的提示：设 $ w \in C([0,1],X) $ 是 $ X $ 中满足 $ w(0) = x_{0} $ 与 $ w(1) = x_{1} $ 的一条道路。则映射 $ [\gamma] \mapsto [w\oplus\gamma\oplus w^{-}] $ 是从 $ \Pi_{1}(X,x_{0}) $ 到 $ \Pi_{1}(X,x_{1}) $ 的群同构。）

\end{items}

\end{exercises}

\section{全纯函数}

线积分在复分析——即复函数的理论——中变得尤为有用。借助上一节的结果，我们几乎毫不费力地就能推出（全局的）柯西积分定理与柯西积分公式。这些定理构成复分析的核心，并具有深远而广泛的影响，其中一部分我们将在本节与下一节中看到。

此处，设

• $ U $ 是 $ \mathbb{C} $ 中的集合，且 $ f : U \to \mathbb{C} $ 连续。

同样在本节中，我们常把 $ z \in U $ 与 $ f $ 分解为实部与虚部，即写作

\[
z = x + iy \in \mathbb{R} + i\mathbb{R} \quad \text{且} \quad f(z) = u(x,y) + iv(x,y) \in \mathbb{R} + i\mathbb{R} .
\]

\needspace{3\baselineskip}\subsection*{复线积分}

设 $ I \subset \mathbb{R} $ 是紧致区间，且设 $ \Gamma $ 是 $ U $ 中由

\[
I \to U , \quad t \mapsto z(t) = x(t) + iy(t) .
\]

参数化的分段 $ C^{1} $ 曲线。则

\[
\int_{\Gamma} f\,dz := \int_{\Gamma} f(z)\,dz := \int_{\Gamma} u\,dx - v\,dy + i\int_{\Gamma} u\,dy + v\,dx
\]

称为 $ f $ 沿 $ \Gamma $ 的\textbf{复线积分}。\footnote{在此语境下，曲线 $ \Gamma $ 也可称为\textbf{围道}。}

\begin{remark}

\begin{items}
\item

我们以 $ \Omega(U,\mathbb{C}) $ 表示\textbf{连续复 1-形式}的空间，并将其定义如下：
在乘积群 $ (\Omega_{(0)}(U)\times\Omega_{(0)}(U),+) $ 上，我们定义一个外乘法

\[
C(U,\mathbb{C})\times\left[\Omega_{(0)}(U)\times\Omega_{(0)}(U)\right] \to \Omega_{(0)}(U)\times\Omega_{(0)}(U) , \quad (f,\alpha) \mapsto f\alpha ,
\]

其中

\[
f\alpha := \left((ua_{1} - vb_{1})\,dx + (ua_{2} - vb_{2})\,dy, (ub_{1} + va_{1})\,dx + (va_{2} + ub_{2})\,dy\right)
\]

对 $ \alpha = (a_{1}\,dx + a_{2}\,dy, b_{1}\,dx + b_{2}\,dy) $。于是立刻可验证

\[
\Omega(U,\mathbb{C}) := (\Omega_{(0)}(U)\times\Omega_{(0)}(U),+,\cdot)
\]

是 $ C(U,\mathbb{C}) $ 上的自由模。此外，我们有

\begin{equation*}
\begin{aligned}
1(a_{1}\,dx + a_{2}\,dy,0) &= (a_{1}\,dx + a_{2}\,dy,0) , \\ i(a_{1}\,dx + a_{2}\,dy,0) &= (0, a_{1}\,dx + a_{2}\,dy)
\end{aligned}
\end{equation*}

对 $ a_{1}\,dx + a_{2}\,dy \in \Omega_{(0)}(U) $。因此我们可将 $ \Omega_{(0)}(U) $ 等同于 $ \Omega(U,\mathbb{C}) $ 中的 $ \Omega_{(0)}(U)\times\{0\} $，并把 $ (a_{1}\,dx + a_{2}\,dy, b_{1}\,dx + b_{2}\,dy) \in \Omega(U,\mathbb{C}) $ 唯一地表示成

\[
a_{1}\,dx + a_{2}\,dy + i(b_{1}\,dx + b_{2}\,dy) ,
\]

我们通过记号

\[
\Omega(U,\mathbb{C}) = \Omega_{(0)}(U) + i\Omega_{(0)}(U) .
\]

来表示。最后，我们有

\begin{equation*}
\begin{aligned}
(a_{1} + ib_{1})(dx,0) &= (a_{1}\,dx, b_{1}\,dx) , \\ (a_{2} + ib_{2})(dy,0) &= (a_{2}\,dy, b_{2}\,dy) ,
\end{aligned}
\end{equation*}

从而

\[
(a_{1} + ib_{1})(dx,0) + (a_{2} + ib_{2})(dy,0) = (a_{1}\,dx + a_{2}\,dy, b_{1}\,dx + b_{2}\,dy) .
\]

这意味着我们也可以把 $ (a_{1}\,dx + a_{2}\,dy, b_{1}\,dx + b_{2}\,dy) \in \Omega(U,\mathbb{C}) $ 唯一地写成

\[
(a_{1} + ib_{1})\,dx + (a_{2} + ib_{2})\,dy ,
\]

即

\[
\Omega(U,\mathbb{C}) = \{\, a\,dx + b\,dy ;\ a,b \in C(U,\mathbb{C}) \,\} .
\]
\item

令 $ u_{x} := \partial_{1}u $ 等，我们把

\[
df := (u_{x} + iv_{x})\,dx + (u_{y} + iv_{y})\,dy \in \Omega(U,\mathbb{C})
\]

称为 $ f $ 的\textbf{复微分}。显然，$ dz = dx + i\,dy $，且我们得到\footnote{将 (5.1) 与 $ f $ 的复线积分的定义相比较。}

\begin{equation}
f\,dz = (u + iv)(dx + i\,dy) = u\,dx - v\,dy + i(u\,dy + v\,dx)\tag{5.1}
\end{equation}

对 $ f = u + iv \in C(U,\mathbb{C}) $。\nobreak\hspace{0.35em}■

\end{items}

\end{remark}

\begin{proposition}

设 $ \Gamma $ 是由 $ I \to U $，$ t \mapsto z(t) $ 参数化的分段 $ C^{1} $ 曲线。则

\begin{subitems}
\item

$ \int_{\Gamma}f(z)\,dz = \int_{I}f(z(t))\dot{z}(t)\,dt $；

\item

$ \left|\int_{\Gamma}f(z)\,dz\right| \le \max_{z\in\Gamma}|f(z)|\,L(\Gamma) $。

\proofhead
(i) 对分段 $ C^{1} $ 道路 $ \gamma : I \to \mathbb{R}^{2} $，$ t \mapsto (x(t),y(t)) $，我们有
\begin{equation*}
\begin{aligned}
\int_{\Gamma} f\,dz &= \int_{I}\gamma^{*}(u\,dx - v\,dy) + i\int_{I}\gamma^{*}(u\,dy + v\,dx) \\ &= \int_{I}\left[(u\circ\gamma)\dot{x} - (v\circ\gamma)\dot{y}\right]dt + i\int_{I}\left[(u\circ\gamma)\dot{y} + (v\circ\gamma)\dot{x}\right]dt \\ &= \int_{I}(u\circ\gamma + iv\circ\gamma)(\dot{x} + i\dot{y})\,dt \\ &= \int_{I}f(z(t))\dot{z}(t)\,dt .
\end{aligned}
\end{equation*}
\end{subitems}

(ii) 该断言由 (i) 推出：利用定理 1.3 与命题 VI.4.3 估计最后一个积分。\nobreak\hspace{0.35em}■

\end{proposition}

\begin{example}

\begin{items}
\item

设 $ z_{0} \in \mathbb{C} $ 且 $ r > 0 $，并令 $ z(t) := z_{0} + re^{it} $ 对 $ t \in [0,2\pi] $。则对 $ \Gamma := [z] $，我们有

\begin{equation*}
\int_{\Gamma}(z - z_{0})^{m}\,dz = \begin{cases} 0 , & m \in \mathbb{Z}\setminus\{-1\} , \\ 2\pi i , & m = -1 . \end{cases}
\end{equation*}

\proofhead
因为
\[
\int_{\Gamma}(z - z_{0})^{m}\,dz = \int_{0}^{2\pi}r^{m}e^{itm}ire^{it}\,dt = ir^{m+1}\int_{0}^{2\pi}e^{i(m+1)t}\,dt ,
\]

该断言由指数函数的 $ 2\pi i $-周期性推出。\nobreak\hspace{0.35em}■

\item

设 $ \Gamma $ 同 (a) 中一样，取 $ z_{0} = 0 $。则对 $ z \in \mathbb{C} $，我们有

\[
\frac{1}{2\pi i}\int_{\Gamma}\lambda^{-k-1}e^{\lambda z}\,d\lambda = \frac{z^{k}}{k!} \quad \text{对 } k \in \mathbb{N} .
\]

\proofhead
由习题 4.9 可得
\[
\int_{\Gamma}\lambda^{-k-1}e^{\lambda z}\,d\lambda = \sum_{n=0}^{\infty}\frac{z^{n}}{n!}\int_{\Gamma}\lambda^{n-k-1}\,d\lambda ,
\]

而该断言由 (a) 推出。\nobreak\hspace{0.35em}■

\item

设 $ \Gamma $ 是 $ \mathbb{R} $ 中由 $ I \to \mathbb{C} $，$ t \mapsto t $ 参数化的曲线。则

\[
\int_{\Gamma}f(z)\,dz = \int_{I}f(t)\,dt ,
\]

即在此情形下复线积分与第 VI 章的柯西–黎曼积分一致。

\proofhead
这由命题 5.2(i) 推出。\nobreak\hspace{0.35em}■
\item

对分别由

\[
[0,\pi] \to \mathbb{C} , \quad t \mapsto e^{i(\pi - t)} ,
\]

与

\[
[-1,1] \to \mathbb{C} , \quad t \mapsto t ,
\]

参数化的曲线 $ \Gamma_{1} $ 与 $ \Gamma_{2} $，

\end{items}

\bookfigure[0.41]{figures/77f1d4126962389d.png}

我们有

\[
\int_{\Gamma_{1}}|z|\,dz = -i\int_{0}^{\pi}e^{i(\pi - t)}\,dt = e^{i(\pi - t)}\big|_{0}^{\pi} = 2
\]

且

\[
\int_{\Gamma_{2}}|z|\,dz = \int_{-1}^{1}|t|\,dt = 1 .
\]

因此复线积分一般依赖于积分路径，而不仅依赖于起点与终点。\nobreak\hspace{0.35em}■

\end{example}

\needspace{3\baselineskip}\subsection*{全纯性}

\textbf{全纯}函数 $ f $ 是具有连续复导数的函数，即 $ f \in C^{1}(U,\mathbb{C}) $。若

\[
U \to \mathbb{R}^{2} , \quad (x,y) \mapsto (u(x,y),v(x,y))
\]

属于 $ C^{1}(U,\mathbb{R}^{2}) $（见节 VII.2），则我们称 $ f $ \textbf{连续实可微}。

\begin{remark}

\begin{items}
\item

在本节末尾我们将证明连续复导数这一假设是不必要的，即每个复可微函数都有连续导数，从而都是全纯的。这里我们采用更强的连续复可微假设，是因为它能让我们更快地到达\textbf{复分析}——即复变量的复值函数理论——的基本定理。

\item

函数 $ f $ 全纯，当且仅当它连续实可微且满足柯西–黎曼方程

\[
u_{x} = v_{y} \quad \text{且} \quad u_{y} = -v_{x} .
\]

在此情形下，$ f' = u_{x} + iv_{x} $。

\proofhead
这是定理 VII.2.18 的重新表述。\nobreak\hspace{0.35em}■
\item

根据注 5.1(b)，我们有

\[
f\,dz = u\,dx - v\,dy + i(u\,dy + v\,dx) .
\]

若 $ f $ 全纯，则柯西–黎曼方程表明 1-形式 $ u\,dx - v\,dy $ 与 $ u\,dy + v\,dx $ 都是闭的。

\item

设 $ U $ 是区域且 $ f $ 全纯。则 $ f $ 为常数当且仅当满足下列条件：\footnote{见习题 V.3.5。}

\end{items}

\begin{subitems}
\item

$ u = \mathrm{const} $；

\item

$ v = \mathrm{const} $；

\item

$ \overline{f} $ 全纯；

\item

$ |f| = \mathrm{const} $。

\proofhead
若 $ f $ 为常数，则 (i)–(iv) 显然成立。
因为 $ f' = u_{x} + iv_{x} $，由柯西–黎曼方程与注 VII.3.11(c) 可得 (i) 或 (ii) 蕴含 $ f $ 为常数。
\end{subitems}

\begin{subitems}
\item[（iii）]

若 $ f $ 与 $ \overline{f} $ 全纯，则 $ u = (f + \overline{f})/2 $ 也全纯。由于 $ \operatorname{Im}(u) = 0 $，由 (ii) 可得 $ u $ 为常数，于是由 (i) 可得 $ f $ 也为常数。

\item[（iv）]

只需考虑 $ f \ne 0 $ 的情形。因为 $ |f| $ 为常数，$ f $ 处处不为零。因此 $ 1/f $ 在 $ U $ 上有定义且全纯。于是 $ \overline{f} = |f|^{2}/f $ 在 $ U $ 中也全纯，而该断言由 (iii) 推出。\nobreak\hspace{0.35em}■

\end{subitems}

\end{remark}

\needspace{3\baselineskip}\subsection*{柯西积分定理}

我们在例 5.3(d) 中已经看到，复线积分一般依赖的不只是积分路径的起点与终点。然而，在被积函数全纯的情形下，前面各节的定理给出下列关于路径无关性的重要结果。

\begin{theorem}

（柯西积分定理）设 $ U $ 单连通且 $ f $ 全纯。则对 $ U $ 中每条闭的分段 $ C^{1} $ 曲线 $ \Gamma $，

\[
\int_{\Gamma}f\,dz = 0 .
\]

\proofhead
根据注 5.4(c)，1-形式 $ \alpha_{1} := u\,dx - v\,dy $ 与 $ \alpha_{2} := u\,dy + v\,dx $ 都是闭的。因为 $ U $ 单连通，由定理 4.8 可得 $ \alpha_{1} $ 与 $ \alpha_{2} $ 是恰当的。于是定理 4.4 蕴含
\[
\int_{\Gamma}f\,dz = \int_{\Gamma}\alpha_{1} + i\int_{\Gamma}\alpha_{2} = 0
\]

对 $ U $ 中每条闭的分段 $ C^{1} $ 曲线 $ \Gamma $。\nobreak\hspace{0.35em}■

\end{theorem}

\begin{theorem}

设 $ U $ 单连通且 $ f $ 全纯。则 $ f $ 有全纯原函数。$ f $ 的每个原函数 $ \varphi $ 满足

\[
\int_{\Gamma}f\,dz = \varphi(E_{\Gamma}) - \varphi(A_{\Gamma})
\]

对 $ U $ 中每条分段 $ C^{1} $ 曲线 $ \Gamma $。

\proofhead
我们沿用前面证明中的记号。因为 $ \alpha_{1} $ 与 $ \alpha_{2} $ 是恰当的，存在 $ h_{1}, h_{2} \in C^{2}(U,\mathbb{R}) $ 使得 $ dh_{1} = \alpha_{1} $ 且 $ dh_{2} = \alpha_{2} $。由此我们读出
\[
(h_{1})_{x} = u , \quad (h_{1})_{y} = -v , \quad (h_{2})_{x} = v , \quad (h_{2})_{y} = u .
\]

因此 $ \varphi := h_{1} + ih_{2} $ 满足柯西–黎曼方程。于是 $ \varphi $ 全纯，且

\[
\varphi' = (h_{1})_{x} + i(h_{2})_{x} = u + iv = f .
\]

这表明 $ \varphi $ 是 $ f $ 的原函数。第二个断言由例 4.2(b) 推出。\nobreak\hspace{0.35em}■

\end{theorem}

\begin{proposition}

设 $ f $ 全纯，且 $ \gamma_{1} $ 与 $ \gamma_{2} $ 是 $ U $ 中同伦的分段 $ C^{1} $ 闭路。则

\[
\int_{\gamma_{1}}f\,dz = \int_{\gamma_{2}}f\,dz .
\]

\proofhead
这由注 5.4(c) 与命题 4.7 推出（也可参见定理 5.5 的证明）。\nobreak\hspace{0.35em}■
\end{proposition}

\needspace{3\baselineskip}\subsection*{圆周的定向}

\begin{remark}

我们回忆记号 $ \mathbb{D}(a,r) = a + r\mathbb{D} $ 表示 $ \mathbb{C} $ 中以 $ a \in \mathbb{C} $ 为圆心、$ r > 0 $ 为半径的开圆盘。下文中，我们以 $ \partial\mathbb{D}(a,r) = a + r\partial\mathbb{D} $ 表示以 $ a $ 为圆心、$ r $ 为半径的\textbf{正向圆周}。这意味着 $ \partial\mathbb{D}(a,r) $ 是曲线 $ \Gamma = [\gamma] $，其中 $ \gamma : [0,2\pi] \to \mathbb{C} $，$ t \mapsto a + re^{it} $。在这个定向下，圆周被这样遍历：圆盘 $ \mathbb{D}(a,r) $ 保持在左侧，或者说它被逆时针遍历。这等价于说负单位法向量 $ -\mathbf{n} $ 总是指向外部。

\proofhead
由注 2.5(c) 与 $ \mathbb{C} $ 与 $ \mathbb{R}^{2} $ 的典范等同，可得 Frenet 二标架由 $ \mathbf{e}_{1} = (-\sin,\cos) $ 与 $ \mathbf{e}_{2} = (-\cos,-\sin) $ 给出。令 $ x(t) = a + r(\cos t,\sin t) \in \partial\mathbb{D}(a,r) $，我们有关系式
\[
x(t) + re_{2}(t) = a \quad \text{对 } 0 \le t \le 2\pi .
\]

于是 $ e_{2}(t) $ 沿 $ \partial\mathbb{D}(a,r) $ 指向内部，而其负向量指向外部。\nobreak\hspace{0.35em}■

\end{remark}

\needspace{3\baselineskip}\subsection*{柯西积分公式}

全纯函数有引人注目的积分表示，我们在下一个定理中给出。这一公式的含义将在后面的应用中进一步探讨。

\begin{theorem}

（柯西积分公式）设 $ f $ 全纯且 $ \overline{\mathbb{D}}(z_{0},r) \subset U $。则

\[
f(z) = \frac{1}{2\pi i}\int_{\partial\mathbb{D}(z_{0},r)}\frac{f(\zeta)}{\zeta - z}\,d\zeta \quad \text{对 } z \in \mathbb{D}(z_{0},r) .
\]

\proofhead
设 $ z \in \mathbb{D}(z_{0},r) $ 且 $ \varepsilon > 0 $。则存在 $ \delta > 0 $ 使得 $ \overline{\mathbb{D}}(z,\delta) \subset U $ 且
\begin{equation}
|f(\zeta) - f(z)| \le \varepsilon \quad \text{对 } \zeta \in \overline{\mathbb{D}}(z,\delta) .\tag{5.2}
\end{equation}

我们令 $ \Gamma_{\delta} := \partial\mathbb{D}(z,\delta) $ 且 $ \Gamma := \partial\mathbb{D}(z_{0},r) $。习题 6 表明 $ \Gamma_{\delta} $ 与 $ \Gamma $ 在 $ U\setminus\{z\} $ 中同伦。此外，

\[
U\setminus\{z\} \to \mathbb{C} , \quad \zeta \mapsto 1/(\zeta - z)
\]

全纯。因此由命题 5.7 与例 5.3(a) 可得

\begin{equation}
\int_{\Gamma}\frac{d\zeta}{\zeta - z} = \int_{\Gamma_{\delta}}\frac{d\zeta}{\zeta - z} = 2\pi i .\tag{5.3}
\end{equation}

因为

\[
U\setminus\{z\} \to \mathbb{C} , \quad \zeta \mapsto (f(\zeta) - f(z))/(\zeta - z)
\]

也全纯，由命题 5.7 我们知道

\begin{equation}
\int_{\Gamma}\frac{f(\zeta) - f(z)}{\zeta - z}\,d\zeta = \int_{\Gamma_{\delta}}\frac{f(\zeta) - f(z)}{\zeta - z}\,d\zeta .\tag{5.4}
\end{equation}

将 (5.3) 与 (5.4) 结合，我们得到

\begin{equation*}
\begin{aligned}
\frac{1}{2\pi i}\int_{\Gamma}\frac{f(\zeta)}{\zeta - z}\,d\zeta &= \frac{1}{2\pi i}\int_{\Gamma}\frac{f(z)}{\zeta - z}\,d\zeta + \frac{1}{2\pi i}\int_{\Gamma}\frac{f(\zeta) - f(z)}{\zeta - z}\,d\zeta \\ &= f(z) + \frac{1}{2\pi i}\int_{\Gamma_{\delta}}\frac{f(\zeta) - f(z)}{\zeta - z}\,d\zeta .
\end{aligned}
\end{equation*}

然后命题 5.2(ii) 与 (5.2) 蕴含估计

\[
\left|\frac{1}{2\pi i}\int_{\Gamma_{\delta}}\frac{f(\zeta) - f(z)}{\zeta - z}\,d\zeta\right| \le \frac{\varepsilon}{2\pi\delta}2\pi\delta = \varepsilon ,
\]

并且我们得到

\[
\left|\frac{1}{2\pi i}\int_{\Gamma}\frac{f(\zeta)}{\zeta - z}\,d\zeta - f(z)\right| \le \varepsilon .
\]

因为 $ \varepsilon > 0 $ 是任意的，该断言得证。\nobreak\hspace{0.35em}■

\end{theorem}

\begin{remark}

\begin{items}
\item

在定理 5.9 的假设下，我们有

\[
f(z) = \frac{r}{2\pi}\int_{0}^{2\pi}\frac{f(z_{0} + re^{it})}{z_{0} + re^{it} - z}e^{it}\,dt \quad \text{对 } z \in \mathbb{D}(z_{0},r)
\]

特别地，

\[
f(z_{0}) = \frac{1}{2\pi}\int_{0}^{2\pi}f(z_{0} + re^{it})\,dt .
\]
\item

若对每个 $ z_{0} \in U $ 都存在 $ r_{0} > 0 $ 使得

\[
f(z_{0}) = \frac{1}{2\pi}\int_{0}^{2\pi}f(z_{0} + re^{it})\,dt \quad \text{对 } r \in [0,r_{0}] ,
\]

则函数 $ f $ 具有\textbf{平均值性质}。
由 (a) 可得全纯函数具有平均值性质。

\item

若 $ f $ 具有平均值性质，则 $ \operatorname{Re} f $、$ \operatorname{Im} f $、$ \overline{f} $ 以及 $ \lambda f $（$ \lambda \in \mathbb{C} $）也具有平均值性质。若 $ g \in C(U,\mathbb{C}) $ 也具有平均值性质，则 $ f + g $ 也具有平均值性质。

\proofhead
这由定理 III.1.5 与 III.1.10 以及推论 VI.4.10 推出。\nobreak\hspace{0.35em}■
\item

柯西积分公式对本质上一般的曲线也成立，我们将在下一节证明。\nobreak\hspace{0.35em}■

\end{items}

\end{remark}

\needspace{3\baselineskip}\subsection*{解析函数}

作为柯西积分公式的另一推论，我们现在证明基本定理，即全纯函数都是解析的。我们如下进行：设 $ z_{0} \in U $ 且 $ r > 0 $ 满足 $ \overline{\mathbb{D}}(z_{0},r) \subset U $。则函数

\[
\mathbb{D}(z_{0},r) \to \mathbb{C} , \quad z \mapsto f(\zeta)/(\zeta - z)
\]

对每个 $ \zeta \in \partial\mathbb{D}(z_{0},r) $ 都是解析的，且可作幂级数展开。借助柯西积分公式，我们可以把这个展开转移到 $ f $ 上。

\qyxffntextdef{1}{另见习题 4.9。}
\begin{theorem}

函数 $ f $ 是全纯的当且仅当它是解析的。因此

\[
C^{1}(U,\mathbb{C}) = C^{\omega}(U,\mathbb{C}) .
\]

证明

\begin{subitems}
\item

设 $ f $ 是全纯的。令 $ z_{0} \in U $ 且 $ r > 0 $ 满足 $ \overline{\mathbb{D}}(z_{0},r) \subset U $。取 $ z \in \mathbb{D}(z_{0},r) $ 并令 $ r_{0} := |z - z_{0}| $。则 $ r_{0} < r $，且

\[
\frac{|z - z_{0}|}{|\zeta - z_{0}|} = \frac{r_{0}}{r} < 1 \quad \text{对 } \zeta \in \Gamma := \partial\mathbb{D}(z_{0},r) .
\]

因此几何级数 $ \sum a^{k} $（其中 $ a := (z - z_{0})/(\zeta - z_{0}) $）收敛，其值为

\[
\sum_{k=0}^{\infty}\left(\frac{z - z_{0}}{\zeta - z_{0}}\right)^{k} = \frac{1}{1 - (z - z_{0})/(\zeta - z_{0})} = \frac{\zeta - z_{0}}{\zeta - z} .
\]

于是

\begin{equation}
\frac{f(\zeta)}{\zeta - z} = \frac{f(\zeta)}{\zeta - z_{0}}\sum_{k=0}^{\infty}\left(\frac{z - z_{0}}{\zeta - z_{0}}\right)^{k} = \sum_{k=0}^{\infty}\frac{f(\zeta)}{(\zeta - z_{0})^{k+1}}(z - z_{0})^{k} .\tag{5.5}
\end{equation}

因为 $ \Gamma $ 是紧致的，存在 $ M \ge 0 $ 使得对 $ \zeta \in \Gamma $ 有 $ |f(\zeta)| \le M $。由此可得

\[
\left|\frac{f(\zeta)}{(\zeta - z_{0})^{k+1}}(z - z_{0})^{k}\right| \le \frac{M}{r^{k+1}}r_{0}^{k} = \frac{M}{r}\left(\frac{r_{0}}{r}\right)^{k} \quad \text{对 } \zeta \in \Gamma .
\]

因为 $ r_{0}/r < 1 $，魏尔斯特拉斯优级数判别法（定理 V.1.6）表明 (5.5) 中的级数关于 $ \zeta \in \Gamma $ 一致收敛。令

\begin{equation}
a_{k} := \frac{1}{2\pi i}\int_{\Gamma}\frac{f(\zeta)}{(\zeta - z_{0})^{k+1}}\,d\zeta \quad \text{对 } k \in \mathbb{N} ,\tag{5.6}
\end{equation}

由命题 VI.4.1 可得\qyxffnmark{1}

\begin{equation*}
\begin{aligned}
f(z) &= \frac{1}{2\pi i}\int_{\Gamma}\frac{f(\zeta)}{\zeta - z}\,d\zeta = \frac{1}{2\pi i}\int_{\Gamma}\sum_{k=0}^{\infty}\frac{f(\zeta)}{(\zeta - z_{0})^{k+1}}(z - z_{0})^{k}\,d\zeta \\ &= \sum_{k=0}^{\infty}\left[\frac{1}{2\pi i}\int_{\Gamma}\frac{f(\zeta)}{(\zeta - z_{0})^{k+1}}\,d\zeta\right](z - z_{0})^{k} = \sum_{k=0}^{\infty}a_{k}(z - z_{0})^{k} .
\end{aligned}
\end{equation*}

因为这对每个 $ z_{0} \in U $ 都成立，我们发现 $ f $ 可在 $ U $ 中每点的邻域内展成收敛幂级数。因此 $ f $ 是解析的。

\item

当 $ f $ 是解析的时，它具有连续复导数，从而是全纯的。\nobreak\hspace{0.35em}■

\end{subitems}

\end{theorem}

\begin{corollary}

（柯西导数公式）设 $ f $ 是全纯的，$ z \in U $，且 $ r > 0 $ 满足 $ \overline{\mathbb{D}}(z,r) \subset U $。则

\[
f^{(n)}(z) = \frac{n!}{2\pi i}\int_{\partial\mathbb{D}(z,r)}\frac{f(\zeta)}{(\zeta - z)^{n+1}}\,d\zeta \quad \text{对 } n \in \mathbb{N} .
\]

\proofhead
由定理 5.11、注 V.3.4(b) 以及解析函数的唯一性定理可得，$ f $ 在 $ \mathbb{D}(z,r) $ 中由它的泰勒级数表示，即 $ f = \mathcal{T}(f,z) $。结论现在由 (5.6) 和幂级数的唯一性定理推出。\nobreak\hspace{0.35em}■
\end{corollary}

\begin{remark}

\begin{items}
\item

设 $ f $ 是全纯的，$ z \in U $。则上面的证明表明泰勒级数 $ \mathcal{T}(f,z) $ 至少在完全落在 $ U $ 内的最大圆盘上重现函数 $ f $。

\item

若 $ f $ 是全纯的，则 $ u $ 与 $ v $ 属于 $ C^{\infty}(U,\mathbb{R}) $。

\proofhead
这由定理 5.11 推出。\nobreak\hspace{0.35em}■
\item

若 $ f $ 是解析的，则 $ 1/f $ 也是（在 $ U\setminus f^{-1}(0) $ 中）。

\proofhead
由商法则，$ 1/f $ 在 $ U\setminus f^{-1}(0) $ 中有连续导数，因此是全纯的。于是结论由定理 5.11 推出。\nobreak\hspace{0.35em}■
\item

设 $ V $ 是 $ \mathbb{C} $ 中的开集。若 $ f : U \to \mathbb{C} $ 与 $ g : V \to \mathbb{C} $ 是解析的且 $ f(U) \subset V $，则复合 $ g \circ f : U \to \mathbb{C} $ 也是解析的。

\proofhead
这由链式法则和定理 5.11 推出。\nobreak\hspace{0.35em}■
\end{items}

\end{remark}

\needspace{3\baselineskip}\subsection*{刘维尔定理}

在全部 $ \mathbb{C} $ 上全纯的函数称为\textbf{整函数}。

\begin{theorem}

（刘维尔）每个有界整函数都是常值函数。

\proofhead
根据注 5.13(a)，我们有
\[
f(z) = \sum_{k=0}^{\infty}\frac{f^{(k)}(0)}{k!}z^{k} \quad \text{对 } z \in \mathbb{C} .
\]

由假设，存在 $ M < \infty $ 使得对 $ z \in \mathbb{C} $ 有 $ |f(z)| \le M $。于是由命题 5.2(ii) 和推论 5.12 可得

\[
\left|\frac{f^{(k)}(0)}{k!}\right| \le \frac{M}{r^{k}} \quad \text{对 } r > 0 .
\]

对 $ k \ge 1 $，令 $ r \to \infty $ 取极限即得 $ f^{(k)}(0) = 0 $。因此对每个 $ z \in \mathbb{C} $，$ f(z) $ 等于 $ f(0) $。\nobreak\hspace{0.35em}■

\end{theorem}

\begin{application}

刘维尔定理给出代数基本定理的一个简单证明，即 $ \mathbb{C}[X] $ 上的每个非常值多项式都有零点。

\proofhead
我们把 $ p \in \mathbb{C}[X] $ 写成 $ p(z) = \sum_{k=0}^{n}a_{k}z^{k} $ 的形式，其中 $ n \ge 1 $ 且 $ a_{n} \ne 0 $。因为
\[
p(z) = z^{n}\left(a_{n} + \frac{a_{n-1}}{z} + \cdots + \frac{a_{0}}{z^{n}}\right) ,
\]

当 $ |z| \to \infty $ 时有 $ |p(z)| \to \infty $。于是存在 $ R > 0 $ 使得对 $ z \notin R\overline{\mathbb{D}} $ 有 $ |p(z)| \ge 1 $。我们假设 $ p $ 在 $ R\overline{\mathbb{D}} $ 中没有零点。因为 $ R\overline{\mathbb{D}} $ 是紧致的，由极值定理（推论 III.3.8）可得存在正数 $ \varepsilon $ 使得 $ |p(z)| \ge \varepsilon $。因此 $ 1/p $ 是整函数，且对 $ z \in \mathbb{C} $ 满足 $ |1/p(z)| \le \max\{1,1/\varepsilon\} $。刘维尔定理现在推出 $ 1/p $ 是常值的，从而 $ p $ 也是常值的，矛盾。\nobreak\hspace{0.35em}■

\end{application}

\needspace{3\baselineskip}\subsection*{菲涅耳积分}

柯西积分定理也可用来计算被积函数与全纯函数相关的实积分。我们可以把具有实积分限的积分看作实路径上的复线积分。于是，柯西积分定理所保证的选择积分曲线的更大自由，将使我们能够计算出许多新的积分。

下一个定理展示这一方法，我们将在后续各节中推广这些技巧。从习题中还可以学到更多。

\begin{proposition}

下列反常菲涅耳积分收敛，并且

\[
\int_{0}^{\infty}\cos(t^{2})\,dt = \int_{0}^{\infty}\sin(t^{2})\,dt = \frac{\sqrt{2\pi}}{4} .
\]

\proofhead
这些积分的收敛性由习题 VI.8.7 推出。

考虑整函数 $ z \mapsto e^{-z^{2}} $ 沿闭分段 $ C^{1} $ 曲线 $ \Gamma = \Gamma_{1} + \Gamma_{2} + \Gamma_{3} $ 的积分，该曲线沿直线段从 0 到 $ \alpha > 0 $，再到 $ \alpha + i\alpha $，最后回到 0。
\bookfigure[0.31]{figures/69e0ce7c02b19f54.png}

由柯西积分定理，我们有

\[
-\int_{\Gamma_{3}}e^{-z^{2}}\,dz = \int_{\Gamma_{1}}e^{-z^{2}}\,dz + \int_{\Gamma_{2}}e^{-z^{2}}\,dz .
\]

应用 VI.9.7 表明

\[
\int_{\Gamma_{1}}e^{-z^{2}}\,dz = \int_{0}^{\alpha}e^{-t^{2}}\,dt \to \int_{0}^{\infty}e^{-t^{2}}\,dt = \frac{\sqrt{\pi}}{2} \quad (\alpha \to \infty) .
\]

沿 $ \Gamma_{2} $ 的积分可估计为

\[
\left|\int_{\Gamma_{2}}e^{-z^{2}}\,dz\right| = \left|\int_{0}^{\alpha}e^{-(\alpha + it)^{2}}i\,dt\right| \le \int_{0}^{\alpha}e^{-\operatorname{Re}(\alpha + it)^{2}}\,dt = e^{-\alpha^{2}}\int_{0}^{\alpha}e^{t^{2}}\,dt .
\]

注意到

\[
\int_{0}^{\alpha}e^{t^{2}}\,dt \le \int_{0}^{\alpha}e^{\alpha t}\,dt = \frac{1}{\alpha}(e^{\alpha^{2}} - 1) ,
\]

我们发现

\[
\left|\int_{\Gamma_{2}}e^{-z^{2}}\,dz\right| \le \frac{1}{\alpha}(1 - e^{-\alpha^{2}}) \to 0 \quad (\alpha \to \infty) .
\]

于是我们得到

\begin{equation}
\lim_{\alpha\to\infty}\left(-\int_{\Gamma_{3}}e^{-z^{2}}\,dz\right) = \frac{\sqrt{\pi}}{2} .\tag{5.7}
\end{equation}

用 $ -\Gamma_{3} $ 的参数化 $ t \mapsto t + it $，我们有

\begin{equation*}
\begin{aligned}
-\int_{\Gamma_{3}}e^{-z^{2}}\,dz &= \int_{0}^{\alpha}e^{-(1+i)^{2}t^{2}}(1+i)\,dt = (1+i)\int_{0}^{\alpha}e^{-2it^{2}}\,dt \\ &= (1+i)\left(\int_{0}^{\alpha}\cos(2t^{2})\,dt - i\int_{0}^{\alpha}\sin(2t^{2})\,dt\right) .
\end{aligned}
\end{equation*}

由代换 $ \sqrt{2}\,t = \tau $ 和极限 $ \alpha \to \infty $，利用 (5.7) 我们发现

\[
\int_{0}^{\infty}\cos(\tau^{2})\,d\tau - i\int_{0}^{\infty}\sin(\tau^{2})\,d\tau = \frac{\sqrt{2}\sqrt{\pi}}{2(1+i)} = \frac{\sqrt{2\pi}}{4}(1 - i) .\quad\blacksquare
\]
\end{proposition}

\needspace{3\baselineskip}\subsection*{最大模原理}

我们已经在注 5.10(b) 中看到全纯函数具有平均值性质。特别地，全纯函数 $ f $ 在圆盘中心的绝对值不可能大于 $ f $ 在边界上的最大绝对值。

\begin{theorem}

（广义最大模原理）设函数 $ f $ 具有平均值性质。若 $ |f| $ 在 $ z_{0} \in U $ 处有局部极大值，则 $ f $ 在 $ z_{0} $ 的一个邻域内是常值的。

证明

\begin{subitems}
\item

情形 $ f(z_{0}) = 0 $ 是显然的。因为 $ f(z_{0}) \ne 0 $，存在 $ c $ 使得 $ |c| = 1 $ 且 $ cf(z_{0}) > 0 $。因为 $ cf $ 也具有平均值性质，我们可以不失一般性地假设 $ f(z_{0}) $ 是正实数。由假设，存在 $ r_{0} > 0 $ 使得 $ \overline{\mathbb{D}}(z_{0},r_{0}) \subset U $ 且对 $ z \in \overline{\mathbb{D}}(z_{0},r_{0}) $ 有 $ |f(z)| \le f(z_{0}) $。此外

\[
f(z_{0}) = \frac{1}{2\pi}\int_{0}^{2\pi}f(z_{0} + re^{it})\,dt \quad \text{对 } r \in [0,r_{0}] .
\]
\item

函数 $ h : U \to \mathbb{R} $，$ z \mapsto \operatorname{Re} f(z) - f(z_{0}) $ 满足 $ h(z_{0}) = 0 $ 且

\[
h(z) \le |f(z)| - f(z_{0}) \le 0 \quad \text{对 } z \in \mathbb{D}(z_{0},r_{0}) .
\]

根据注 5.10(c)，$ h $ 也具有平均值性质。因此

\begin{equation}
0 = h(z_{0}) = \frac{1}{2\pi}\int_{0}^{2\pi}h(z_{0} + re^{it})\,dt \quad \text{对 } 0 \le r \le r_{0} .\tag{5.8}
\end{equation}

因为对 $ r \in [0,r_{0}] $ 和 $ t \in [0,2\pi] $ 有 $ h(z_{0} + re^{it}) \le 0 $，命题 VI.4.8 和 (5.8) 推出 $ h $ 在 $ \mathbb{D}(z_{0},r_{0}) $ 上恒等于零。因此对 $ z \in \mathbb{D}(z_{0},r_{0}) $ 有 $ \operatorname{Re} f(z) = f(z_{0}) $。现在由 $ |f(z)| \le |f(z_{0})| = \operatorname{Re} f(z_{0}) $ 可得 $ \operatorname{Im}(f(z)) = 0 $，因此对每个 $ z \in \mathbb{D}(z_{0},r_{0}) $，$ f(z) $ 等于 $ f(z_{0}) $。\nobreak\hspace{0.35em}■

\end{subitems}

\end{theorem}

\begin{corollary}

（最大模原理）设 $ U $ 是连通的，且 $ f $ 具有平均值性质。

\begin{subitems}
\item

若 $ |f| $ 在 $ z_{0} \in U $ 处有局部极大值，则 $ f $ 是常值的。

\item

若 $ U $ 是有界的且 $ f \in C(\overline{U},\mathbb{C}) $，则 $ |f| $ 在 $ \partial U $ 上取到其最大值，即存在 $ z_{0} \in \partial U $ 使得 $ |f(z_{0})| = \max_{z\in\overline{U}}|f(z)| $。

\proofhead
(i) 设 $ f(z_{0}) = w_{0} $ 且 $ M := f^{-1}(w_{0}) $。$ f $ 的连续性表明 $ M $ 在 $ U $ 中是闭的（见例 III.2.22(a)）。根据定理 5.17，每个 $ z_{1} \in M $ 都有一个邻域 $ V $ 使得对 $ z \in V $ 有 $ f(z) = f(z_{0}) = w_{0} $。因此 $ M $ 在 $ U $ 中是开的。于是由注 III.4.3 可得 $ M $ 与 $ U $ 重合。因此对每个 $ z \in U $ 有 $ f(z) = w_{0} $。
\end{subitems}

(ii) 因为 $ f $ 在紧致集 $ \overline{U} $ 上连续，我们知道 $ |f| $ 在某个点 $ z_{0} \in \overline{U} $ 处取到其最大值。若 $ z_{0} $ 属于 $ \partial U $，则无须证明。若 $ z_{0} $ 落在 $ U $ 内部，则结论由 (i) 推出。\nobreak\hspace{0.35em}■

\end{corollary}

\needspace{3\baselineskip}\subsection*{调和函数}

设 $ X $ 是 $ \mathbb{R}^{n} $ 中的开集且非空。线性映射

\[
\Delta : C^{2}(X,\mathbb{K}) \to C(X,\mathbb{K}) , \quad f \mapsto \sum_{j=1}^{n}\partial_{j}^{2}f
\]

称为 $ X $ 上的\textbf{拉普拉斯算子}（\textbf{Laplacian}）。若 $ \Delta g = 0 $，则称函数 $ g \in C^{2}(X,\mathbb{K}) $ 在 $ X $ 上\textbf{调和}。我们用 $ \mathcal{H}\mathrm{arm}(X,\mathbb{K}) $ 表示 $ X $ 上全体调和函数构成的集合。

\begin{remark}

\begin{items}
\item

我们有 $ \Delta \in \operatorname{Hom}\left(C^{2}(X,\mathbb{K}),C(X,\mathbb{K})\right) $，且

\[
\mathcal{H}\mathrm{arm}(X,\mathbb{K}) = \Delta^{-1}(0) .
\]

因此调和函数构成 $ C^{2}(X,\mathbb{K}) $ 的向量子空间。

\item

对 $ f \in C^{2}(X,\mathbb{C}) $，我们有

\[
f \in \mathcal{H}\mathrm{arm}(X,\mathbb{C}) \iff \operatorname{Re} f, \operatorname{Im} f \in \mathcal{H}\mathrm{arm}(X,\mathbb{R}) .
\]
\item

$ U $ 中每个全纯函数都是调和的，即 $ C^{\omega}(U,\mathbb{C}) \subset \mathcal{H}\mathrm{arm}(U,\mathbb{C}) $。

\proofhead
若 $ f $ 全纯，则由柯西–黎曼方程得
\[
\partial_{x}^{2}f = \partial_{x}\partial_{y}v - i\,\partial_{x}\partial_{y}u , \quad \partial_{y}^{2}f = -\partial_{y}\partial_{x}v + i\,\partial_{y}\partial_{x}u .
\]

因此由推论 VII.5.5(ii) 得 $ \Delta f = 0 $。\nobreak\hspace{0.35em}■

\item

$ \mathcal{H}\mathrm{arm}(U,\mathbb{C}) \ne C^{\omega}(U,\mathbb{C}) $。

\proofhead
函数 $ U \to \mathbb{C} $，$ x + iy \mapsto x $ 是调和的但不是全纯的（见注 5.4(b)）。\nobreak\hspace{0.35em}■
\end{items}

前面的讨论表明全纯函数的实部是调和的。下一个定理表明，在单连通区域上，每个调和实值函数都是某个全纯函数的实部。

\end{remark}

\begin{proposition}

设 $ u : U \to \mathbb{R} $ 调和。则我们有如下结论：

\begin{subitems}
\item

设 $ V := \mathbb{D}(z_{0},r) \subset U $ 对某个 $ (z_{0},r) \in U\times(0,\infty) $ 成立。则存在 $ V $ 中的全纯函数 $ g $ 使得 $ u = \operatorname{Re} g $。

\item

若 $ U $ 单连通，则存在 $ g \in C^{\omega}(U,\mathbb{C}) $ 使得 $ u = \operatorname{Re} g $。

\proofhead
由于 $ u $ 调和，1-形式 $ \alpha := -u_{y}\,dx + u_{x}\,dy $ 满足可积性条件。因此 $ \alpha $ 是闭的。
\end{subitems}

(i) 由于 $ V $ 单连通，定理 4.8 给出存在 $ v \in C^{1}(V,\mathbb{R}) $ 使得 $ dv = \alpha|V $。因此 $ v_{x} = -u_{y}|V $ 且 $ v_{y} = u_{x}|V $。令 $ g(z) := u(x,y) + i\,v(x,y) $，由注 5.4(b) 与定理 5.11 得 $ g $ 属于 $ C^{\omega}(V,\mathbb{C}) $。

(ii) 此时，我们可以把 (i) 的证明中的圆盘 $ V $ 换成 $ U $。\nobreak\hspace{0.35em}■

\end{proposition}

\begin{corollary}

设 $ u : U \to \mathbb{R} $ 调和。则

\begin{subitems}
\item

$ u \in C^{\infty}(U,\mathbb{R}) $；

\item

$ u $ 具有平均值性质；

\item

若 $ U $ 是区域且存在 $ U $ 的非空子集 $ V $ 使得 $ u|V = 0 $，则 $ u = 0 $。

\proofhead
(i) 与 (ii) 设 $ V := \mathbb{D}(z_{0},r) \subset U $ 对某个 $ r > 0 $ 成立。由于可微性与平均值性质都是局部性质，只需在 $ V $ 上考虑 $ u $。利用命题 5.20，我们找到 $ g \in C^{\omega}(V,\mathbb{C}) $ 使得 $ \operatorname{Re} g = u|V $，而结论由注 5.13(b) 以及注 5.10(b)、(c) 推出。
\end{subitems}

(iii) 设 $ M $ 是 $ U $ 中所有这样的 $ z $ 的集合：存在邻域 $ V $ 使得 $ u|V = 0 $。则 $ M $ 是开集，且由假设非空。设 $ z_{0} \in U $ 是 $ M $ 的聚点。由命题 5.20，存在 $ r > 0 $ 与 $ g \in C^{\omega}\left(\mathbb{D}(z_{0},r),\mathbb{C}\right) $ 使得 $ \operatorname{Re} g = u|\mathbb{D}(z_{0},r) $。进一步，由于 $ z_{0} $ 是 $ M $ 的聚点，$ M\cap\mathbb{D}(z_{0},r) $ 非空。对 $ z_{1} \in M\cap\mathbb{D}(z_{0},r) $，存在 $ z_{1} $ 在 $ U $ 中的邻域 $ V $ 使得 $ u|V = 0 $。因此由注 5.4(d) 得 $ g $ 在 $ V\cap\mathbb{D}(z_{0},r) $ 上为常值。于是解析函数的唯一性定理（定理 V.3.13）表明 $ g $ 在 $ \mathbb{D}(z_{0},r) $ 上为常值。因此 $ u = 0 $ 在 $ \mathbb{D}(z_{0},r) $ 上成立，即 $ z_{0} $ 属于 $ M $。从而 $ M $ 在 $ U $ 中闭，而注 III.4.3 给出 $ M = U $。\nobreak\hspace{0.35em}■

\end{corollary}

\begin{corollary}

（调和函数的极大值极小值原理）设 $ U $ 是区域且 $ u : U \to \mathbb{R} $ 调和。则

\begin{subitems}
\item

若 $ u $ 在 $ U $ 中有局部极值，则 $ u $ 为常值；

\item

若 $ U $ 有界且 $ u \in C(\overline{U},\mathbb{R}) $，则 $ u $ 在 $ \partial U $ 上取到其最大值与最小值。

\proofhead
根据推论 5.21(ii)，$ u $ 具有平均值性质。
\end{subitems}

(i) 设 $ z_{0} \in U $ 是 $ u $ 的局部极值点。当 $ u(z_{0}) $ 是 $ u $ 的正的极大值时，结论由定理 5.17 与推论 5.18、5.21(iii) 推出。若 $ u(z_{0}) $ 是 $ u $ 的正的极小值，则 $ z \mapsto \left[2u(z_{0}) - u(z)\right] $ 在 $ z_{0} $ 处取得正的极大值，结论如第一种情形那样推出。其余情形可类似处理。

(ii) 这由 (i) 推出。\nobreak\hspace{0.35em}■

\end{corollary}

\begin{remark}

\begin{items}
\item

全纯函数的零点集是离散的。\footnote{度量空间 $ X $ 的子集 $ D $ 称为\textbf{离散的}，若每个 $ d \in D $ 在 $ X $ 中有邻域 $ U $ 使得 $ U\cap D = \{d\}$。} 然而，实值调和函数的零点集一般不是离散的。

\proofhead
第一个断言由定理 5.11 与解析函数的唯一性定理（定理 V.3.13）推出。第二个断言通过考虑调和函数 $ \mathbb{C} \to \mathbb{R} $，$ x + iy \mapsto x $ 得到。\nobreak\hspace{0.35em}■
\item

可以证明，函数调和的当且仅当它具有平均值性质（例如见 [Con78, 定理 X.2.11]）。\nobreak\hspace{0.35em}■

\end{items}

\end{remark}

\needspace{3\baselineskip}\subsection*{古尔萨定理}

现在，正如注 5.4(a) 所承诺的那样，我们来证明每个复可微函数都有连续导数，因而是全纯的。作为准备，我们先证明莫雷拉定理，它给出了判定函数全纯的准则。正如我们将看到的，这一准则在其他地方也有用。

设 $ X \subset \mathbb{C} $。$ X $ 中每条由线段组成且恰好有三条边的闭道路 $ \Delta $，若由 $ \Delta $ 围成的三角形的闭包含于 $ X $，则称为 $ X $ 中的\textbf{三角形道路}。

\begin{theorem}

（莫雷拉）设函数 $ f $ 对 $ U $ 中每条三角形道路 $ \Delta $ 满足 $ \int_{\Delta} f\,dz = 0 $。则 $ f $ 是解析的。

\proofhead
设 $ a \in U $ 且 $ r > 0 $ 满足 $ \mathbb{D}(a,r) \subset U $。只需验证 $ f|\mathbb{D}(a,r) $ 是解析的。故设 $ z_{0} \in \mathbb{D}(a,r) $ 且
\[
F : \mathbb{D}(a,r) \to \mathbb{C} , \quad z \mapsto \int_{[[a,z]]}f(w)\,dw .
\]

我们的假设给出恒等式

\[
F(z) = \int_{[[a,z_{0}]]}f(w)\,dw + \int_{[[z_{0},z]]}f(w)\,dw = F(z_{0}) + \int_{[[z_{0},z]]}f(w)\,dw .
\]

因此

\begin{equation}
\frac{F(z) - F(z_{0})}{z - z_{0}} - f(z_{0}) = \frac{1}{z - z_{0}}\int_{[[z_{0},z]]}\left(f(w) - f(z_{0})\right)dw \quad \text{对 } z \ne z_{0} .\tag{5.9}
\end{equation}

现在设 $ \varepsilon > 0 $。则存在 $ \delta \in (0, r - |z_{0} - a|) $ 使得对 $ w \in \mathbb{D}(z_{0},\delta) $ 有 $ |f(w) - f(z_{0})| < \varepsilon $。于是由 (5.9) 可得

\[
\left|\frac{F(z) - F(z_{0})}{z - z_{0}} - f(z_{0})\right| < \varepsilon \quad \text{对 } 0 < |z - z_{0}| < \delta .
\]

因此 $ F'(z_{0}) = f(z_{0}) $，这表明 $ F $ 有连续导数。于是定理 5.11 表明 $ F $ 属于 $ C^{\omega}\left(\mathbb{D}(a,r),\mathbb{C}\right) $，从而我们得到 $ F' = f|\mathbb{D}(a,r) $ 属于 $ C^{\omega}\left(\mathbb{D}(a,r),\mathbb{C}\right) $。

\end{theorem}

\begin{theorem}

（古尔萨）设 $ f $ 可微。则对 $ U $ 中每条三角形道路 $ \Delta $ 有 $ \int_{\Delta} f\,dz = 0 $。

\proofhead
(i) 设 $ \Delta $ 是 $ U $ 中的三角形道路。不失一般性，我们可以假设 $ \Delta $ 围成一个面积为正的三角形。则 $ \Delta $ 有顶点 $ z_{0} $、$ z_{1} $ 与 $ z_{2} $，且 $ z_{2} \in [[z_{0},z_{1}]] $。于是 $ \int_{\Delta} f(z)\,dz = 0 $，这很容易验证。我们用 $ K $ 表示由 $ \Delta $ 围成的三角形的闭包。
\bookfigure[0.41]{figures/b7ecf59647bd9ee7.png}

通过连接 $ \Delta $ 三边的中点，我们得到 $ K $ 的四个全等闭子三角形 $ K_{1},\ldots,K_{4} $。我们对这些 $ K_{j} $ 的（拓扑）边界定向，使得公共边方向相反，并把所得的三角形道路记为 $ \Delta_{1},\ldots,\Delta_{4} $。则

\[
\left|\int_{\Delta} f(z)\,dz\right| = \left|\sum_{j=1}^{4}\int_{\Delta_{j}}f(z)\,dz\right| \le 4\max_{1\le j\le 4}\left|\int_{\Delta_{j}}f(z)\,dz\right| .
\]

在四条三角形道路 $ \Delta_{1},\ldots,\Delta_{4} $ 中，有一条——我们记作 $ \Delta^{1} $——满足

\[
\left|\int_{\Delta^{1}}f(z)\,dz\right| = \max_{1\le j\le 4}\left|\int_{\Delta_{j}}f(z)\,dz\right| .
\]

则

\[
\left|\int_{\Delta} f(z)\,dz\right| \le 4\left|\int_{\Delta^{1}}f(z)\,dz\right| .
\]

\begin{subitems}
\item[（ii）]

对 $ \Delta^{1} $，我们应用与对 $ \Delta $ 相同的分解与选取程序。于是我们归纳地得到三角形道路的序列 $ (\Delta^{n}) $ 以及相应的闭三角形 $ K^{n} $，使得

\begin{equation}
K^{1} \supset K^{2} \supset \cdots \supset K^{n} \supset \cdots \quad \text{且} \quad \left|\int_{\Delta^{n}}f(z)\,dz\right| \le 4\left|\int_{\Delta^{n+1}}f(z)\,dz\right|\tag{5.10}
\end{equation}

对 $ n \in \mathbb{N}^{\times} $ 成立。显然 $ \{K^{n} ; n \in \mathbb{N}^{\times}\} $ 具有有限交性质。因此由 $ K^{1} $ 的紧致性得 $ \bigcap_{n}K^{n} $ 非空（见习题 III.3.5）。我们在 $ \bigcap_{n}K^{n} $ 中取定一个 $ z_{0} $。

\item[（iii）]

(5.10) 中的不等式给出

\begin{equation}
\left|\int_{\Delta} f(z)\,dz\right| \le 4^{n}\left|\int_{\Delta^{n}}f(z)\,dz\right| .\tag{5.11}
\end{equation}

此外，我们有初等几何关系

\[
L(\Delta^{n+1}) = L(\Delta^{n})/2 \quad \text{且} \quad \operatorname{diam}(K^{n+1}) = \operatorname{diam}(K^{n})/2 \quad \text{对 } n \in \mathbb{N}^{\times} .
\]

由此可得

\[
L(\Delta^{n}) = \ell/2^{n} \quad \text{且} \quad \operatorname{diam}(K^{n}) = d/2^{n} \quad \text{对 } n \in \mathbb{N}^{\times} ,
\]

其中 $ \ell := L(\Delta) $ 且 $ d := \operatorname{diam}(K) $。

\item[（iv）]

设 $ \varepsilon > 0 $。$ f $ 在 $ z_{0} $ 处的可微性给出存在 $ \delta > 0 $ 使得 $ \mathbb{D}(z_{0},\delta) \subset U $ 且

\[
\left|f(z) - f(z_{0}) - f'(z_{0})(z - z_{0})\right| \le \frac{\varepsilon}{d\ell}|z - z_{0}| \quad \text{对 } z \in \mathbb{D}(z_{0},\delta) .
\]

现在我们选取 $ n \in \mathbb{N}^{\times} $ 使得 $ \operatorname{diam}(K^{n}) = d/2^{n} < \delta $。由于 $ z_{0} \in K^{n} $，于是有 $ \Delta^{n} \subset \mathbb{D}(z_{0},\delta) $。柯西积分定理给出

\[
0 = \int_{\Delta^{n}}dz = \int_{\Delta^{n}}z\,dz .
\]

从而

\begin{equation*}
\begin{aligned}
\left|\int_{\Delta^{n}}f(z)\,dz\right| &= \left|\int_{\Delta^{n}}\left(f(z) - f(z_{0}) - f'(z_{0})(z - z_{0})\right)dz\right| \\ &\le \frac{\varepsilon}{d\ell}\max_{z\in\Delta^{n}}|z - z_{0}|\,L(\Delta^{n}) \le \frac{\varepsilon}{d\ell}\operatorname{diam}(K^{n})L(\Delta^{n}) = \frac{\varepsilon}{4^{n}} .
\end{aligned}
\end{equation*}

现在 (5.11) 完成证明。\nobreak\hspace{0.35em}■

\end{subitems}

\end{theorem}

\begin{corollary}

若 $ f $ 可微，则 $ f $ 是全纯的。

\proofhead
这由定理 5.24 与 5.25 推出。\nobreak\hspace{0.35em}■
\end{corollary}

\needspace{3\baselineskip}\subsection*{魏尔斯特拉斯收敛定理}

作为莫雷拉定理的进一步应用，我们证明魏尔斯特拉斯关于全纯函数局部一致收敛序列的极限的一个定理。我们已经——结合定理 5.11——在应用 VI.7.23(b) 中运用了这一结果来证明正弦的乘积表示。

\begin{theorem}

（魏尔斯特拉斯收敛定理）设 $ (g_{n}) $ 是 $ U $ 中全纯函数的局部一致收敛序列。则 $ g := \lim g_{n} $ 在 $ U $ 中全纯。

\proofhead
根据定理 V.2.1，$ g $ 连续。由注 V.2.3(c) 我们知道 $ (g_{n}) $ 在 $ U $ 的每个紧致子集上一致收敛。因此 $ (g_{n}) $ 在 $ U $ 中每条三角形道路 $ \Delta $ 上一致收敛于 $ g $。因此由命题 VI.4.1(i) 可得
\begin{equation}
\int_{\Delta} g\,dz = \lim_{n\to\infty}\int_{\Delta} g_{n}\,dz = 0 ,\tag{5.12}
\end{equation}

其中最后一个等号由应用于 $ g_{n} $ 的柯西积分定理得到。由于 (5.12) 对 $ U $ 中每条三角形道路成立，莫雷拉定理完成证明。\nobreak\hspace{0.35em}■

\end{theorem}

\exerciseshead

\begin{exercises}
\item

设 $ f : U \to \mathbb{C} $ 实可微，

\[
\partial_W f := \frac{1}{2}(\partial_x f - i\partial_y f) \quad \text{与} \quad \overline{\partial}_W f := \frac{1}{2}(\partial_x f + i\partial_y f)
\]

称为 $ f $ 的 \textbf{维尔廷格导数}\footnote{通常，维尔廷格导数将记作 $ \partial $（或 $ \overline{\partial} $）。然而，当这种记号可能与我们表示导数算子的记号 $ \partial $ 相冲突时，我们将写 $ \partial_W $（或相应地 $ \overline{\partial}_W $）。}。

证明

\begin{items}
\item

$ \overline{\partial}_W f = \overline{\partial_W \overline{f}} $，且 $ \overline{\partial}_W \overline{f} = \overline{\partial_W f} $；

\item

$ f $ 全纯 $ \iff \overline{\partial}_W f = 0 $；

\item

$ 4\partial_W \overline{\partial}_W f = 4\overline{\partial}_W \partial_W f = \Delta f $，若 $ f $ 二次实可微；

\item

$ \det\begin{bmatrix} u_x & u_y \\ v_x & v_y \end{bmatrix} = \det\begin{bmatrix} \partial_W f & \overline{\partial}_W f \\ \partial_W \overline{f} & \overline{\partial}_W \overline{f} \end{bmatrix} = \left|\partial_W f\right|^2 - \left|\overline{\partial}_W f\right|^2 $。

\end{items}

\item

设 $ d\overline{z} := d(z \mapsto \overline{z}) $。则 $ d\overline{z} = dx - i\,dy $，且

\[
df = \partial_x f\,dx + \partial_y f\,dy = f'\,dz = \partial_W f\,dz + \overline{\partial}_W f\,d\overline{z} \quad \text{对 } f \in C^1(U,\mathbb{C}) .
\]
\item

1-形式 $ \alpha \in \Omega_{(1)}(U,\mathbb{C}) $ 称为\textbf{全纯的}，若每个 $ z_0 \in U $ 都有一个邻域 $ V $，使得存在 $ f \in C^1(V,\mathbb{C}) $ 满足 $ df = \alpha|V $。

\begin{items}
\item

证明下列陈述等价：

\begin{subitems}
\item

$ \alpha $ 是全纯的。

\item

存在实可微函数 $ a \in C(U,\mathbb{C}) $，使得 $ \alpha = a\,dz $ 且 $ \alpha $ 是闭的。

\item

存在 $ a \in C^1(U,\mathbb{C}) $，使得 $ \alpha = a\,dz $。

\end{subitems}

\item

证明

\[
\alpha := \frac{x\,dx + y\,dy}{x^2 + y^2} + i\,\frac{x\,dy - y\,dx}{x^2 + y^2}
\]

是 $ \mathbb{C}^{\times} $ 中的全纯 1-形式。$ \alpha $ 是否全局恰当，即是否存在 $ \mathbb{C}^{\times} $ 中的全纯函数 $ f $，使得 $ \alpha = f\,dz $？

\end{items}

\item

设 $ U $ 连通且 $ u : U \to \mathbb{R} $ 调和。若 $ v \in C^1(U,\mathbb{R}) $ 满足 $ v_x = -u_y $ 与 $ v_y = u_x $，我们称 $ v $ 与 $ u $ \textbf{共轭}。证明

\begin{items}
\item

若 $ v $ 与 $ u $ 共轭，则 $ v $ 调和；

\item

若 $ U $ 单连通，则对 $ U $ 中每个调和函数都存在一个调和共轭。（提示：考虑 $ u_x - i\,u_y $。）

\end{items}

\item

用柯西积分公式证明魏尔斯特拉斯收敛定理。并证明，在定理 5.27 的假设下，对每个 $ k \in \mathbb{N} $，$ k $ 阶导数的序列 $ (f_n^{(k)})_{n \in \mathbb{N}} $ 在 $ U $ 上局部一致收敛到 $ f^{(k)} $。（提示：命题 VII.6.7。）

\item

设 $ z_0 \in U $ 且 $ r > 0 $ 满足 $ \overline{\mathbb{D}}(z_0,r) \subset U $。进一步设 $ z \in \mathbb{D}(z_0,r) $ 且 $ \delta > 0 $ 满足 $ \mathbb{D}(z,\delta) \subset U $。验证 $ \partial\mathbb{D}(z,\delta) $ 与 $ \partial\mathbb{D}(z_0,r) $ 在 $ U \setminus \{z\} $ 中同伦。

\item

借助柯西积分定理，证明 $ \int_{-\infty}^{\infty} dx/(1 + x^2) = \pi $。

\item

设 $ p = \sum_{k=0}^{n} a_k X^k \in \mathbb{C}[X] $ 满足 $ a_n = 1 $。令 $ R := \max\{1, 2\sum_{k=0}^{n-1}\left|a_k\right|\} $。证明

\[
|z|^n/2 \le |p(z)| \le 2|z|^n \quad \text{对 } z \in R\mathbb{D}^{c} .
\]
\item

设 $ 0 \le r_0 < r < r_1 \le \infty $ 且 $ f $ 在 $ \mathbb{D}(z_0,r_1) $ 中全纯。证明

\[
\left|f^{(n)}(z)\right| \le \frac{r\,n!}{(r - r_0)^{n+1}} \max_{w \in \partial\mathbb{D}(z_0,r)} |f(w)| \quad \text{对 } z \in \mathbb{D}(z_0,r_0) \text{ 且 } n \in \mathbb{N} .
\]
\item

给定整函数 $ f $，设存在 $ M,R > 0 $ 与 $ n \in \mathbb{N} $，使得 $ |f(z)| \le M|z|^n $ 对 $ z \in R\mathbb{D}^{c} $。证明 $ f $ 是满足 $ \deg(f) \le n $ 的多项式。（提示：习题 9。）

\item

设 $ \Gamma_1 $ 与 $ \Gamma_2 $ 是 $ U $ 中的紧致分段 $ C^1 $ 曲线，且 $ f \in C(U \times U,\mathbb{C}) $。证明

\[
\int_{\Gamma_2}\left(\int_{\Gamma_1} f(w,z)\,dw\right)dz = \int_{\Gamma_1}\left(\int_{\Gamma_2} f(w,z)\,dz\right)dw .
\]

（提示：习题 VII.6.2。）

\item

设 $ \Gamma $ 是 $ U $ 中的紧致分段 $ C^1 $ 曲线，且 $ f \in C(U \times U,\mathbb{C}) $。又设 $ f(w,\cdot) $ 对每个 $ w \in U $ 全纯。证明

\[
F : U \to \mathbb{C} , \quad z \mapsto \int_{\Gamma} f(w,z)\,dw
\]

全纯，且 $ F' = \int_{\Gamma} \partial_2 f(w,\cdot)\,dw $。

（提示：习题 11、莫雷拉定理与命题 VII.6.7。）

\item

设 $ L = a + \mathbb{R}b $，$ a,b \in \mathbb{C} $ 是 $ \mathbb{C} $ 中的直线，且 $ f \in C(U,\mathbb{C}) $ 在 $ U \setminus L $ 中全纯。证明 $ f $ 在整个 $ U $ 中全纯。（提示：莫雷拉定理。）

\item

设 $ U $ 连通且 $ f $ 在 $ U $ 中全纯。证明下列命题。

\begin{items}
\item

若 $ |f| $ 在 $ z_0 $ 处取得全局极小值，则要么 $ f $ 是常值，要么 $ f(z_0) = 0 $。

\item

设 $ U $ 有界且 $ f \in C(\overline{U},\mathbb{C}) $。证明要么 $ f $ 在 $ U $ 中有零点，要么 $ |f| $ 在 $ \partial U $ 上取得其极小值。

\end{items}

\item

对 $ R > 0 $，

\[
P_R : R\partial\mathbb{D} \times R\mathbb{D} \to \mathbb{C} , \quad (\zeta,z) \mapsto \frac{R^2 - |z|^2}{|\zeta - z|^2}
\]

称为 $ R\mathbb{D} $ 的\textbf{泊松核}。

证明

\begin{items}
\item

$ P_R(\zeta,z) = \mathrm{Re}\left((\zeta + z)/(\zeta - z)\right) $ 对 $ (\zeta,z) \in R\partial\mathbb{D} \times R\mathbb{D} $；

\item

对每个 $ \zeta \in R\partial\mathbb{D} $，函数 $ P_R(\zeta,\cdot) $ 在 $ R\mathbb{D} $ 中调和；

\item

对 $ r \in [0,R] $ 与 $ t,\theta \in [0,2\pi) $，

\[
P_R(Re^{i\theta}, re^{it}) = \frac{R^2 - r^2}{R^2 - 2Rr\cos(\theta - t) + r^2} ;
\]
\item

$ P_1(1, re^{it}) = \sum_{n=-\infty}^{\infty} r^{|n|}e^{int} $ 对 $ r \in [0,1) $ 与 $ t \in \mathbb{R} $；

\item

$ \int_0^{2\pi} P_R(Re^{i\theta}, z)\,d\theta = 2\pi $ 对 $ z \in R\mathbb{D} $。
（(d) 的提示：$ \sum_{k=0}^{\infty}(re^{it})^k = 1/(1 - re^{it}) $。）

\end{items}

\item

设 $ \rho > 1 $ 且 $ f $ 在 $ \rho\mathbb{D} $ 中全纯。证明

\[
f(z) = \frac{1}{2\pi}\int_0^{2\pi} P_1(e^{i\theta}, z)f(e^{i\theta})\,d\theta \quad \text{对 } z \in \mathbb{D} .
\]

（提示：

\begin{subitems}
\item

对 $ g \in C^1(\rho_0\mathbb{D},\mathbb{C}) $ 且 $ \rho_0 > 1 $，证明

\[
g(z) = \frac{1}{2\pi}\int_0^{2\pi} \frac{g(e^{i\theta})}{1 - e^{-i\theta}z}\,d\theta \quad \text{对 } z \in \mathbb{D} .
\]
\item

对 $ z \in \mathbb{D} $，$ \rho_0 := \min(\rho, 1/|z|) $，$ g : \rho_0\mathbb{D} \to \mathbb{C} $，证明 $ w \mapsto f(w)/(1 - w\overline{z}) $ 全纯。）

\end{subitems}

\item

证明下列命题：

\begin{items}
\item

对 $ g \in C(\partial\mathbb{D},\mathbb{R}) $，

\[
\mathbb{D} \to \mathbb{C} , \quad z \mapsto \int_0^{2\pi} P_1(e^{i\theta}, z)g(e^{i\theta})\,d\theta
\]

调和。

\item

若 $ f \in C(\overline{\mathbb{D}},\mathbb{R}) $ 在 $ \mathbb{D} $ 中调和，则

\[
f(z) = \frac{1}{2\pi}\int_0^{2\pi} P_1(e^{i\theta}, z)f(e^{i\theta})\,d\theta \quad \text{对 } z \in \mathbb{D} .
\]

（提示：(a) 习题 15(b) 与命题 VII.6.7。(b) 设 $ 0 < r_k < 1 $ 满足 $ \lim r_k = 1 $，且 $ f_k(z) := f(r_k z) $ 对 $ z \in r_k^{-1}\mathbb{D} $。习题 16 给出

\[
f_k(z) = \frac{1}{2\pi}\int_0^{2\pi} P_1(e^{i\theta}, z)f_k(e^{i\theta})\,dz \quad \text{对 } z \in \mathbb{D} .
\]

现在考虑极限 $ k \to \infty $。）

\end{items}

\item

设 $ a \in \mathbb{C} $ 且 $ \alpha \ne \mathrm{Re}\,a $。证明对 $ \gamma_\alpha : \mathbb{R} \to \mathbb{C} $，$ s \mapsto \alpha + is $，我们有

\[
e^{ta} = \frac{1}{2\pi i}\int_{\gamma_\alpha} e^{\lambda t}(\lambda - a)^{-1}\,d\lambda \quad \text{对 } t > 0 .
\]

（提示：柯西积分公式给出

\[
e^{ta} = \frac{1}{2\pi i}\int_{\partial\mathbb{D}(a,r)} e^{\lambda t}(\lambda - a)^{-1}\,d\lambda \quad \text{对 } t \in \mathbb{R} \text{ 且 } r > 0 .
\]

现在应用柯西积分定理。）

\end{exercises}

\section{亚纯函数}

在本章的核心各节中，我们研究了除孤立点外全纯的复函数。典型的例子是

\[
\mathbb{C}^{\times} \to \mathbb{C} , \quad z \mapsto e^{1/z} , \quad \mathbb{C}\setminus\{\pm i\} \to \mathbb{C} , \quad z \mapsto 1/(1 + z^{2}) ,
\]
\[
\mathbb{C}^{\times} \to \mathbb{C} , \quad z \mapsto \sin(z)/z , \quad \mathbb{C}\setminus\pi\mathbb{Z} \to \mathbb{C} , \quad z \mapsto \cot z .
\]

结果表明，这些“例外点”导致一个出奇简单的分类。这一分类借助洛朗级数来完成，洛朗级数把一个函数展开成既含正次幂又含负次幂的幂级数。因此，它把泰勒级数推广为能够展开全纯函数的级数。

在此语境下，我们还将推广柯西积分定理，并证明留数定理；留数定理有许多重要的应用，这里我们只给出其中少数几个。

\needspace{3\baselineskip}\subsection*{洛朗展开}

对 $ c := (c_{n}) \in \mathbb{C}^{\mathbb{Z}} $，我们考虑幂级数

\[
nc := \sum_{n\ge 0} c_{n}X^{n} \quad\text{且}\quad hc := \sum_{n\ge 1} c_{-n}X^{n} .
\]

设它们的收敛半径分别为 $ \rho_{1} $ 与 $ 1/\rho_{0} $，且 $ 0 \le \rho_{0} < \rho_{1} \le \infty $。则由定理 V.3.1，$ nc $ [或 $ hc $] 在 $ \rho_{1}\mathbb{D} $ [或 $ (1/\rho_{0})\mathbb{D} $] 中所表示的函数 $ \underline{nc} $ [或 $ \underline{hc} $] 是全纯的。因为 $ z \mapsto 1/z $ 在 $ \mathbb{C}^{\times} $ 中全纯，且当 $ |z| > \rho_{0} $ 时有 $ |1/z| < 1/\rho_{0} $，注 5.13(d) 保证了函数

\[
z \mapsto \underline{hc}(1/z) = \sum_{n=1}^{\infty} c_{-n}z^{-n}
\]

在 $ |z| > \rho_{0} $ 上全纯。因此

\[
z \mapsto \sum_{n=-\infty}^{\infty} c_{n}z^{n} := \sum_{n=0}^{\infty} c_{n}z^{n} + \sum_{n=1}^{\infty} c_{-n}z^{-n}
\]

是由

\[
\Omega(\rho_{0},\rho_{1}) := \rho_{1}\mathbb{D}\setminus\rho_{0}\overline{\mathbb{D}} = \{ z \in \mathbb{C} ;\ \rho_{0} < |z| < \rho_{1} \} .
\]

给出的以 0 为中心的圆环域中的全纯函数。

现在设 $ z_{0} \in \mathbb{C} $。则

\begin{equation}
\sum_{n\in\mathbb{Z}} c_{n}(z - z_{0})^{n} := \sum_{n\ge 0} c_{n}(z - z_{0})^{n} + \sum_{n\ge 1} c_{-n}(z - z_{0})^{-n}\tag{6.1}
\end{equation}

称为关于\textbf{展开点} $ z_{0} $ 的\textbf{洛朗级数}，而 $ (c_{n}) $ 是系数序列。在 (6.1) 中，$ \sum_{n\ge 1} c_{-n}(z - z_{0})^{-n} $ 是\textbf{主部}，而 $ \sum_{n\ge 0} c_{n}(z - z_{0})^{n} $ 是\textbf{正则部}。若主部与正则部都在 $ M \subset \mathbb{C} $ 中收敛 [或依范数收敛]，则称洛朗级数 (6.1) 在 $ M $ 中\textbf{收敛} [或\textbf{依范数收敛}]。按定义，它在 $ z \in M $ 处的\textbf{值}是主部与正则部在 $ z $ 处的值之和，即

\[
\sum_{n=-\infty}^{\infty} c_{n}(z - z_{0})^{n} := \sum_{n=0}^{\infty} c_{n}(z - z_{0})^{n} + \sum_{n=1}^{\infty} c_{-n}(z - z_{0})^{-n} .
\]

由这些预备知识及定理 V.1.8 可知，洛朗级数 $ \sum_{n\in\mathbb{Z}} c_{n}(z - z_{0})^{n} $ 在由

\[
z_{0} + \Omega(\rho_{0},\rho_{1}) = \{ z \in \mathbb{C} ;\ \rho_{0} < |z - z_{0}| < \rho_{1} \}
\]

给出的以 $ z_{0} $ 为中心的圆环域的每个紧致子集中正规收敛，且函数

\[
z_{0} + \Omega(\rho_{0},\rho_{1}) \to \mathbb{C} , \quad z \mapsto \sum_{n=-\infty}^{\infty} c_{n}(z - z_{0})^{n}
\]

是全纯的。上面的讨论表明，反之，在圆环域中全纯的每个函数都可以由一个洛朗级数来表示。

\begin{lemma}

对 $ \rho > 0 $ 与 $ a \in \mathbb{C} $，

\begin{equation*}
\int_{\rho\partial\mathbb{D}} \frac{dz}{z - a} = \begin{cases} 2\pi i , & |a| < \rho , \\ 0 , & |a| > \rho . \end{cases}
\end{equation*}

证明

\begin{subitems}
\item

设 $ |a| < \rho $，且 $ \delta > 0 $ 满足 $ \mathbb{D}(a,\delta) \subset \rho\mathbb{D} $。因为 $ \partial\mathbb{D}(a,\delta) $ 与 $ \rho\partial\mathbb{D} $ 在 $ \mathbb{C}\setminus\{a\} $ 中同伦（见习题 5.6），该断言由命题 5.7 与例 5.3(a) 推出。

\item

若 $ |a| > \rho $，则 $ \rho\partial\mathbb{D} $ 在 $ \mathbb{C}\setminus\{a\} $ 中零伦，且该断言同样由命题 5.7 推出。\nobreak\hspace{0.35em}■

\end{subitems}

\end{lemma}

\begin{lemma}

设 $ f : \Omega(r_{0},r_{1}) \to \mathbb{C} $ 全纯。

\begin{subitems}
\item

对 $ r,s \in (r_{0},r_{1}) $，我们有

\[
\int_{r\partial\mathbb{D}} f(z)\,dz = \int_{s\partial\mathbb{D}} f(z)\,dz .
\]
\item

设 $ a \in \Omega(\rho_{0},\rho_{1}) $，且 $ r_{0} < \rho_{0} < \rho_{1} < r_{1} $。则

\[
f(a) = \frac{1}{2\pi i}\int_{\rho_{1}\partial\mathbb{D}}\frac{f(z)}{z - a}\,dz - \frac{1}{2\pi i}\int_{\rho_{0}\partial\mathbb{D}}\frac{f(z)}{z - a}\,dz .
\]

\proofhead
(i) 因为 $ r\partial\mathbb{D} $ 与 $ s\partial\mathbb{D} $ 在 $ \Omega := \Omega(r_{0},r_{1}) $ 中同伦，该断言由命题 5.7 推出。
\end{subitems}

(ii) 设 $ g : \Omega \to \mathbb{C} $ 由

\begin{equation*}
g(z) := \begin{cases} \left(f(z) - f(a)\right)/(z - a) , & z \in \Omega\setminus\{a\} , \\ f'(a) , & z = a . \end{cases}
\end{equation*}

定义。显然 $ g $ 在 $ \Omega\setminus\{a\} $ 中全纯，且

\begin{equation}
g'(z) = \frac{f'(z)(z - a) - f(z) + f(a)}{(z - a)^{2}} \quad\text{对 } z \in \Omega\setminus\{a\} .\tag{6.2}
\end{equation}

由泰勒公式（推论 IV.3.3）

\begin{equation}
f(z) = f(a) + f'(a)(z - a) + \frac{1}{2}f''(a)(z - a)^{2} + o(|z - a|^{2})\tag{6.3}
\end{equation}

（当 $ z \to a $），我们得到

\[
\frac{g(z) - g(a)}{z - a} = \frac{1}{z - a}\left(\frac{f(z) - f(a)}{z - a} - f'(a)\right) = \frac{1}{2}f''(a) + \frac{o(|z - a|^{2})}{(z - a)^{2}}
\]

（当 $ z $ 在 $ \Omega\setminus\{a\} $ 中趋于 $ a $ 时）。于是 $ g $ 在 $ a $ 处可微，且 $ g'(a) = f''(a)/2 $。因此 $ g $ 在 $ \Omega $ 中全纯。将引理 6.1 与 (i) 应用于 $ g $ 即得该断言。\nobreak\hspace{0.35em}■

\end{lemma}

现在我们可以证明前面提到的展开定理。

\begin{theorem}

（洛朗展开）在 $ \Omega := \Omega(r_{0},r_{1}) $ 中全纯的每个函数 $ f $ 都有唯一的洛朗展开

\begin{equation}
f(z) = \sum_{n=-\infty}^{\infty} c_{n}z^{n} \quad\text{对 } z \in \Omega .\tag{6.4}
\end{equation}

洛朗级数在 $ \Omega $ 的每个紧致子集上正规收敛，且其系数由

\begin{equation}
c_{n} = \frac{1}{2\pi i}\int_{r\partial\mathbb{D}}\frac{f(\zeta)}{\zeta^{n+1}}\,d\zeta \quad\text{对 } n \in \mathbb{Z} \text{ 且 } r_{0} < r < r_{1} .\tag{6.5}
\end{equation}

给出。

证明

\begin{subitems}
\item

我们首先验证 $ f $ 由系数按 (6.5) 给出的洛朗级数表示。由引理 6.2(i) 可知 $ c_{n} $ 对 $ n \in \mathbb{Z} $ 良定义，即它与 $ r $ 无关。
设 $ r_{0} < s_{0} < s_{1} < r_{1} $ 且 $ z \in \Omega(s_{0},s_{1}) $。对满足 $ |\zeta| = s_{1} $ 的 $ \zeta \in \mathbb{C} $，我们有 $ |z/\zeta| < 1 $，从而

\[
\frac{1}{\zeta - z} = \frac{1}{\zeta}\cdot\frac{1}{1 - z/\zeta} = \sum_{n=0}^{\infty}\frac{z^{n}}{\zeta^{n+1}} ,
\]

在 $ s_{1}\partial\mathbb{D} $ 上正规收敛。因此我们有

\[
\frac{1}{2\pi i}\int_{s_{1}\partial\mathbb{D}}\frac{f(\zeta)}{\zeta - z}\,d\zeta = \sum_{n=0}^{\infty} c_{n}z^{n}
\]

（见习题 4.9）。对满足 $ |\zeta| = s_{0} $ 的 $ \zeta \in \mathbb{C} $，我们有 $ |\zeta/z| < 1 $，从而

\[
\frac{1}{\zeta - z} = -\frac{1}{z}\cdot\frac{1}{1 - \zeta/z} = -\sum_{m=0}^{\infty}\frac{\zeta^{m}}{z^{m+1}} ,
\]

在 $ s_{0}\partial\mathbb{D} $ 上正规收敛。因此

\begin{equation*}
\begin{aligned}
\frac{1}{2\pi i}\int_{s_{0}\partial\mathbb{D}}\frac{f(\zeta)}{\zeta - z}\,d\zeta &= -\sum_{m=0}^{\infty}\left(\frac{1}{2\pi i}\int_{s_{0}\partial\mathbb{D}} f(\zeta)\zeta^{m}\,d\zeta\right)z^{-m-1} \\ &= -\sum_{n=1}^{\infty} c_{-n}z^{-n} .
\end{aligned}
\end{equation*}

于是由引理 6.2(ii) 我们得到表示 (6.4)。

\item

现在我们证明 (6.4) 中的系数是唯一的。设 $ f(z) = \sum_{n=-\infty}^{\infty} a_{n}z^{n} $ 在 $ \Omega $ 的紧致子集上正规收敛。对 $ r \in (r_{0},r_{1}) $ 与 $ m \in \mathbb{Z} $，我们有（见例 5.3(a)）

\[
\frac{1}{2\pi i}\int_{r\partial\mathbb{D}} f(z)z^{-m-1}\,dz = \frac{1}{2\pi i}\sum_{n=-\infty}^{\infty} a_{n}\int_{r\partial\mathbb{D}} z^{n-m-1}\,dz = a_{m} .
\]

这表明对 $ m \in \mathbb{Z} $ 有 $ a_{m} = c_{m} $。

\item

因为我们已知洛朗级数在 $ \Omega $ 的紧致子集上正规收敛，定理得证。\nobreak\hspace{0.35em}■

\end{subitems}

\end{theorem}

作为该定理的一个简单推论，我们得到在开\textbf{穿孔圆盘}

\[
\mathbb{D}^{\bullet}(z_{0},r) := z_{0} + \Omega(0,r) = \{ z \in \mathbb{C} ;\ 0 < |z - z_{0}| < r \} .
\]

中全纯的函数的洛朗展开。

\begin{corollary}

设 $ f $ 在 $ \mathbb{D}^{\bullet}(z_{0},r) $ 中全纯。则 $ f $ 有唯一的洛朗展开

\[
f(z) = \sum_{n=-\infty}^{\infty} c_{n}(z - z_{0})^{n} \quad\text{对 } z \in \mathbb{D}^{\bullet}(z_{0},r) ,
\]

其中

\[
c_{n} := \frac{1}{2\pi i}\int_{\partial\mathbb{D}(z_{0},\rho)}\frac{f(z)}{(z - z_{0})^{n+1}}\,dz \quad\text{对 } n \in \mathbb{Z} \text{ 且 } \rho \in (0,r) .
\]

该级数在 $ \mathbb{D}^{\bullet}(z_{0},r) $ 的每个紧致子集上正规收敛，且

\begin{equation}
|c_{n}| \le \rho^{-n}\max_{z\in\partial\mathbb{D}(z_{0},\rho)}|f(z)| \quad\text{对 } n \in \mathbb{Z} \text{ 且 } \rho \in (0,r) .\tag{6.6}
\end{equation}

\proofhead
除 (6.6) 外，所有断言都由定理 6.3 应用于 $ z \mapsto f(z + z_{0}) $ 得到。估计 (6.6) 由 (6.5) 与命题 5.2(ii) 推出。\nobreak\hspace{0.35em}■
\end{corollary}

\needspace{3\baselineskip}\subsection*{可去奇点}

在下文中设

• $ U $ 是 $ \mathbb{C} $ 的开子集，且 $ z_{0} $ 是 $ U $ 中的一个点。

给定全纯函数 $ f : U\setminus\{z_{0}\} \to \mathbb{C} $，若 $ f $ 有全纯延拓 $ F : U \to \mathbb{C} $，则点 $ z_{0} $ 是\textbf{可去奇点}。当不致混淆时，我们对这个延拓沿用符号 $ f $。

\begin{example}

设 $ f : U \to \mathbb{C} $ 全纯。则 $ z_{0} $ 是

\[
g : U\setminus\{z_{0}\} \to \mathbb{C} , \quad z \mapsto \left(f(z) - f(z_{0})\right)/(z - z_{0})
\]

的可去奇点。特别地，0 是

\[
z \mapsto \sin(z)/z , \quad z \mapsto (\cos z - 1)/z , \quad z \mapsto \log(z + 1)/z
\]

的可去奇点。

\proofhead
这由引理 6.2(ii) 的证明推出。\nobreak\hspace{0.35em}■

函数 $ f $ 的可去奇点可以由 $ f $ 的局部有界性来刻画：
\end{example}

\begin{theorem}

（黎曼可去奇点定理）设 $ f : U\setminus\{z_{0}\} \to \mathbb{C} $ 全纯。则点 $ z_{0} $ 是 $ f $ 的可去奇点，当且仅当 $ f $ 在 $ z_{0} $ 的一个邻域中有界。

\proofhead
设 $ r > 0 $ 满足 $ \overline{\mathbb{D}}(z_{0},r) \subset U $。

若 $ z_{0} $ 是 $ f $ 的可去奇点，则存在 $ F \in C^{\omega}(U) $ 使得 $ F \supset f $。由 $ \overline{\mathbb{D}}(z_{0},r) $ 的紧致性可得
\[
\sup_{z\in\mathbb{D}^{\bullet}(z_{0},r)}|f(z)| = \sup_{z\in\mathbb{D}^{\bullet}(z_{0},r)}|F(z)| \le \max_{z\in\overline{\mathbb{D}}(z_{0},r)}|F(z)| < \infty .
\]

因此 $ f $ 在 $ U\setminus\{z_{0}\} $ 中 $ z_{0} $ 的邻域 $ \mathbb{D}^{\bullet}(z_{0},r) $ 内有界。

为证明其逆，我们令

\[
M(\rho) := \max_{z\in\partial\mathbb{D}(z_{0},\rho)}|f(z)| \quad\text{对 } \rho \in (0,r) .
\]

由假设，存在 $ M \ge 0 $ 使得对 $ \rho \in (0,r) $ 有 $ M(\rho) \le M $（因为 $ |f| $ 在 $ \overline{\mathbb{D}}(z_{0},r)\setminus\{z_{0}\} $ 上连续，从而对每个 $ 0 < r_{0} < r $ 在紧致集 $ \overline{\mathbb{D}}(z_{0},r)\setminus\mathbb{D}(z_{0},r_{0}) $ 上有界）。于是由 (6.6) 可得

\[
|c_{n}| \le M(\rho)\rho^{-n} \le M\rho^{-n} \quad\text{对 } n \in \mathbb{Z} \text{ 且 } \rho \in (0,r) .
\]

因此 $ f $ 的洛朗展开的主部为零，由推论 6.4 可得

\[
f(z) = \sum_{n=0}^{\infty} c_{n}(z - z_{0})^{n} \quad\text{对 } z \in \mathbb{D}^{\bullet}(z_{0},r) .
\]

由

\[
z \mapsto \sum_{n=0}^{\infty} c_{n}(z - z_{0})^{n} \quad\text{对 } z \in \mathbb{D}(z_{0},r) ,
\]

定义的函数在 $ \mathbb{D}(z_{0},r) $ 上全纯，且在 $ \mathbb{D}^{\bullet}(z_{0},r) $ 上与 $ f $ 一致。因此 $ z_{0} $ 是 $ f $ 的可去奇点。\nobreak\hspace{0.35em}■

\end{theorem}

\needspace{3\baselineskip}\subsection*{孤立奇点}

设 $ f : U\setminus\{z_{0}\} \to \mathbb{C} $ 全纯，且 $ r > 0 $ 满足 $ \overline{\mathbb{D}}(z_{0},r) \subset U $。进一步设

\[
f(z) = \sum_{n=-\infty}^{\infty}c_{n}(z - z_{0})^{n} \quad \text{对 } z \in \mathbb{D}^{\bullet}(z_{0},r)
\]

是 $ f $ 在 $ \mathbb{D}^{\bullet}(z_{0},r) $ 中的洛朗展开。若 $ z_{0} $ 是 $ f $ 的奇点，则当它不是可去奇点时，称它是\textbf{孤立的}。由黎曼可去奇点定理（的证明），这成立当且仅当 $ f $ 的洛朗展开的主部不恒等于零。若 $ z_{0} $ 是 $ f $ 的孤立奇点，且存在 $ m \in \mathbb{N}^{\times} $ 使得 $ c_{-m} \ne 0 $ 且对 $ n > m $ 有 $ c_{-n} = 0 $，则称 $ z_{0} $ 是 $ f $ 的\textbf{极点}。此时 $ m $ 是该极点的\textbf{阶}。若洛朗级数主部中无穷多个系数异于零，则 $ z_{0} $ 是 $ f $ 的\textbf{本性奇点}。最后，我们通过下式定义 $ f $ 在 $ z_{0} $ 处的\textbf{留数}

\[
\operatorname{Res}(f,z_{0}) := c_{-1} .
\]

函数 $ g $ 在 $ U $ 中称为\textbf{亚纯的}，若存在 $ U $ 的闭子集 $ P(g) $，使得 $ g $ 在 $ U\setminus P(g) $ 上全纯，且每个 $ z \in P(g) $ 都是 $ g $ 的极点。\footnote{$ P(g) $ 也可以是空集。因此 $ U $ 上的每个全纯函数也是亚纯的。} 此时 $ P(g) $ 是 $ g $ 的\textbf{极点集}。

\begin{remark}

\begin{items}
\item

设 $ f : U\setminus\{z_{0}\} \to \mathbb{C} $ 全纯。则有

\[
\operatorname{Res}(f,z_{0}) = \frac{1}{2\pi i}\int_{\partial\mathbb{D}(z_{0},r)}f(z)\,dz
\]

对每个满足 $ \overline{\mathbb{D}}(z_{0},r) \subset U $ 的 $ r > 0 $ 成立。因此，$ f $ 在 $ z_{0} $ 处的留数（相差因子 $ 1/2\pi i $）就是对 $ f $ 沿道路 $ \partial\mathbb{D}(z_{0},r) $ 积分之后所“剩余”（即“残余”）的量。

\proofhead
这由推论 6.4 与例 5.3(a) 得到。\nobreak\hspace{0.35em}■
\item

$ U $ 中亚纯函数 $ f $ 的极点集 $ P(f) $ 是离散的、可数的，且该集合在 $ U $ 中没有聚点。\footnote{然而，$ f $ 的极点完全可能聚积在 $ U $ 的边界上。}

\end{items}

证明

\begin{subitems}
\item

设 $ z_{0} \in P(f) $。则存在 $ r > 0 $ 使得 $ \overline{\mathbb{D}}(z_{0},r) \subset U $。从而 $ f $ 在 $ \mathbb{D}^{\bullet}(z_{0},r) $ 中全纯。因此 $ P(f)\cap\mathbb{D}(z_{0},r) = \{z_{0}\} $，这表明 $ P(f) $ 是离散的。

\item

假设 $ P(f) $ 在 $ U $ 中有聚点 $ z_{0} $。由于 $ P(f) $ 是离散的，$ z_{0} $ 不属于 $ P(f) $。因此 $ z_{0} $ 位于开集 $ U\setminus P(f) $ 中，且存在 $ r > 0 $ 使得 $ \mathbb{D}(z_{0},r) \subset U\setminus P(f) $。因此 $ z_{0} $ 不是 $ P(f) $ 的聚点，矛盾。

\item

对每个 $ z \in P(f) $，存在 $ r_{z} > 0 $ 使得 $ \mathbb{D}^{\bullet}(z,r_{z})\cap P(f) = \varnothing $。若 $ K $ 是 $ U $ 的紧致子集，则 $ K\cap P(f) $ 也是紧致的。因此存在 $ z_{0},\ldots,z_{m} \in P(f) $ 使得

\[
K\cap P(f) \subset \bigcup_{j=0}^{m}\mathbb{D}(z_{j},r_{z_{j}}) .
\]

因此 $ K\cap P(f) $ 是有限集。

\item

为验证 $ P(f) $ 可数，令

\[
K_{j} := \{\, x \in U ;\ d(x,\partial U) \ge 1/j,\ |x| \le j \,\} \quad \text{对 } j \in \mathbb{N}^{\times} .
\]

由例 III.1.3(l) 与 III.2.22(c)，以及海涅–博雷尔定理，每个 $ K_{j} $ 都是紧致的。此外 $ \bigcup_{j}K_{j} = U $，且对每个 $ j \in \mathbb{N}^{\times} $，$ K_{j}\cap P(f) $ 是有限集。于是由命题 I.6.8 可知 $ P(f) = \bigcup_{j}\left(K_{j}\cap P(f)\right) $ 可数。\nobreak\hspace{0.35em}■

\end{subitems}

下一个定理表明，函数亚纯当且仅当它局部是两个全纯函数之商。\footnote{这解释了“亚纯”这一名称，其意为“分式形式”，而“全纯”可译为“整形式”。}

\end{remark}

\begin{proposition}

函数 $ f $ 在 $ U $ 中亚纯当且仅当存在 $ U $ 的闭子集 $ A $ 且满足以下条件：

\begin{subitems}
\item

$ f $ 在 $ U\setminus A $ 中全纯。

\item

对每个 $ a \in A $，存在 $ r > 0 $ 使得 $ \mathbb{D}(a,r) \subset U $，存在 $ g, h \in C^{\omega}\left(\mathbb{D}(a,r)\right) $ 且 $ h \ne 0 $，且 $ f = g/h $ 在 $ \mathbb{D}^{\bullet}(a,r) $ 中成立。

\end{subitems}

证明

\begin{items}
\item

设 $ f $ 在 $ U $ 中亚纯。则对 $ a \in P(f) =: A $，存在 $ r > 0 $ 满足 $ \mathbb{D}(a,r) \subset U $，使得 $ f $ 有洛朗展开

\[
f(z) = \sum_{n=-m}^{\infty}c_{n}(z - a)^{n} \quad \text{对 } z \in \mathbb{D}^{\bullet}(a,r)
\]

（在 $ \mathbb{D}^{\bullet}(a,r) $ 中），其中 $ m \in \mathbb{N}^{\times} $ 适当。全纯函数

\begin{equation}
\mathbb{D}^{\bullet}(a,r) \to \mathbb{C} , \quad z \mapsto (z - a)^{m}f(z) = \sum_{k=0}^{\infty}c_{k-m}(z - a)^{k}\tag{6.7}
\end{equation}

在 $ a $ 处有可去奇点。因此存在 $ g \in C^{\omega}\left(\mathbb{D}(a,r)\right) $ 使得

\[
g(z) = (z - a)^{m}f(z) \quad \text{对 } 0 < |z - a| < r .
\]

因此 $ f $ 在 $ \mathbb{D}^{\bullet}(a,r) $ 中有表示 $ f = g/h $，其中 $ h := (X - a)^{m} \in C^{\omega}(\mathbb{C}) $。

\item

设给定条件满足。若 $ h(a) \ne 0 $，则可（通过缩小 $ r $）假设对 $ z \in \mathbb{D}(a,r) $ 有 $ h(z) \ne 0 $。则 $ f = g/h $ 在 $ \mathbb{D}(a,r) $ 中全纯。在 $ h(a) = 0 $ 的情形，可假设存在 $ m \in \mathbb{N}^{\times} $ 使得 $ h $ 有幂级数展开

\[
h(z) = \sum_{k=m}^{\infty}c_{k}(z - a)^{k} = (z - a)^{m}\sum_{n=0}^{\infty}c_{n+m}(z - a)^{n}
\]

（在 $ \mathbb{D}(a,r) $ 中），其中 $ c_{m} $ 异于零，因为 $ h \ne 0 $。因此由下式定义的函数

\[
\varphi(z) := \sum_{n=0}^{\infty}c_{n+m}(z - a)^{n} \quad \text{对 } z \in \mathbb{D}(a,r)
\]

在 $ \mathbb{D}(a,r) $ 上全纯（见命题 V.3.5），且 $ \varphi(a) = c_{m} \ne 0 $。因此存在 $ \rho \in (0,r) $ 使得对 $ |z - a| \le \rho $ 有 $ \varphi(z) \ne 0 $，这意味着 $ g/\varphi \in C^{\omega}\left(\mathbb{D}(a,\rho)\right) $。记 $ g/\varphi $ 的泰勒级数为 $ \sum_{n\ge 0}b_{n}(z - a)^{n} $，则得

\[
f(z) = \frac{g(z)}{h(z)} = \frac{1}{(z - a)^{m}}\frac{g}{\varphi}(z) = \sum_{n=0}^{\infty}b_{n}(z - a)^{n-m} \quad \text{对 } z \in \mathbb{D}^{\bullet}(a,\rho) .
\]

由这一表达式以及洛朗展开的唯一性可知，$ a $ 是 $ f $ 的极点。\nobreak\hspace{0.35em}■

\end{items}

\end{proposition}

\begin{example}

\begin{items}
\item

每个有理函数在 $ \mathbb{C} $ 中亚纯，且只有有限多个极点。

\item

正切与余切在 $ \mathbb{C} $ 中亚纯。它们的极点集分别为 $ \pi(\mathbb{Z} + 1/2) $ 与 $ \pi\mathbb{Z} $。余切在 $ \pi\mathbb{D}^{\bullet} $ 中的洛朗展开为

\[
\cot z = \frac{1}{z} - 2\sum_{k=1}^{\infty}\frac{\zeta(2k)}{\pi^{2k}}z^{2k-1} \quad \text{对 } z \in \pi\mathbb{D}^{\bullet} .
\]

\proofhead
因为 $ \tan = \sin/\cos $ 且 $ \cot = 1/\tan $，前两个断言由命题 6.8 得到。余切的上述表示由 (VI.7.18)、定理 VI.6.15 以及洛朗展开的唯一性推出。\nobreak\hspace{0.35em}■
\item

伽马函数在 $ \mathbb{C} $ 中亚纯。它的极点集是 $ -\mathbb{N} $。

\proofhead
由命题 VI.9.5 的魏尔斯特拉斯乘积表示与魏尔斯特拉斯收敛定理（定理 5.27）可知，$ \Gamma $ 是某个整函数的倒数，且该整函数的零点集为 $ -\mathbb{N} $。于是断言由命题 6.8 得到。\nobreak\hspace{0.35em}■
\item

黎曼 $ \zeta $ 函数在 $ \mathbb{C} $ 中亚纯。

\proofhead
这由定理 VI.6.15 得到。\nobreak\hspace{0.35em}■
\item

函数 $ z \mapsto e^{1/z} $ 在 $ \mathbb{C} $ 中不是亚纯的。

\proofhead
因为 $ e^{1/z} = \exp(1/z) $，有
\[
e^{1/z} = \sum_{k=0}^{\infty}\frac{1}{k!\,z^{k}} = 1 + \frac{1}{z} + \frac{1}{2z^{2}} + \cdots \quad \text{对 } z \in \mathbb{C}^{\times} .
\]

因此 $ 0 $ 是 $ z \mapsto e^{1/z} $ 的本性奇点。\nobreak\hspace{0.35em}■

\end{items}

\end{example}

\needspace{3\baselineskip}\subsection*{单极点}

我们将看到，亚纯函数的留数起着特别重要的作用。因此，我们应设法在不显式作洛朗展开的情况下确定留数。对一阶极点——即\textbf{单极点}——这一点尤其容易，我们接下来说明。

\begin{proposition}

全纯函数 $ f : U\setminus\{z_{0}\} \to \mathbb{C} $ 在 $ z_{0} $ 处有单极点当且仅当 $ z \mapsto g(z) := (z - z_{0})f(z) $ 在 $ z_{0} $ 处有可去奇点且 $ g(z_{0}) \ne 0 $。此时

\[
\operatorname{Res}(f,z_{0}) = \lim_{z\to z_{0}}(z - z_{0})f(z) .
\]

\proofhead
设 $ g $ 在 $ U $ 中全纯且 $ g(z_{0}) \ne 0 $。则存在 $ r > 0 $ 满足 $ \mathbb{D}(z_{0},r) \subset U $，以及 $ \mathbb{C} $ 中的序列 $ (b_{n}) $，使得
\[
g(z) = \sum_{n=0}^{\infty}b_{n}(z - z_{0})^{n} \quad \text{对 } z \in \mathbb{D}(z_{0},r) \text{ 且 } g(z_{0}) = b_{0} \ne 0 .
\]

因为

\[
f(z) = \frac{g(z)}{z - z_{0}} = \sum_{n=-1}^{\infty}c_{n}(z - z_{0})^{n} \quad \text{对 } z \in \mathbb{D}^{\bullet}(z_{0},r) ,
\]

其中 $ c_{n} := b_{n+1} $ 且 $ c_{-1} := b_{0} \ne 0 $，点 $ z_{0} $ 是单极点，且

\[
\operatorname{Res}(f,z_{0}) = c_{-1} = b_{0} = g(z_{0}) = \lim_{z\to z_{0}}(z - z_{0})f(z) .
\]

反之，若 $ z_{0} $ 是 $ f $ 的单极点，则存在 $ r > 0 $ 使得 $ \mathbb{D}(z_{0},r) \subset U $，且

\[
f(z) = \sum_{n=-1}^{\infty}c_{n}(z - z_{0})^{n} \quad \text{对 } z \in \mathbb{D}^{\bullet}(z_{0},r) \text{ 且 } c_{-1} \ne 0 .
\]

由此可得

\[
(z - z_{0})f(z) = \sum_{n=0}^{\infty}c_{n-1}(z - z_{0})^{n} \quad \text{对 } z \in \mathbb{D}^{\bullet}(z_{0},r) .
\]

现在断言由黎曼可去奇点定理与解析函数的唯一性定理推出。\nobreak\hspace{0.35em}■

\end{proposition}

\begin{example}

\begin{items}
\item

设 $ g $ 与 $ h $ 在 $ U $ 中全纯，且 $ h $ 在 $ z_{0} $ 处有\textbf{单零点}，即\footnote{单零点即阶为 1 的零点（见习题 IV.3.10）。} $ h(z_{0}) = 0 $ 且 $ h'(z_{0}) \ne 0 $。则 $ f := g/h $ 在 $ U $ 中亚纯，且若 $ g(z_{0}) \ne 0 $，则 $ z_{0} $ 是 $ f $ 的单极点，其留数为 $ \operatorname{Res}(f,z_{0}) = g(z_{0})/h'(z_{0}) $。

\proofhead
泰勒公式给出
\[
h(z) = (z - z_{0})h'(z_{0}) + o(|z - z_{0}|) \quad (z \to z_{0}) ,
\]

因此当 $ z \to z_{0} $ 时 $ h(z)(z - z_{0})^{-1} \to h'(z_{0}) $。这意味着

\[
\lim_{z\to z_{0}}(z - z_{0})f(z) = \lim_{z\to z_{0}}g(z)/\left(h(z)(z - z_{0})^{-1}\right) = g(z_{0})/h'(z_{0}) ,
\]

于是断言由定理 6.8 与 6.10 以及黎曼可去奇点定理得到。\nobreak\hspace{0.35em}■

\item

正切与余切只有单极点。它们的留数为

\[
\operatorname{Res}\left(\tan,\pi(k + 1/2)\right) = -\operatorname{Res}(\cot,k\pi) = -1 \quad \text{对 } k \in \mathbb{Z} .
\]

\proofhead
这立即由 (a) 得到。\nobreak\hspace{0.35em}■
\item

伽马函数只有一阶极点，且

\[
\operatorname{Res}(\Gamma,-n) = (-1)^{n}/n! \quad \text{对 } n \in \mathbb{N} .
\]

\proofhead
由 (VI.9.2)，对满足 $ \operatorname{Re} z > -n - 1 $ 的 $ z \in \mathbb{C}\setminus(-\mathbb{N}) $，得表示
\[
(z + n)\Gamma(z) = \frac{\Gamma(z + n + 1)}{z(z + 1)\cdot\cdot\cdot\cdot\cdot(z + n - 1)} .
\]

因此当 $ z \to -n $ 时 $ (z + n)\Gamma(z) $ 收敛到 $ (-1)^{n}\Gamma(1)/n! $。因为 $ \Gamma(1) = 1 $，断言由例 6.9(c) 与命题 6.10 推出。\nobreak\hspace{0.35em}■

\item

黎曼 $ \zeta $ 函数在 $ 1 $ 处有单极点，留数为 $ 1 $。

\proofhead
这由定理 VI.6.15 得到。\nobreak\hspace{0.35em}■
\item

设 $ p \in \mathbb{R} $ 且 $ a > 0 $。则由 $ f(z) := e^{-ipz}/(z^{2} + a^{2}) $ 定义的函数在 $ \mathbb{C} $ 中亚纯，且在 $ \pm ia $ 处有单极点，其留数为

\[
\operatorname{Res}(f,\pm ia) = \pm e^{\pm pa}/2ia .
\]

\proofhead
显然 $ (z \mp ia)f(z) = e^{-ipz}/(z \pm ia) $ 在 $ \pm ia $ 处全纯，且
\[
\lim_{z\to\pm ia}(z \mp ia)f(z) = \pm e^{\pm pa}/2ia .
\]

因此断言由命题 6.10 得到。\nobreak\hspace{0.35em}■

\item

对 $ p \in \mathbb{R} $ 与 $ f(z) := e^{-ipz}/(z^{4} + 1) $，函数 $ f $ 在 $ \mathbb{C} $ 中亚纯，且

\[
P(f) = \{\, z_{j} := e^{i(\pi/4+j\pi/2)} ;\ j = 0,\ldots,3 \,\} .
\]

每个极点都是单极点，且每个留数为

\[
\operatorname{Res}(f,z_{j}) = e^{-ipz_{j}}/4z_{j}^{3} \quad \text{对 } 0 \le j \le 3 .
\]

\proofhead
初等计算给出
\[
\prod_{k=0,\,k\ne j}^{3}(z_{j} - z_{k}) = 4z_{j}^{3} .
\]

因此这一断言也由命题 6.10 得到。\nobreak\hspace{0.35em}■

\end{items}

\end{example}

\needspace{3\baselineskip}\subsection*{绕数}

我们已在上一节看到，全纯函数的线积分的同伦不变性（命题 5.7）可用来有效地计算其他积分。现在我们来对亚纯函数推导类似的结果。为此，我们必须更多地了解极点相对于积分曲线的位置。

在下文中

• $ \Gamma $ 总是 $ \mathbb{C} $ 中的一条闭的紧致分段 $ C^{1} $ 曲线。

对 $ a \in \mathbb{C}\setminus\Gamma $，

\[
w(\Gamma,a) := \frac{1}{2\pi i}\int_{\Gamma}\frac{dz}{z - a}
\]

称为 $ \Gamma $ 关于 $ a $ 的\textbf{绕数}、\textbf{循环数}或\textbf{指数}。

\begin{example}

\begin{items}
\item

设 $ m \in \mathbb{Z}^{\times} $ 且 $ r > 0 $。则对 $ z_{0} \in \mathbb{C} $，参数化 $ \gamma_{m} : [0,2\pi] \to z_{0} + re^{imt} $ 是一条光滑曲线 $ \Gamma_{m} := \Gamma_{m}(z_{0},r) $ 的参数化，其像是定向圆 $ \partial\mathbb{D}(z_{0},r) $。若 $ m $ 为正 [或负]，则 $ \Gamma_{m} $ 与 $ \partial\mathbb{D}(z_{0},r) $ 定向相同 [或相反]。因此当 $ t $ 从 $ 0 $ 变到 $ 2\pi $ 时，$ \partial\mathbb{D}(z_{0},r) $ 将沿正 [或负] 方向环绕 $ |m| $ 次。基于这些理由，我们称 $ \Gamma_{m} $ 为以 $ z_{0} $ 为中心、以 $ r $ 为半径的 \textbf{$ m $ 次遍历圆}。我们有

\begin{equation*}
w(\Gamma_{m},a) = \begin{cases} m , & |a - z_{0}| < r , \\ 0 , & |a - z_{0}| > r . \end{cases}
\end{equation*}

\proofhead
同引理 6.1 的证明一样，在情形 $ |a - z_{0}| < r $ 下可得
\[
\int_{\Gamma_{m}}\frac{dz}{z - a} = \int_{\Gamma_{m}}\frac{dz}{z - z_{0}} = \int_{0}^{2\pi}\frac{i\,mre^{imt}}{re^{imt}}\,dt = 2\pi i m .
\]

当 $ |a - z_{0}| > r $ 时，曲线 $ \Gamma_{m} $ 在 $ \mathbb{C}\setminus\{a\} $ 中零伦。\nobreak\hspace{0.35em}■

\item

对 $ z_{0} \in \mathbb{C} $，设 $ \Gamma_{z_{0}} $ 是像 $ \operatorname{im}(\Gamma_{z_{0}}) = \{z_{0}\} $ 的点曲线。则对 $ a \in \mathbb{C}\setminus\{z_{0}\} $，$ w(\Gamma_{z_{0}},a) = 0 $。

\item

若 $ \gamma_{1} $ 与 $ \gamma_{2} $ 是 $ U $ 中同伦的分段 $ C^{1} $ 闭路，且 $ a \in U^{c} $，则

\[
w(\Gamma_{1},a) = w(\Gamma_{2},a)
\]

其中 $ \Gamma_{1} = [\gamma_{1}] $ 且 $ \Gamma_{2} = [\gamma_{2}] $。

\proofhead
对 $ a \in U^{c} $，函数 $ z \mapsto 1/(z - a) $ 在 $ U $ 中全纯。因此结论由命题 5.7 可得。\nobreak\hspace{0.35em}■
\item

若 $ U $ 单连通且 $ \Gamma \subset U $，则对 $ a \in U^{c} $，$ w(\Gamma,a) = 0 $。

\proofhead
这由 (b) 和 (c) 可得。\nobreak\hspace{0.35em}■
\end{items}

根据例 6.12(a)，$ m $ 次遍历圆 $ \Gamma_{m} := \Gamma_{m}(z_{0},r) $ 的绕数 $ w(\Gamma_{m},a) $ 是 $ \Gamma_{m} $ “绕”点 $ a $ 的（带符号的）次数。下面我们证明，绕数的这种几何解释对每条闭的分段 $ C^{1} $ 曲线都成立，但我们必须先证明一个技术性结果。

\end{example}

\begin{lemma}

设 $ I $ 是一个完全紧致区间，且 $ \gamma : I \to \mathbb{C}^{\times} $ 有分段连续导数。则存在一个连续且分段连续可微的函数 $ \varphi : I \to \mathbb{C} $ 使得 $ \exp\circ\varphi = \gamma $。这里，$ \gamma $ 与 $ \varphi $ 在同一个区间 $ I $ 上有连续导数。

证明

\begin{subitems}
\item

对 $ \alpha \in \mathbb{R} $，我们定义 $ \log_{\alpha} := \operatorname{Log}|\left(\mathbb{C}\setminus\mathbb{R}^{+}e^{i\alpha}\right) $。则由命题 III.6.19 与习题 III.6.9，$ \log_{\alpha} $ 是 $ \mathbb{C}\setminus\mathbb{R}^{+}e^{i\alpha} $ 到 $ \mathbb{R} + i(\alpha,\alpha + 2\pi) $ 上的拓扑映射，且满足

\begin{equation}
\exp(\log_{\alpha}z) = z \quad \text{对 } z \in \mathbb{C}\setminus\mathbb{R}^{+}e^{i\alpha} .\tag{6.8}
\end{equation}

由此我们得到（见例 IV.1.13(e)）$ \log_{\alpha} $ 是 $ \mathbb{C}\setminus\mathbb{R}^{+}e^{i\alpha} $ 到 $ \mathbb{R} + i(\alpha,\alpha + 2\pi) $ 上的 $ C^{1} $ 微分同胚，且满足

\begin{equation}
(\log_{\alpha})'(z) = 1/z \quad \text{对 } z \in \mathbb{C}\setminus\mathbb{R}^{+}e^{i\alpha} .\tag{6.9}
\end{equation}
\item

因为 $ \gamma(I) $ 是紧致的，存在 $ r > 0 $ 使得 $ r\overline{\mathbb{D}}\cap\gamma(I) = \varnothing $。由 $ \gamma $ 的一致连续性，存在 $ I $ 的划分 $ (t_{0},\ldots,t_{m}) $ 使得对 $ 1 \le j \le m $ 有 $ \operatorname{diam}\left(\gamma|[t_{j-1},t_{j}]\right) < r/2 $。因为 $ \gamma $ 有分段连续导数，我们可以选取该划分使得对 $ 1 \le j \le m $，$ \gamma|[t_{j-1},t_{j}] $ 连续可微。因为圆盘 $ D_{j} := \mathbb{D}(\gamma(t_{j}),r) $ 不含零点，我们可以固定一个 $ \alpha_{j} \in (-\pi,\pi) $ 使得对 $ 1 \le j \le m $ 有 $ \mathbb{R}^{+}e^{i\alpha_{j}}\cap\overline{D}_{j} = \varnothing $。然后我们令 $ \log_{j} := \log_{\alpha_{j}} $ 与 $ \varphi_{j} := \log_{j}\circ\gamma|[t_{j-1},t_{j}] $。由 (i)，我们知道 $ \varphi_{j} $ 属于 $ C^{1}\left([t_{j-1},t_{j}]\right) $，又因为对 $ t \in [t_{j-1},t_{j}] $ 有 $ \gamma(t) \in D_{j}\cap D_{j+1} $，利用 (6.8) 我们得到

\[
\exp\left(\varphi_{j}(t_{j})\right) = \gamma(t_{j}) = \exp\left(\varphi_{j+1}(t_{j})\right) \quad \text{对 } 1 \le j \le m - 1 .
\]

因此，命题 III.6.13 与指数函数的加法定理保证存在 $ k_{j} \in \mathbb{Z} $ 使得

\begin{equation}
\varphi_{j}(t_{j}) - \varphi_{j+1}(t_{j}) = 2\pi ik_{j} \quad \text{对 } 1 \le j \le m - 1 .\tag{6.10}
\end{equation}

现在我们定义 $ \varphi : I \to \mathbb{C} $ 为

\begin{equation}
\varphi(t) := \varphi_{j}(t) + 2\pi i\sum_{n=0}^{j-1}k_{n} \quad \text{对 } t_{j-1} \le t \le t_{j} \text{ 且 } 1 \le j \le m ,\tag{6.11}
\end{equation}

其中 $ k_{0} := 0 $。则由 (6.10) 以及 $ \log_{j} $ 与 $ \gamma|[t_{j-1},t_{j}] $ 的连续可微性可得，$ \varphi $ 有分段连续导数。最后，由 (6.8)、$ \varphi_{j} $ 的定义以及指数函数的 $ 2\pi i $ 周期性，我们得到 $ \exp\circ\varphi = \gamma $。\nobreak\hspace{0.35em}■

\end{subitems}

\end{lemma}

\begin{theorem}

对每个 $ a \in \mathbb{C}\setminus\Gamma $，指数 $ w(\Gamma,a) $ 是整数。

\proofhead
设 $ \gamma $ 是 $ \Gamma $ 的一个分段 $ C^{1} $ 参数化，且 $ (t_{0},\ldots,t_{m}) $ 是参数区间 $ I $ 的划分，使得对 $ 1 \le j \le m $，$ \gamma|[t_{j-1},t_{j}] $ 连续可微。则根据引理 6.13，对 $ a \in \mathbb{C}\setminus\Gamma $ 存在一个 $ \varphi \in C(I) $ 使得对 $ 1 \le j \le m $ 有 $ \varphi|[t_{j-1},t_{j}] \in C^{1}[t_{j-1},t_{j}] $，且 $ e^{\varphi} = \gamma - a $。由此可得，对 $ t_{j-1} \le t \le t_{j} $ 与 $ 1 \le j \le m $ 有 $ \dot\gamma(t) = \dot\varphi(t)\left(\gamma(t) - a\right) $。因此我们得到
\begin{equation}
\int_{\Gamma}\frac{dz}{z - a} = \sum_{j=1}^{m}\int_{t_{j-1}}^{t_{j}}\frac{\dot\gamma(t)\,dt}{\gamma(t) - a} = \sum_{j=1}^{m}\int_{t_{j-1}}^{t_{j}}\dot\varphi(t)\,dt = \varphi(t_{m}) - \varphi(t_{0}) .\tag{6.12}
\end{equation}

因为 $ \Gamma $ 是闭的，我们有

\[
\exp\left(\varphi(t_{m})\right) = E_{\Gamma} - a = A_{\Gamma} - a = \exp\left(\varphi(t_{0})\right) ,
\]

因此根据命题 III.6.13(i)，$ \varphi(t_{m}) - \varphi(t_{0}) \in 2\pi i\mathbb{Z} $。所以 $ w(\Gamma,a) $ 是整数。\nobreak\hspace{0.35em}■

\end{theorem}

\begin{remark}

\begin{items}
\item

设引理 6.13 的假设满足。用其证明中的记号，对 $ 1 \le j \le m $ 与 $ t \in [t_{j-1},t_{j}] $，我们用 $ \arg_{j}(\gamma(t)) $ 记满足 $ \operatorname{Im}(\varphi_{j}(t)) = \eta $ 的唯一数 $ \eta \in (\alpha_{j},\alpha_{j} + 2\pi) $。则由

\[
e^{\log|\gamma(t)|} = |\gamma(t)| = |e^{\varphi(t)}| = e^{\operatorname{Re}\varphi(t)}
\]

可得

\begin{equation}
\varphi_{j}(t) = \log_{j}\gamma(t) = \log|\gamma(t)| + i\arg_{j}(\gamma(t))\tag{6.13}
\end{equation}

对 $ t_{j-1} \le t \le t_{j} $ 与 $ 1 \le j \le m $ 成立。现在我们令

\[
\arg_{c,\gamma}(t) := \arg_{j}(\gamma(t)) + 2\pi\sum_{n=0}^{j-1}k_{n}
\]

对 $ t_{j-1} \le t \le t_{j} $ 与 $ 1 \le j \le m $ 成立。则 (6.11) 与 (6.13) 蕴含

\begin{equation}
\varphi = \log\circ|\gamma| + i\arg_{c,\gamma} .\tag{6.14}
\end{equation}

因为 $ \varphi $ 分段连续可微，(6.14) 表明这对 $ \arg_{c,\gamma} $ 也成立，其中 $ \arg_{c,\gamma} $ 与 $ \gamma $ 在 $ I $ 的相同子区间上连续可微。换言之，$ \arg_{c,\gamma} $ 是一个分段连续可微函数，它“从”集值函数 $ \operatorname{Arg}\circ\gamma $ “选取”，即

\[
\arg_{c,\gamma} \in \operatorname{Arg}(\gamma(t)) \quad \text{对 } t \in I .
\]

同样，$ \varphi $ 是集值对数 $ \operatorname{Log}\circ\gamma $ 的一个分段连续可微的“选取”。

\item

设 $ \gamma $ 是 $ \Gamma $ 的一个分段 $ C^{1} $ 参数化，且 $ a \in \mathbb{C}\setminus\Gamma $。又设

\[
\varphi = \log\circ|\gamma - a| + i\arg_{c,\gamma-a}
\]

是 $ \operatorname{Log}\circ(\gamma - a) $ 的一个分段 $ C^{1} $ 选取，其中对 $ 1 \le j \le m $，$ \varphi $ 属于 $ C^{1}[t_{j-1},t_{j}] $。则我们有

\[
\int_{t_{j-1}}^{t_{j}}\dot\varphi(t)\,dt = \log\circ|\gamma - a|\big|_{t_{j-1}}^{t_{j}} + i\arg_{c,\gamma-a}\big|_{t_{j-1}}^{t_{j}} = \varphi(t_{j}) - \varphi(t_{j-1}) .
\]

因为 $ \log|\gamma(t_{m}) - a| = \log|\gamma(t_{0}) - a| $，因此由 (6.12) 可得

\[
w(\Gamma,a) = \frac{1}{2\pi i}\int_{\Gamma}\frac{dz}{z - a} = \frac{1}{2\pi}\left(\arg_{c,\gamma-a}(t_{m}) - \arg_{c,\gamma-a}(t_{0})\right) .
\]

这表明 $ \Gamma $ 关于 $ a $ 的绕数的 $ 2\pi $ 倍是当 $ \Gamma $ 从 $ \gamma(t_{0}) $ 走到 $ \gamma(t_{m}) $ 时 $ \gamma - a $ 的辐角 $ \arg_{c,\gamma-a} $ 的“累积改变”。因此 $ w(\Gamma,a) $ 刻画了 $ \Gamma $ 绕点 $ a $ 的次数，其中 $ w(\Gamma,a) > 0 $ 表示它逆时针绕行，$ w(\Gamma,a) < 0 $ 表示它顺时针绕行。\nobreak\hspace{0.35em}■

\end{items}

\bookfigure[0.80]{figures/bdf78ad8b4e1786c.png}

一条绕数 $ w(\Gamma,0) = 3 $ 的曲线

\end{remark}

\needspace{3\baselineskip}\subsection*{绕数的连续性}

我们首先对离散空间中的映射证明一个简单的引理。

\begin{lemma}

设 $ X $ 与 $ Y $ 是度量空间且 $ Y $ 离散。若 $ f : X \to Y $ 连续，则 $ f $ 在 $ X $ 的每个连通分支上为常数。

\proofhead
设 $ Z $ 是 $ X $ 的一个连通分支。由定理 III.4.5，$ f(Z) $ 连通。因为 $ Y $ 离散，$ Y $ 的每个连通分支恰好由一个点构成。因此 $ f $ 在 $ Z $ 上为常数。\nobreak\hspace{0.35em}■
\end{lemma}

因为 $ \Gamma $ 紧致，存在 $ R > 0 $ 使得 $ \Gamma \subset R\mathbb{D} $。因此 $ \mathbb{C}\setminus\Gamma $ 包含集合 $ R\mathbb{D}^{c} $，这表明 $ \mathbb{C}\setminus\Gamma $ 恰有一个无界连通分支。

\begin{corollary}

映射 $ w(\Gamma,\cdot) : \mathbb{C}\setminus\Gamma \to \mathbb{Z} $ 在每个连通分支上为常数。若 $ a $ 属于 $ \mathbb{C}\setminus\Gamma $ 的无界连通分支，则 $ w(\Gamma,a) = 0 $。

\proofhead
设 $ a \in \mathbb{C}\setminus\Gamma $，并定义 $ d := d(a,\Gamma) $。然后我们选取一个 $ \varepsilon > 0 $ 并定义 $ \delta := \min\{\varepsilon\pi d^{2}\backslash L(\Gamma),d/2\} $。则我们有
\[
|z - a| \ge d > 0 \quad \text{且}\quad |z - b| \ge d/2 \quad \text{对 } z \in \Gamma \text{ 且 } b \in \mathbb{D}(a,\delta) .
\]

由此可得

\[
|w(\Gamma,a) - w(\Gamma,b)| \le \frac{1}{2\pi}\int_{\Gamma}\left|\frac{a - b}{(z - a)(z - b)}\right|dz < \frac{1}{2\pi}L(\Gamma)\frac{2}{d^{2}}\delta \le \varepsilon
\]

对 $ b \in \mathbb{D}(a,\delta) $ 成立。因为 $ \varepsilon $ 是任意的，$ w(\Gamma,\cdot) $ 在 $ a \in \mathbb{C}\setminus\Gamma $ 处连续。第一个陈述现在由定理 6.14 与引理 6.16 可得。

以 $ Z $ 记 $ \mathbb{C}\setminus\Gamma $ 的无界连通分支。又设已选取 $ R > 0 $ 使得 $ Z $ 包含集合 $ R\mathbb{D}^{c} $。最后，我们选取 $ a \in Z $ 使得 $ |a| > R $ 且 $ |a| > L(\Gamma)/\pi + \max\{|z| ; z \in \Gamma\} $。则我们有

\[
|w(\Gamma,a)| \le \frac{1}{2\pi}\int_{\Gamma}\frac{dz}{|z - a|} < \frac{1}{2} .
\]

因为 $ w(\Gamma,a) \in \mathbb{Z} $，我们有 $ w(\Gamma,a) = 0 $，因此对每个 $ b \in Z $ 有 $ w(\Gamma,b) = 0 $。\nobreak\hspace{0.35em}■

\end{corollary}

\bookfigure[0.44]{figures/459a4d89d537d254.png}

\begin{corollary}

设 $ f $ 在 $ U $ 中亚纯，且对 $ a \in U^{c} $ 有 $ w(\Gamma,a) = 0 $。则 $ \{z \in P(f)\setminus\Gamma ; w(\Gamma,z) \ne 0\} $ 是有限集。

\proofhead
由推论 6.17，$ B := \{z \in U\setminus\Gamma ; w(\Gamma,z) \ne 0\} $ 有界。当 $ B\cap P(f) $ 不是有限集时，该集合有聚点 $ z_{0} \in \overline{B} $。根据注 6.7(b)，$ P(f) $ 在 $ U $ 中有聚点，因此 $ z_{0} $ 属于 $ U^{c} $。于是由假设，$ w(\Gamma,z_{0}) = 0 $。另一方面，由 $ w(\Gamma,\cdot) $ 的连续性与 $ w(\Gamma,B) \subset \mathbb{Z}^{\times} $ 可得，$ w(\Gamma,z_{0}) $ 异于零。因此 $ B\cap P(f) $ 是有限集。\nobreak\hspace{0.35em}■
\end{corollary}

\needspace{3\baselineskip}\subsection*{广义柯西积分定理}

绕数的概念使我们得以扩展柯西积分定理与柯西积分公式的适用范围。我们先给出一个引理。

\qyxffntextdef{1}{这是富比尼定理的一个初等版本，将在第 III 卷中以完全一般的形式证明。另见习题 5.12。}
\begin{lemma}

设 $ f : U \to \mathbb{C} $ 全纯且 $ \Gamma \subset U $。

\begin{subitems}
\item

函数

\begin{equation*}
g : U\times U \to \mathbb{C} , \quad (z,w) \mapsto \begin{cases} (f(w) - f(z))/(w - z) , & z \ne w , \\ f'(z) , & z = w , \end{cases}
\end{equation*}

是连续的。

\item

映射

\[
h : U \to \mathbb{C} , \quad z \mapsto \int_{\Gamma}g(z,w)\,dw
\]

是解析的。

\proofhead
(i) 显然，$ g $ 在满足 $ z \ne w $ 的每个点 $ (z,w) $ 处连续。
设 $ z_{0} \in U $ 且 $ \varepsilon > 0 $。则存在 $ r > 0 $ 使得 $ \mathbb{D}(z_{0},r) \subset U $ 且对 $ \zeta \in \mathbb{D}(z_{0},r) $ 有 $ |f'(\zeta) - f'(z_{0})| < \varepsilon $。对 $ z,w \in \mathbb{D}(z_{0},r) $ 与 $ \gamma(t) := (1 - t)z + tw $（其中 $ t \in [0,1] $），由 $ \operatorname{im}(\gamma) \subset \mathbb{D}(z_{0},r) $ 与中值定理可得
\[
g(z,w) - g(z_{0},z_{0}) = \int_{0}^{1}\left[f'(\gamma(t)) - f'(z_{0})\right]dt .
\]

因此 $ |g(z,w) - g(z_{0},z_{0})| < \varepsilon $，这表明 $ g $ 在 $ (z_{0},z_{0}) $ 处连续。

\end{subitems}

(ii) 由 (i)，$ h $ 是良定义的。我们将借助莫雷拉定理（定理 5.24）验证 $ h $ 是解析的。设 $ \Delta $ 是 $ U $ 中的三角形道路。由习题 5.11 可得\qyxffnmark{1}

\begin{equation}
\int_{\Delta}\left(\int_{\Gamma}g(z,w)\,dw\right)dz = \int_{\Gamma}\left(\int_{\Delta}g(z,w)\,dz\right)dw .\tag{6.15}
\end{equation}

此外，引理 6.2 的证明表明，对每个 $ w \in U $，$ g(\cdot,w) $ 属于 $ C^{\omega}(U) $。因此由定理 5.25 可得，对 $ w \in \Gamma $ 有 $ \int_{\Delta}g(z,w)\,dz = 0 $，并由 (6.15) 得到 $ \int_{\Delta}h(z)\,dz = 0 $。因此 $ h $ 是解析的。\nobreak\hspace{0.35em}■

\end{lemma}

称 $ U $ 中的曲线 $ \Gamma $ \textbf{在 $ U $ 中同调于零}，如果它是闭的、分段连续可微的，并且对 $ a \in U^{c} $ 有 $ w(\Gamma,a) = 0 $。\footnote{当 $ U = \mathbb{C} $ 时，每条闭的分段 $ C^{1} $ 曲线都同调于零。}

\begin{theorem}

（柯西积分定理与柯西积分公式的同调版本）设 $ U $ 是 $ \mathbb{C} $ 中的开集且 $ f $ 在 $ U $ 中全纯。若曲线 $ \Gamma $ 在 $ U $ 中同调于零，则

\begin{equation}
\frac{1}{2\pi i}\int_{\Gamma}\frac{f(\zeta)}{\zeta - z}\,d\zeta = w(\Gamma,z)f(z) \quad \text{对 } z \in U\setminus\Gamma ,\tag{6.16}
\end{equation}

且

\begin{equation}
\int_{\Gamma}f(z)\,dz = 0 .\tag{6.17}
\end{equation}

\proofhead
我们沿用引理 6.19 的记号。

\begin{subitems}
\item

显然，(6.16) 等价于陈述

\begin{equation}
h(z) = 0 \quad \text{对 } z \in U\setminus\Gamma .\tag{6.18}
\end{equation}

为验证这一点，设 $ U_{0} := \{z \in \mathbb{C}\setminus\Gamma ; w(\Gamma,z) = 0\} $ 且

\[
h_{0}(z) := \frac{1}{2\pi i}\int_{\Gamma}\frac{f(\zeta)}{\zeta - z}\,d\zeta \quad \text{对 } z \in U_{0} .
\]

根据定理 6.14，且由于 $ w(\Gamma,\cdot) $ 连续，$ U_{0} $ 是开集。我们还有

\[
h_{0}(z) = \frac{1}{2\pi i}\int_{\Gamma}\frac{f(\zeta)}{\zeta - z}\,d\zeta - f(z)w(\Gamma,z) = h(z) \quad \text{对 } z \in U\cap U_{0} .
\]

由于 $ h_{0} $ 全纯，解析函数的唯一性定理（定理 V.3.13）表明，存在一个在 $ U\cup U_{0} $ 上全纯的函数 $ H $，使得 $ H \supset h_{0} $ 且 $ H \supset h $。由假设，对 $ a \in U^{c} $ 有 $ w(\Gamma,a) = 0 $。于是 $ U^{c} $ 属于 $ U_{0} $，从而得知 $ H $ 是整函数。

\item

设 $ R > 0 $ 且 $ \Gamma \subset R\mathbb{D} $。由于 $ R\mathbb{D}^{c} $ 位于 $ \mathbb{C}\setminus\Gamma $ 的无界连通分支中，我们有 $ R\mathbb{D}^{c} \subset U_{0} $。现设 $ \varepsilon > 0 $。令 $ M := \max_{\zeta\in\Gamma}|f(\zeta)| $ 且 $ R' := R + L(\Gamma)M/2\pi\varepsilon $。对 $ z \in R'\mathbb{D}^{c} $，我们便有

\[
|\zeta - z| \ge |z| - |\zeta| > L(\Gamma)M/2\pi\varepsilon \quad \text{对 } \zeta \in \Gamma ,
\]

并且我们得到

\begin{equation}
|h_{0}(z)| \le \frac{1}{2\pi}\int_{\Gamma}\left|\frac{f(\zeta)}{\zeta - z}\right|d\zeta < \varepsilon \quad \text{对 } z \in R'\mathbb{D}^{c} .\tag{6.19}
\end{equation}

由于 $ h_{0} \subset H $，且 $ H $ 作为整函数在有界集上有界，由此可得 $ H $ 在整个 $ \mathbb{C} $ 上有界。因此由刘维尔定理与 (6.19) 可得 $ H $ 处处为零。现在 $ h \subset H $ 蕴含 (6.18)。

\item

设 $ a \in \mathbb{C}\setminus\Gamma $ 且

\[
F : U \to \mathbb{C} , \quad z \mapsto (z - a)f(z) .
\]

由于 $ F $ 全纯且 $ F(a) = 0 $，(6.16) 表明

\[
\frac{1}{2\pi i}\int_{\Gamma}f(z)\,dz = \frac{1}{2\pi i}\int_{\Gamma}\frac{F(z)}{z - a}\,dz = w(\Gamma,a)F(a) = 0 .\quad\blacksquare
\]
\end{subitems}

\end{theorem}

\begin{remark}

\begin{items}
\item

若 $ U $ 单连通，则每条闭的分段 $ C^{1} $ 曲线都在 $ U $ 中同调于零。于是定理 6.20 是定理 5.5 与 5.9 的推广。

\proofhead
这由例 6.12(d) 推出。\nobreak\hspace{0.35em}■
\item

在定理 6.20 的假设下，我们有广义柯西导数公式

\[
w(\Gamma,z)f^{(k)}(z) = \frac{k!}{2\pi i}\int_{\Gamma}\frac{f(\zeta)}{(\zeta - z)^{k+1}}\,d\zeta \quad \text{对 } z \in U\setminus\Gamma , \quad k \in \mathbb{N} .
\]

\proofhead
这由定理 6.20 容易推出。\nobreak\hspace{0.35em}■
\end{items}

\end{remark}

\needspace{3\baselineskip}\subsection*{留数定理}

下一个定理把柯西积分定理的同调版本进一步推广到亚纯函数的情形。

\begin{theorem}

（留数定理）设 $ U $ 是 $ \mathbb{C} $ 中的开集且 $ f $ 在 $ U $ 中亚纯。又设 $ \Gamma $ 是 $ U\setminus P(f) $ 中在 $ U $ 中同调于零的曲线。则

\begin{equation}
\int_{\Gamma}f(z)\,dz = 2\pi i\sum_{p\in P(f)}\operatorname{Res}(f,p)w(\Gamma,p) ,\tag{6.20}
\end{equation}

其中和式中只有有限多项不等于零。

\proofhead
由推论 6.18 我们知道 $ A := \{a \in P(f) ; w(\Gamma,a) \ne 0\} $ 是有限集。于是 (6.20) 中的和式只有有限多项不等于零。

设 $ A = \{a_{0},\ldots,a_{m}\} $，且对 $ 0 \le j \le m $，$ f_{j} $ 是 $ f $ 在 $ a_{j} $ 处洛朗展开的主部。则 $ f_{j} $ 在 $ \mathbb{C}\setminus\{a_{j}\} $ 中全纯，并且 $ F := f - \sum_{j=0}^{m}f_{j} $ 在 $ a_{0},\ldots,a_{m} $ 处的奇点都是可去的（因为 $ F $ 在 $ a_{j} $ 处局部具有形式
\begin{equation*}
\begin{aligned}
F = g_{j} - \sum_{\substack{k=0 \\ k \ne j}}^{m}f_{j}
\end{aligned}
\end{equation*}

其中 $ g_{j} $ 是 $ f $ 在 $ a_{j} $ 处的正则部）。因此，由黎曼可去奇点定理，$ F $ 在

\[
U_{0} := U\setminus\left(P(f)\setminus A\right) = A\cup\left(U\setminus P(f)\right) .
\]

上有全纯延拓（仍记作 $ F $）。由于 $ \Gamma $ 位于 $ U\setminus P(f) $ 中且在 $ U $ 中同调于零，$ \Gamma $ 位于 $ U_{0} $ 中且在其中同调于零。因此由广义柯西积分定理可得 $ \int_{\Gamma}F\,dz = 0 $，这蕴含

\begin{equation}
\int_{\Gamma}f\,dz = \sum_{j=0}^{m}\int_{\Gamma}f_{j}\,dz .\tag{6.21}
\end{equation}

由于 $ a_{j} $ 是 $ f $ 的极点，存在 $ n_{j} \in \mathbb{N}^{\times} $ 与 $ c_{jk} \in \mathbb{C} $（对 $ 1 \le k \le n_{j} $ 与 $ 0 \le j \le m $）使得

\[
f_{j}(z) = \sum_{k=1}^{n_{j}}c_{jk}(z - a_{j})^{-k} \quad \text{对 } 0 \le j \le m .
\]

因此由注 6.21(b)（取 $ f = \mathbf{1} $）可得

\[
\int_{\Gamma}f_{j}\,dz = \sum_{k=1}^{n_{j}}c_{jk}\int_{\Gamma}\frac{dz}{(z - a_{j})^{k}} = 2\pi i\,c_{j1}w(\Gamma,a_{j}) .
\]

由于 $ c_{j1} = \operatorname{Res}(f,a_{j}) $，利用 $ A $ 的定义，结论由 (6.21) 推出。\nobreak\hspace{0.35em}■

\end{theorem}

\needspace{3\baselineskip}\subsection*{傅里叶积分}

留数定理有许多重要应用，我们现在通过计算若干反常积分来举例说明。我们集中于一个特别重要的类，即傅里叶积分。

设 $ f : \mathbb{R} \to \mathbb{C} $ 绝对可积。对 $ p \in \mathbb{R} $，有 $ |e^{-ipx}| = 1 $（对 $ x \in \mathbb{R} $），因此函数 $ x \mapsto e^{-ipx}f(x) $ 绝对可积。于是 $ f $ 在 $ p $ 处的\textbf{傅里叶积分}，由

\begin{equation}
\widehat{f}(p) := \int_{-\infty}^{\infty}e^{-ipx}f(x)\,dx \in \mathbb{C} ,\tag{6.22}
\end{equation}

给出，对每个 $ p \in \mathbb{R} $ 都有定义。通过 (6.22) 定义的函数 $ \widehat{f} : \mathbb{R} \to \mathbb{C} $ 称为 $ f $ 的\textbf{傅里叶变换}。

下一个定理与随后的例子表明，留数定理有助于计算许多傅里叶变换。

\begin{proposition}

设 $ \mathbb{C} $ 上的亚纯函数 $ f $ 具有下列性质

\begin{subitems}
\item

$ P(f) $ 有限，

\item

$ P(f)\cap\mathbb{R} = \varnothing $，

\item

$ \lim_{|z|\to\infty}zf(z) = 0 $。
则

\begin{equation*}
\widehat{f}(p) = \begin{cases} -2\pi i\sum_{\substack{z\in P(f) \\ \operatorname{Im} z < 0}}\operatorname{Res}(f e^{-ip\cdot},z) , & p \ge 0 , \\ 2\pi i\sum_{\substack{z\in P(f) \\ \operatorname{Im} z > 0}}\operatorname{Res}(f e^{-ip\cdot},z) , & p \le 0 . \end{cases}
\end{equation*}

\proofhead
设 $ p \le 0 $。由假设 (i)，存在 $ r > 0 $ 使得 $ P(f) \subset r\mathbb{D} $。我们选取 $ \Gamma $ 作为 $ V := r\mathbb{D}\cap\{z \in \mathbb{C} ; \operatorname{Im} z > 0\} $ 的正向边界。
\end{subitems}

\bookfigure[0.41]{figures/70c33af6755ed277.png}

于是我们有

\begin{equation*}
w(\Gamma,z) = \begin{cases} 1 , & z \in V , \\ 0 , & z \in \left(\overline{V}\right)^{c} , \end{cases}
\end{equation*}

（见习题 13）。因此由留数定理可得

\[
\int_{-r}^{r}f(x)e^{-ipx}\,dx + i\int_{0}^{\pi}f(re^{it})e^{-ipre^{it}}re^{it}\,dt = 2\pi i\sum_{z\in V}\operatorname{Res}(f e^{-ip\cdot},z) .
\]

由于 $ p \le 0 $，半圆上的积分可估计为

\[
\left|\int_{0}^{\pi}f(re^{it})e^{-ipre^{it}}re^{it}\,dt\right| \le \pi\max_{0\le t\le\pi}\left|rf(re^{it})e^{pr\sin t}\right| \le \pi\max_{|z|=r}|zf(z)| .
\]

假设 (iii) 现在推出 $ p \le 0 $ 情形的结论。

$ p \ge 0 $ 的情形类似，但积分路径在 $ \{z \in \mathbb{C} ; \operatorname{Im} z \le 0\} $ 中。\nobreak\hspace{0.35em}■

\end{proposition}

\begin{example}

\begin{items}
\item

对 $ f(x) := 1/(x^{2} + a^{2}) $（其中 $ a > 0 $），我们有

\[
\widehat{f}(p) = \pi e^{-|p|a}/a \quad \text{对 } p \in \mathbb{R} .
\]

\proofhead
函数 $ f(z) := 1/(z^{2} + a^{2}) $ 在 $ \mathbb{C} $ 中亚纯。进一步，$ f|\mathbb{R} $ 绝对可积（见例 VI.8.4(e)），且 $ \lim_{|z|\to\infty}zf(z) = 0 $。对 $ z \mapsto e^{-ipz}f(z) $ 在单极点 $ \pm ia $ 处的留数，我们在例 6.11(e) 中已求得
\[
\operatorname{Res}\left(e^{-ipz}f(z),\pm ia\right) = \pm e^{\pm pa}/2ia .
\]

因此结论由命题 6.23 推出。\nobreak\hspace{0.35em}■

\item

我们有

\[
\int_{-\infty}^{\infty}\frac{dx}{x^{4} + 1} = \frac{\pi}{\sqrt{2}} .
\]

\proofhead
我们考虑 $ \mathbb{C} $ 中的亚纯函数 $ f $，由 $ f(z) := 1/(z^{4} + 1) $ 定义。$ f $ 在 $ \{z \in \mathbb{C} ; \operatorname{Im} z > 0\} $ 中的极点是 $ z_{0} := (1 + i)/\sqrt{2} $ 与 $ z_{1} := iz_{0} = (-1 + i)/\sqrt{2} $。由例 6.11(f)，我们知道
\[
\operatorname{Res}(f e^{-ip\cdot},z_{0}) + \operatorname{Res}(f e^{-ip\cdot},z_{1}) = \frac{1}{4z_{0}^{3}}\left(e^{-ipz_{0}} + i\,e^{-ipz_{1}}\right) .
\]

现在由命题 6.23 可得

\[
\int_{-\infty}^{\infty}\frac{dx}{x^{4} + 1} = \widehat{f}(0) = 2\pi i\left(\frac{1}{4z_{0}^{3}}(1 + i)\right) = \frac{\pi i}{2}(-z_{0})(1 + i) = \frac{\pi}{\sqrt{2}} .\quad\blacksquare
\]

即使 $ 0 $ 是 $ x \mapsto \sin(x)/x $ 的可去奇点且该积分不绝对收敛（见习题 VI.8.1），我们仍可用留数定理计算收敛的反常积分 $ \int_{\mathbb{R}}\sin(x)\,dx/x $。

\end{items}

\end{example}

\begin{proposition}

对 $ p \in \mathbb{R} $，我们有

\begin{equation*}
\lim_{R\to\infty}\int_{-R}^{R}\frac{\sin x}{x}e^{-ipx}\,dx = \begin{cases} \pi , & |p| < 1 , \\ \pi/2 , & |p| = 1 , \\ 0 , & |p| > 1 . \end{cases}
\end{equation*}

证明

\begin{subitems}
\item

设 $ R > 1 $。我们对整函数

\[
z \mapsto \frac{\sin z}{z}e^{-ipz}
\]

沿道路 $ \gamma_{R} $ 积分，该道路从 $ -R $ 沿实轴到 $ -1 $，再沿单位圆的上半部分到 $ +1 $，最后沿实轴到 $ R $。柯西积分定理给出

\[
\int_{-R}^{R}\frac{\sin x}{x}e^{-ipx}\,dx = \int_{\gamma_{R}}\frac{\sin z}{z}e^{-ipz}\,dz .
\]

由于 $ \sin z = (e^{iz} - e^{-iz})/2i $，可得

\[
\int_{-R}^{R}\frac{\sin x}{x}e^{-ipx}\,dx = \frac{1}{2i}\int_{\gamma_{R}}\left(e^{-iz(p-1)} - e^{-iz(p+1)}\right)\frac{dz}{z} .
\]
\item

现在我们希望计算积分

\[
h_{R}(q) := \frac{1}{2\pi i}\int_{\gamma_{R}}\frac{e^{-izq}}{z}\,dz \quad \text{对 } q \in \mathbb{R} .
\]

于是考虑由闭路 $ \gamma_{R} + k_{R}^{\pm} $ 参数化的曲线 $ \Gamma_{R}^{\pm} $，其中

\[
k_{R}^{\pm} : [0,\pi] \to \mathbb{C} , \quad t \mapsto Re^{\pm it} .
\]
\end{subitems}

\bookfigure[0.37]{figures/c5909d5be7eab704.png}

则 $ w(\Gamma_{R}^{+},0) = 0 $ 且 $ w(\Gamma_{R}^{-},0) = -1 $。由于 $ \operatorname{Res}(e^{-izq}/z,0) = 1 $，由留数定理可得

\[
h_{R}(q) = -\frac{1}{2\pi i}\int_{k_{R}^{+}}\frac{e^{-izq}}{z}\,dz = -\frac{1}{2\pi}\int_{0}^{\pi}e^{-iqRe^{it}}\,dt
\]

且

\[
h_{R}(q) = -1 - \frac{1}{2\pi i}\int_{k_{R}^{-}}\frac{e^{-izq}}{z}\,dz = -1 + \frac{1}{2\pi}\int_{-\pi}^{0}e^{-iqRe^{it}}\,dt .
\]

\begin{subitems}
\item[（iii）]

接下来我们要证明

\begin{equation}
\int_{0}^{\pi}e^{-iqRe^{it}}\,dt \to 0 \quad (R \to \infty) \quad \text{对 } q < 0 .\tag{6.23}
\end{equation}

设 $ \varepsilon \in (0,\pi/2) $。由于对 $ t \in [0,\pi] $ 有 $ q\sin t \le 0 $，可得

\[
\left|e^{-iqRe^{it}}\right| = e^{qR\sin t} \le 1 \quad \text{对 } q < 0 , \quad R > 1 , \quad t \in [0,\pi] ,
\]

并且我们得到

\begin{equation}
\left|\left(\int_{0}^{\varepsilon} + \int_{\pi - \varepsilon}^{\pi}\right)e^{-iqRe^{it}}\,dt\right| \le 2\varepsilon \quad \text{对 } q < 0 \text{ 且 } R > 1 .\tag{6.24}
\end{equation}

由于 $ q\sin\varepsilon < 0 $，存在 $ R_{0} > 1 $ 使得

\[
e^{qR\sin t} \le e^{qR\sin\varepsilon} \le \varepsilon \quad \text{对 } R \ge R_{0} \text{ 且 } \varepsilon \le t \le \pi - \varepsilon .
\]

因此

\[
\left|\int_{\varepsilon}^{\pi - \varepsilon}e^{-iqRe^{it}}\,dt\right| \le \varepsilon\pi \quad \text{对 } R \ge R_{0} ,
\]

它与 (6.24) 一起证明了 (6.23)。类似地，对 $ q > 0 $ 可验证

\[
\int_{-\pi}^{0}e^{-iqRe^{it}}\,dt \to 0 \quad (R \to \infty) .
\]
\item[（iv）]

容易验证，对 $ R > 1 $ 有 $ h_{R}(0) = -1/2 $。因此由 (ii) 与 (iii)，我们有

\begin{equation*}
\lim_{R\to\infty}h_{R}(q) = \begin{cases} 0 , & q < 0 , \\ -1/2 , & q = 0 , \\ -1 , & q > 0 . \end{cases}
\end{equation*}

由 (i)，我们知道

\[
\lim_{R\to\infty}\int_{-R}^{R}\frac{\sin x}{x}e^{-ipx}\,dx = \pi\lim_{R\to\infty}\left(h_{R}(p - 1) - h_{R}(p + 1)\right) ,
\]

结论由此推出。\nobreak\hspace{0.35em}■

\end{subitems}

\end{proposition}

\begin{corollary}
\[
\int_{-\infty}^{\infty}\frac{\sin x}{x}\,dx = \pi .
\]

\proofhead
反常积分 $ \int_{-\infty}^{\infty}(\sin x/x)\,dx $ 的收敛性由习题 VI.4.11(ii) 与 VI.8.1(ii) 推出。其值由命题 6.25 得到。\nobreak\hspace{0.35em}■
\end{corollary}

\begin{remark}

\begin{items}
\item

由命题 6.25 \textbf{不能}推出反常积分

\[
\int_{-\infty}^{\infty}\frac{\sin x}{x}e^{ipx}\,dx
\]

对 $ p \in \mathbb{R}^{\times} $ 收敛。要使其成立，我们必须确认 $ \lim_{R\to\infty}\int_{0}^{R}\cdots\,dx $ 与 $ \lim_{R\to\infty}\int_{-R}^{0}\cdots\,dx $ 各自独立地存在。当极限

\[
\mathrm{PV}\int_{-\infty}^{\infty}f := \lim_{R\to\infty}\int_{-R}^{R}f(x)\,dx
\]

对一个（连续）函数 $ f : \mathbb{R} \to \mathbb{C} $ 存在时，称它为 $ \int_{-\infty}^{\infty}f $ 的\textbf{柯西主值}\footnote{见习题 VI.8.9。}。若 $ f $ 反常可积，则柯西主值存在且与 $ \int_{-\infty}^{\infty}f $ 一致。然而，该陈述的逆命题不成立，如例 $ f(t) := t $（对 $ t \in \mathbb{R} $）所示。

\item

由于函数 $ g : \mathbb{R} \to \mathbb{R} $，$ x \mapsto \sin(x)/x $ 不绝对可积，我们不能给 $ g $ 指派一个傅里叶变换 $ \widehat{g} $。取而代之，可以（而且必须！）使用傅里叶变换的一个更一般的定义——来自分布理论\footnote{关于分布理论的入门及其许多重要应用，见例如 [Sch65]。}的框架——正如我们在这里所做的。在一般理论中，$ g $ 的傅里叶变换有定义，且 $ \widehat{g} $ 就是由命题 6.25 中给出的柯西主值所定义的同一个分段常值函数。

\item

\textbf{傅里叶变换} $ f \mapsto \widehat{f} $ 在数学与物理的许多领域都非常重要。然而，对这一映射的更精确研究有待于勒贝格积分理论。因此，我们将在第 III 卷中回到傅里叶积分。\nobreak\hspace{0.35em}■

\end{items}

关于柯西积分定理与留数定理的大量进一步应用，以及全纯函数与亚纯函数理论的更深入处理，我们请你参考习题与复分析的文献（例如 [Car66]、[Con78]、[FB95]、[Rem92]）。

\end{remark}

\exerciseshead

\begin{exercises}
\item

设 $ f $ 在 $ U $ 中全纯，且 $ z_{0} \in U $ 满足 $ f'(z_{0}) \ne 0 $。此外，设 $ g : \mathbb{C} \to \mathbb{C} $ 亚纯，且在 $ w_{0} := f(z_{0}) $ 处有单极点。证明 $ g \circ f $ 在 $ z_{0} $ 处有单极点，且 $ \operatorname{Res}(g \circ f, z_{0}) = \operatorname{Res}(g, w_{0})/f'(z_{0}) $。

\item

设函数 $ f $ 在 $ \mathbb{C} $ 中亚纯，且在 $ \Omega(r_{0}, r_{1}) $（$ r_{0} > 0 $）中具有洛朗展开 $ \sum_{n=-\infty}^{\infty}c_{n}z^{n} $，对 $ n \in \mathbb{N} $ 满足 $ c_{-n} \ne 0 $。证明或反驳 $ f $ 有本性奇点。（提示：考虑 $ z \mapsto 1/(z - 1) $ 在 $ \Omega(1, 2) $ 中的情形。）

\item

设 $ a, b \in \mathbb{C} $ 满足 $ 0 < \left|a\right| < \left|b\right| $，且

\[
f : \mathbb{C}\setminus\{a,b\} \to \mathbb{C} , \quad z \mapsto \frac{a - b}{(z - a)(z - b)} .
\]

\begin{items}
\item

确定 $ f $ 在 $ \Omega(0,\left|a\right|) $、$ \Omega(\left|a\right|,\left|b\right|) $ 与 $ \Omega(\left|b\right|,\infty) $ 中关于 $ 0 $ 的洛朗展开。

\item

$ f $ 在 $ a $ 处的洛朗展开是什么？（提示：几何级数。）

\end{items}

\item

计算

\[
\int_{\partial\mathbb{D}}\frac{4\,dz}{1 + 4z^{2}} \quad \text{与} \quad \int_{\partial\mathbb{D}}\frac{e^{z}\sin z}{(1 - e^{z})^{2}}\,dz .
\]
\item

全纯函数 $ f : U\setminus\{z_{0}\} \to \mathbb{C} $ 在 $ z_{0} \in U $ 处有孤立奇点。证明下列陈述等价：

\begin{subitems}
\item

$ z_{0} $ 是 $ n $ 阶极点。

\item

$ g : U\setminus\{z_{0}\} \to \mathbb{C} $，$ z \mapsto (z - z_{0})^{n}f(z) $ 在 $ z_{0} $ 处有可去奇点，且 $ g(z_{0}) \ne 0 $。

\item

存在 $ \varepsilon, M_{1}, M_{2} > 0 $ 使得 $ \mathbb{D}(z_{0},\varepsilon) \subset U $，且

\[
M_{1}\left|z - z_{0}\right|^{-n} \le \left|f(z)\right| \le M_{2}\left|z - z_{0}\right|^{-n} \quad \text{对 } z \in \mathbb{D}^{\bullet}(z_{0},\varepsilon) .
\]
\end{subitems}

\item

证明：$ z_{0} \in U $ 是在 $ U\setminus\{z_{0}\} $ 中全纯的函数 $ f $ 的极点，当且仅当 $ z \to z_{0} $ 时 $ \left|f(z)\right| \to \infty $。

\item

设 $ z_{0} $ 是 $ f $ 的孤立奇点，$ f $ 在 $ U\setminus\{z_{0}\} $ 中全纯。证明下列陈述等价：

\begin{subitems}
\item

$ z_{0} $ 是本性奇点。

\item

对每个 $ w_{0} \in \mathbb{C} $，存在 $ U\setminus\{z_{0}\} $ 中的序列 $ (z_{n}) $ 使得 $ \lim z_{n} = z_{0} $ 且 $ \lim f(z_{n}) = w_{0} $。
（提示：“(i)$\Rightarrow$(ii)” 若该陈述为假，则存在 $ w_{0} \in \mathbb{C} $ 与 $ r, s > 0 $ 使得 $ f(\mathbb{D}(z_{0},r)) \cap \mathbb{D}(w_{0},s) = \varnothing $。于是 $ g : \mathbb{D}^{\bullet}(z_{0},r) \to \mathbb{C} $，$ z \mapsto 1/(f(z) - w_{0}) $ 全纯且有界。借助定理 6.6 与习题 6，排除 $ g(z_{0}) = 0 $ 与 $ g(z_{0}) \ne 0 $ 这两种情形。）

\end{subitems}

\item

确定

\[
z \mapsto e^{1/(z - 1)}/(e^{z} - 1) , \quad z \mapsto (z + 1)\sin(1/(z - 1)) .
\]

的奇点。它们是极点还是本性奇点？

\item

验证 $ i $ 是 $ z \mapsto \sin(\pi/(z^{2} + 1)) $ 的本性奇点。（提示：习题 7。）

\item

设 $ U $ 单连通，$ a \in U $ 是 $ f $ 的孤立奇点，$ f $ 是 $ U\setminus\{a\} $ 中的全纯函数。证明 $ f $ 在 $ U\setminus\{a\} $ 中有原函数当且仅当 $ \operatorname{Res}(f,a) = 0 $。

\item

证明注 6.21(a)。

\item

对 $ p \in (0, 1) $，证明

\[
PV\int_{-\infty}^{\infty}\frac{e^{rx}}{1 + e^{x}}\,dx = \frac{\pi}{\sin(p\pi)} .
\]

（提示：沿以 $ \pm R $ 与 $ \pm R + 2\pi i $ 为顶点的矩形路径逆时针对 $ z \mapsto e^{pz}(1 + e^{z})^{-1} $ 积分。）

\item

设 $ \Gamma_{j} $ 是分段 $ C^{1} $ 曲线，由 $ \gamma_{j} \in C(I,\mathbb{C}) $ 参数化，$ j = 0, 1 $。又设 $ z \in \mathbb{C} $ 满足对 $ t \in I $ 有 $ \left|\gamma_{0}(t) - \gamma_{1}(t)\right| < \left|\gamma_{0}(t) - z\right| $。证明 $ w(\Gamma_{0},z) = w(\Gamma_{1},z) $。（提示：对 $ \gamma := (\gamma_{1} - z)/(\gamma_{0} - z) $，证明

\[
w([\gamma],0) = w(\Gamma_{1},z) - w(\Gamma_{0},z)
\]

且对 $ t \in I $ 有 $ \left|1 - \gamma(t)\right| < 1 $。）

\item

令 $ N(f) := \{\, z \in U\setminus P(f) \ ;\ f(z) = 0 \,\} $ 定义亚纯函数 $ f $ 的\textbf{零点集}。证明下列结论：

\begin{subitems}
\item

若 $ f \ne 0 $，则 $ N(f) $ 是 $ U $ 的离散子集。

\item

函数

\[
1/f : U\setminus N(f) \to \mathbb{C} , \quad z \mapsto 1/f(z)
\]

在 $ U $ 中亚纯，且 $ P(1/f) = N(f) $、$ N(1/f) = P(f) $。
（提示：(i) 解析函数的唯一性定理。(ii) 习题 6。）

\end{subitems}

\item

证明 $ U $ 中亚纯函数的全体关于逐点加法与乘法构成一个域。

\item

证明

\[
\int_{-\infty}^{\infty}\frac{x^{2}}{1 + x^{4}}\,dx = \frac{\pi}{\sqrt{2}} .
\]

（提示：命题 6.23。）

\item

验证

\[
\int_{0}^{\infty}\frac{\cos x}{x^{2} + a^{2}}\,dx = \frac{\pi e^{-a}}{2a} \quad \text{对 } a > 0 .
\]

（提示：例 6.24(a)。）

\item

设 $ f $ 在 $ U $ 中亚纯，$ f \ne 0 $，且 $ g := f'/f $。证明

\begin{subitems}
\item

$ g $ 在 $ U $ 中亚纯且只有单极点；

\item

若 $ z_{0} $ 是 $ f $ 的 $ m $ 阶零点，则 $ \operatorname{Res}(g,z_{0}) = m $；

\item

若 $ z_{0} $ 是 $ f $ 的 $ m $ 阶极点，则 $ \operatorname{Res}(g,z_{0}) = -m $。

\end{subitems}

\item

设 $ f $ 在 $ U $ 中亚纯，$ P $ 是 $ U\setminus\left(N(f) \cup P(f)\right) $ 中在 $ U $ 内同调于零的曲线。对 $ z \in N(f) \cup P(f) $，用 $ \nu(z) $ 表示 $ z $ 的重数。则

\[
\frac{1}{2\pi i}\int_{\Gamma}\frac{f'(z)}{f(z)}\,dz = \sum_{n\in N(f)}w(\Gamma,n)\nu(n) - \sum_{p\in P(f)}w(\Gamma,p)\nu(p) .
\]
\item

对 $ 0 < a < b < 1 $，计算

\[
\frac{1}{2\pi i}\int_{\partial\mathbb{D}}\frac{2z - (a + b)}{z^{2} - (a + b)z + ab}\,dz .
\]
\item

确定 $ \int_{-\infty}^{\infty}dx/(x^{4} + 4x^{2} + 3) $。

\end{exercises}
